\documentclass[mathematics,review,submit,moreauthors]{Definitions/mdpi} 

% MDPI internal commands - do not modify
\firstpage{1} 
\makeatletter 
\setcounter{page}{\@firstpage}
\makeatother
\pubvolume{1}
\issuenum{1}
\articlenumber{0}
\pubyear{2026}
\copyrightyear{2026}
%\externaleditor{Firstname Lastname} % More than 1 editor, please add `` and '' before the last editor name
\datereceived{ } 
\daterevised{ } % Comment out if no revised date
\dateaccepted{ } 
\datepublished{ } 
%\datecorrected{} % For corrected papers: "Corrected: XXX" date in the original paper.
%\dateretracted{} % For retracted papers: "Retracted: XXX" date in the original paper.
%\doinum{} % Used for some special journals, like molbank
%\pdfoutput=1 % Uncommented for upload to arXiv.org
%\CorrStatement{yes}  % For updates
%\longauthorlist{yes} % For many authors that exceed the left citation part
%\IsAssociation{yes} % For association journals

%=================================================================
% Add packages and commands here. The following packages are loaded in our class file: fontenc, inputenc, calc, indentfirst, fancyhdr, graphicx, epstopdf, lastpage, ifthen, float, amsmath, amssymb, lineno, setspace, enumitem, mathpazo, booktabs, titlesec, etoolbox, tabto, xcolor, colortbl, soul, multirow, microtype, tikz, totcount, changepage, attrib, upgreek, array, tabularx, pbox, ragged2e, tocloft, marginnote, marginfix, enotez, amsthm, natbib, hyperref, cleveref, scrextend, url, geometry, newfloat, caption, draftwatermark, seqsplit
% Display paragraph titles in bold, with body text on the same line (overrides the non-bold block default of mdpi.cls)
\titleformat{\paragraph}[runin]{\fontsize{10}{10}\selectfont\bfseries}{}{0pt}{}[.]
\titlespacing{\paragraph}{0pt}{12pt}{1em}
\usepackage{rotating} % landscape (sideways) tables
\usepackage{listings} % verbatim prompts (appendix), with line breaking
\usepackage{wrapfig} % half-width inserts with text wrapped around (fig:one-rotation)
\definecolor{promptbg}{HTML}{FEF8EC}    % same orange tint as the prompt bubbles of fig:selection-flow
\definecolor{promptframe}{HTML}{E69F00}
\lstdefinestyle{prompt}{basicstyle=\ttfamily\footnotesize, columns=fullflexible, keepspaces=true,
  escapeinside={(*}{*)}, % used to draw tab characters (TSV column separators) as arrows
  breaklines=true, breakatwhitespace=true, frame=single, rulecolor=\color{promptframe},
  backgroundcolor=\color{promptbg}, framesep=4pt, xleftmargin=4pt, xrightmargin=4pt}
% The MDPI text block is offset to the right by \extralength (the margin column),
% so sideways floats, which are centred on \textwidth, inherit that offset.
% Shifting the reference point left by \extralength/2 re-centres them on the paper.
\setlength{\rotFPtop}{0pt plus 1fil}
\addtolength{\rotFPtop}{-\extralength}

% -------------------------------------------------------------------------
\newcommand{\pnote}[1]{\textcolor{red}{PL: #1}}
\newcommand{\myline}{
  \vspace{5mm}
  \hrule
  \vspace{5mm}
  }

%=================================================================
% Please use the following mathematics environments: Theorem, Lemma, Corollary, Proposition, Characterization, Property, Problem, Example, ExamplesandDefinitions, Hypothesis, Remark, Definition, Notation, Assumption
%% For proofs, please use the proof environment (the amsthm package is loaded by the MDPI class).

%=================================================================
\Title{Symmetries and Neural Networks}
\newcommand{\orcidauthorA}{0009-0004-2050-916X}
\newcommand{\orcidauthorB}{0000-0002-0610-9168}
\newcommand{\orcidauthorC}{0000-0003-1847-5031}
\Author{Pierrick Leroy $^{*}$\orcidA{}, Giovanni Petri\orcidC{} and Francesco Vaccarino\orcidB{}}
% MDPI internal command: Authors, for metadata in PDF
\AuthorNames{Pierrick Leroy, Giovanni Petri and Francesco Vaccarino}
\address[1]{Department of Mathematical Sciences, Politecnico di Torino, Turin, Italy}
\corres{Correspondence: pierrick.leroy@polito.it}

\abstract{Deep learning is a field of symmetries: imposed by design in architectures, augmentations or training frameworks, emerging from the parametrization, unknown and to be discovered, or holding only partially or approximately.
Starting from the concept of symmetry, we systematically review more than 300 works from the literature from a mathematical perspective.
We propose a faceted taxonomy whose six axes (exactness, purpose, object, a priori knowledge, state, explicitness) map, locate and contrast occurrences of symmetry across subfields.
We pay particular attention to which object carries the symmetry, from the learning problem to the training framework.
We end with underexplored areas.}


\keyword{deep learning, symmetry, neural network, invariance, equivariance, group theory}

\setlength{\headheight}{20pt}
\begin{document}

\section{Introduction}

% Symmetries
In plane geometry, familiar symmetries are reflections, which mirror the plane across a line, rotations, which spin it around a point, and translations, which shift it without fixed points.
Symmetry can also be abstract: the symmetries of a set are its bijections onto itself, the permutations.

Group theory is the language of symmetry.
It is rooted in geometry but also, and perhaps first, in the theory of algebraic equations and number theory~\citep{wussing1984}.
Ideas underlying the notion of group go back to Euler and Lagrange in the $18$th century.
In his work on solid geometry, Legendre took the conceptual step of viewing a symmetry as a transformation between objects regarded as equivalent~\citep{legendre1794,hon2008fststmoarsc}.
Galois coined the term ``group'' around $1830$.
Felix Klein's Erlangen program ($1872$) made groups an organizing principle: a geometry is understood through the symmetries that preserve its structures.
Arthur Cayley ($1854$), then Heinrich Weber ($1882$), proposed axiomatic definitions, but the universally accepted modern one is due to Walther von Dyck, Klein's assistant.
The group axioms capture the transformations that leave an object unchanged: composing two of them gives another (closure), doing nothing is one (identity), each can be undone since nothing is lost (inverse), and associativity holds automatically.
Beyond mathematics, symmetry is an important concept in physics, chemistry, biology, and even architecture and art~\citep{weyl1952,makse2025}.

% Deep learning
Deep learning is now the dominant approach in vision, language and the sciences.
A trained model is the output of a pipeline (data and augmentations, architecture, loss, optimizer, training protocol) whose components are often adopted because they work: progress has been largely empirical.
The resulting models and their training dynamics are hard to interpret, and the subfields that study them (NLP, vision, graph learning, self-supervised learning, learning theory, scientific machine learning) have developed separate vocabularies and practices.

% Symmetries and deep learning
Symmetry recurs at each stage of this pipeline.
Some symmetries are imposed on purpose (an equivariant layer, a contrastive loss, an augmentation policy).
Others appear without being asked for, such as parameter-space symmetries (hidden-neuron permutations, rescaling with batch normalization).
Some are unknown in advance and must be discovered, and some hold only approximately or on part of the domain.
Working explicitly with a symmetry can turn a heuristic into a precise statement: invariance, $f(g\cdot x)=f(x)$, says what plays no role, and equivariance, $f(g\cdot x)=g\cdot f(x)$, says how it plays a role.
The notion can play opposite roles: an inductive prior on the data, a confounder on the parameters.
Existing surveys each cover one of these objects or one application domain (\Cref{sec:related-work}).

% Objective
In contrast, we start from the concept of symmetry and go broad: with a mathematical perspective, we list where, why and how it occurs across deep learning, including when it is unknown or incomplete.
Geometric deep learning~\citep{bronstein2021} follows Klein's Erlangen program to classify architectures by the symmetries of their domains.
We extend this perspective from architectures to every object that can carry a symmetry (problem, model, optimization, parameter, data point, training framework, see \Cref{fig:sym-loci}) and classify the occurrences themselves, to tell apart the many faces of symmetry.
As Alan Perlis put it: ``Symmetry is a complexity-reducing concept […]; seek it everywhere.''~\citep{perlis1982}
Our contributions are threefold.
(i)~A faceted taxonomy (\Cref{tab:taxonomy}) whose six axes locate, link and contrast occurrences of symmetry across subfields.
(ii)~An inventory of $322$ works organized along this taxonomy, underpinned by a hybrid systematic review and the coding of the corpus (\Cref{sec:materials-methods}).
(iii)~The identification of underexplored areas and possible gaps (\Cref{sec:discussion}).
We hope that this single perspective will foster the transfer of knowledge and methods between subfields.
The survey can be of interest to researchers in a covered field seeking a bird's-eye view, mathematicians seeking areas where symmetry applies, and deep learning practitioners who want to make the notion of invariance precise.

% Narrative
After the methodology (\Cref{sec:materials-methods}) and a background that the accustomed reader can skip (\Cref{sec:background}), the survey follows a conceptual progression: symmetry imposed by design (\Cref{sec:design-principle}), symmetry as an artifact of the parametrization (\Cref{sec:parameter-space-symmetries}), unknown symmetries to be discovered (\Cref{sec:sym-disco}), and incomplete symmetries (\Cref{sec:incomplete-symmetries}).
This order is one choice among several: equivariant optimizers, for instance, are treated with parameter-space symmetries (\Cref{sec:bug-or-feature}) but also belong to symmetry by design.

\subsection{Related work}\label{sec:related-work}
Previous works have studied the literature on symmetries in deep learning.
\citet{ilina2022sa} surveys symmetric architectures such as Siamese networks.
\citet{zhao2025} focuses on parameter-space symmetries.
\citet{sipiran202523scgpis} reviews symmetry detection methods for 3D shapes.
\citet{bronstein2021} is the reference textbook on geometric deep learning, and \citet{gerken2021,carlsson2025} study its mathematical formalism.
More specialized works focus on equivariant convolutions \citep{basheer2025ijdsa,habte2023}, 3D rotations \citep{fei2024air} and $SE(3)$-equivariant robot learning \citep{seo2025a}.
\citet{rath2024air} reviews methods that integrate geometric prior knowledge into neural networks so that their representations are invariant or equivariant.
\citet{shorten2019jbd,mumuni2022a} survey data augmentation.
Each of these works focuses on a few objects carrying the symmetry or on one application domain.

\section{Materials and Methods}\label{sec:materials-methods}

\subsection{Methodology}
The corpus was collected over a two-year period, starting from an initial collection that was progressively expanded and completed in the final phase.
The selection process is described in \Cref{fig:selection-flow}.
A reference is discarded if it does not satisfy at least one of the following criteria: (A) it is published in a peer-reviewed venue, (B) it is highly cited (more than 50 citations on Google Scholar), (C) it is recent (made available less than two years ago).
Criterion (C) exists because the field is fast-moving and a citation/peer-review threshold alone would systematically remove emerging work.
Following this protocol, $408$ works were identified, out of which $322$ are cited in the current version.
\Cref{fig:corpus-stats} summarizes the corpus: it reports the distribution over publication years and the pace at which the corpus was assembled.

\begin{figure}[ht]
\begin{adjustwidth}{-\extralength}{0cm}
    \centering
    \includegraphics[width=\linewidth]{images/selection-flow.pdf}
\end{adjustwidth}
    \caption{Selection of the corpus and construction of its knowledge base.
    Works are identified by monitoring (daily research, references of other works, recommendations, Scholar Inbox), keyword search (Google Scholar, Scopus) and dedicated tools (Research Rabbit, Google Scholar Labs, Connected Papers, LLMs: ChatGPT, Gemini, Claude, Elicit), then screened on criteria (A)--(C).
    Each retained work is converted to Markdown, then summarized and indexed by an LLM (abridged prompts, full versions in \Cref{app:prompts}).
    The summaries support the coding of the corpus along the axes of \Cref{tab:taxonomy} by the typed language model $\texttt{jev-1.13.0}$.}
    \label{fig:selection-flow}
\end{figure}

\begin{figure}[ht]
    \centering
    \includegraphics[width=\textwidth]{images/zotero_stats.pdf}
    \caption{Composition and construction of the corpus. \textbf{Left.} Number of works per publication year.
    The first bar aggregates all works published between $1985$ and $2010$ (no collected work is in the $2011$--$2013$ period). 
    \textbf{Right.} Cumulative number of works as a function of the date of addition.
    (1)~The project was initialized in July $2024$. (2)~While mapping the field, works were added at a steady rate of about $5$ per month. (3)~From May $2026$, filling gaps and writing the first versions raised the rate to about $30$ per month.}
    \label{fig:corpus-stats}
\end{figure}

\subsection{Taxonomy}
We propose a faceted taxonomy in \Cref{tab:taxonomy} to categorize occurrences of symmetries in works.
\Cref{fig:sym-loci} further details the \emph{Object} axis of the taxonomy.
This taxonomy was used to code the corpus with the typed-output model $\texttt{jev-1.13.0}$, which provides a $6$-dimensional classification for each paper, or a $17$-dimensional vector of coordinates in $[0,1]$ when the output probabilities associated with each choice are recorded.
The coded corpus was used as an organizing tool to guide the framing of each part, to develop the taxonomy, to position the different parts, and to lay out and contrast their distinctive features.
\Cref{app:automated-taxonomy} gives additional details on the coding of the corpus, our method to validate both the coding and the taxonomy, and the way we leverage it to analyze the corpus.
We chose a typed-output model over a standard LLM because typed output matches the requirement of the task (a discrete set of categorical answers) while being significantly faster and more cost-effective, which allowed us to proceed by iterations.

At a high level, the taxonomy maps onto the sections as follows.
\Cref{sec:design-principle} \emph{prescriptively} builds in \emph{known} symmetries (\Cref{tab:taxonomy}).
In \Cref{fig:sym-loci}, equivariant architectures follow the arrow $f^*\dashrightarrow\mathcal{H}$, data augmentation $f^*\dashrightarrow f_\theta$, and framework symmetries fill the bottom row.
\Cref{sec:parameter-space-symmetries} covers the boxes where $G$ acts on $\Theta$, not on $\mathcal{X},\mathcal{Y}$.
\Cref{sec:sym-disco} discovers \emph{unknown} symmetries at the problem ($f^*$, $p$) or parameter ($f_\theta$) locus.
\Cref{sec:incomplete-symmetries} follows the exactness axis away from \emph{exact}: broken, \emph{extent-restricted} and \emph{approximate} symmetries.

\begin{sidewaystable}[p]
    \centering
    \footnotesize
    \renewcommand{\arraystretch}{1.15}
    \begin{tabularx}{\textheight}{|>{\raggedright\arraybackslash}p{1.5cm}|>{\raggedright\arraybackslash}p{3.1cm}|>{\raggedright\arraybackslash}p{2cm}|>{\RaggedRight\arraybackslash}X|>{\RaggedRight\arraybackslash}X|}
        \hline
        \textbf{Axis} & \textbf{Defining question} & \textbf{Values} & \textbf{Definition} & \textbf{Examples} \\
        \hline
        \multirow{3}{=}{Exactness} & \multirow{3}{=}{Does the symmetry hold fully, on a restricted extent, or up to a tolerance?} & Exact & The symmetry holds exactly for all $g\in G$, with zero error. & Geometric Deep Learning: design an equivariant layer. Data augmentations such as horizontal flip. Global parameter-space symmetries (permutation, ReLU rescaling). Conservation laws under gradient flow. \\
        \cline{3-5}
        & & Extent restricted & The symmetry holds only for a restricted set of transformations (a subset $S\subsetneq G$, i.e.\ a subgroup or a non-closed subset, or a weaker structure such as a semigroup or a groupoid), or only on part of the domain (local, instance-level). & Symmetry-breaking due to regularization. Data augmentations such as cropping and blurring (non-invertible) or rotation in MNIST ($6\leftrightarrow 9$). Instance-level parameter-space symmetries.\\
        \cline{3-5}
        & & Approximate & The symmetry holds for all $g\in G$ only up to a tolerance $\epsilon>0$. & Soft equivariance penalty. Conservation laws in discrete time. Quantification of the strength of a symmetry.\\
        \hline
        \multirow{2}{=}{Purpose} & \multirow{2}{=}{How does the work interact with the symmetry?} & Descriptive & The work seeks to understand the effect, validity, existence or consequences of a symmetry. & Understand the consequences of a parameter-space symmetry on the learning dynamics. Symmetry discovery in data or models.\\
        \cline{3-5}
        & & Prescriptive & The work proposes a method or guideline that enforces, exploits or removes a symmetry. & Design an equivariant layer. Canonicalization (e.g.\ extract an orbit representative). Align two trained networks before merging. \\
        \hline
        \multirow{6}{=}{Object} & \multirow{6}{=}{Which object carries the symmetry (its locus)? (more details in \Cref{fig:sym-loci})} & Problem & The symmetry is a property of the task, independent of any model: ground truth $f^*$ or data distribution $p$. & Rotation-invariant molecular properties. Symmetry discovery in data. \\
        \cline{3-5}
        & & Model & The symmetry is a property of the architecture, valid for every parameter $\theta$: hypothesis class $\mathcal{H}$ or implementation map $\varphi$. & Design an equivariant layer. Neuron permutation and rescaling symmetries. \\
        \cline{3-5}
        & & Optimization & The symmetry is a property of the learning rule: loss $\mathcal{L}$ on parameter space or optimizer $U$. & Conservation laws under gradient flow. Design an equivariant optimizer. \\
        \cline{3-5}
        & & Parameter & The symmetry is a property of a fixed parameter $\theta$ (instantiation): the actual implemented function $f_\theta$ or symmetric state $g\cdot\theta=\theta$. & Measure the equivariance error of a trained network. Track symmetric states along training. \\
        \cline{3-5}
        & & Data & The symmetry is a property of an individual data point, $g\cdot x = x$. & Curie principle: symmetric inputs forcing symmetric outputs of an equivariant network. \\
        \cline{3-5}
        & & Framework & The symmetry is a property of the training pipeline as a whole and can involve several forward calls (criterion, data ordering, joint embeddings). & Joint-embedding self-supervised learning. Permutation-invariant criterion (Hungarian loss). Permuting the batch ordering. \\
        \hline
        \multirow{2}{=}{A priori knowledge} & \multirow{2}{=}{Is the symmetry assumed to be known?} & Known & The work starts already knowing the symmetry. & Derive the consequences of a parameter-space symmetry. Design an equivariant optimizer.\\
        \cline{3-5}
        & & Unknown & The symmetry or its validity is not known, or not fully. & Symmetry discovery. Use RL to learn an automated data augmentation policy. Measure an equivariance error. Characterize the fiber of the implementation map. \\
        \hline
        \multirow{2}{=}{State} & \multirow{2}{=}{Does the symmetry interact with a time-indexed process?} & Dynamic & The symmetry is studied along a time-indexed process: its validity or strength varies along it, or it constrains the trajectory. & Conservation laws. Automated data augmentation. Tracking the evolution of symmetric states along training.\\
        \cline{3-5}
        & & Static & The symmetry holds in the same way at all times, or no process is involved. & Standard data augmentation. Quotienting a symmetry in parameter space. Design an equivariant layer.\\
        \hline
        \multirow{2}{=}{Explicitness} & \multirow{2}{=}{Does the work use a standard symmetry-language to formalize its thesis?} & Implicit & There is a symmetry but it is not abstracted, not presented as such. & Intuitive data augmentation. Empirical-only work. \citep{kobayashi2024anips}.\\
        \cline{3-5}
        & & Explicit & The symmetry is abstracted, named and formalized as such. & Works stating a group action, an invariance or equivariance property (most works in this survey). \\
        \hline
    \end{tabularx}
    \small
    \caption{Axes proposed to categorize occurrences of symmetries in a deep learning context. Each axis answers a defining question and takes one of a small set of values.}
    \label{tab:taxonomy}
\end{sidewaystable}

\begin{figure}[!tbp]
    \centering
    \includegraphics[width=\textwidth]{images/figure_sym-loci.pdf}
    \caption{Loci of symmetry, a refinement of the \emph{Object} axis of the taxonomy (\Cref{tab:taxonomy}).
    Each box is an object on which a symmetry can act, with its defining relation.
    Solid arrows indicate necessity, dashed arrows another type of interaction, e.g.\ a design choice.
    The figure is not exhaustive.
    Architectural symmetries (\Cref{sec:architectural-sym}) are not shown, and the restriction in extent is not specific to $\varphi$: it can also be applied, for example, to the ground truth function, as discussed in \Cref{sec:domain-restriction}.
    }
    \label{fig:sym-loci}
\end{figure}

\paragraph{AI statement}
AI tools were used at multiple stages, but all works in the corpus were carefully selected, added and reviewed by humans.
A knowledge base of the retained references was maintained, with a standardized processing of each entry (\Cref{app:prompts}).
It was used for indexing and interactive exploration of the corpus.
Language models were additionally used to polish writing.




\section{Background}\label{sec:background}
% Move background on conservation laws here ? Same for gauges
% useful to write borde2025,carlsson2025,chen2020,bronstein2021
\Cref{sec:what-is-sym} centers on symmetry, invariance and equivariance, without aiming at full rigor.
The following subsections introduce standard material used throughout: abstract groups (\Cref{sec:groups}), their realization on spaces (\Cref{sec:actions}), representations as the linear view of groups and the constraints they place on equivariant layers (\Cref{sec:representations}), and the infinitesimal view of continuous groups through Lie algebras (\Cref{sec:lie-algebra}).
The survey is not mathematically self-contained.

\subsection{What is symmetry?}\label{sec:what-is-sym}
Intuitively, a \emph{symmetry} of an object is a transformation that leaves it unchanged as a whole: it is a statement about identity.
\emph{Invariance} is a broader, functional term used to qualify a property of an object that does not change under a transformation (the property is called an \emph{invariant}).
For example, the area of a square is invariant under all isometries of the plane, yet only those superimposing the square onto itself (quarter rotations and reflections) are symmetries of the object ``square''.
The two notions can be related as follows: given a list of invariants of interest, a symmetry is an invertible transformation that leaves all of them unchanged.
The longer the list, the fewer the symmetries, and those of the object are recovered when the list determines the object.

\begin{Definition}[$\phi$-equivalence, {\citep{cohen2014}}]\label{def:equivalence}
    Let $\phi : \mathcal{X} \to \mathcal{Y}$ be a function from an input space $\mathcal{X}$ to an output space $\mathcal{Y}$.
    Two elements $x, x' \in \mathcal{X}$ are said to be \emph{$\phi$-equivalent}, denoted $x \sim_{\phi} x'$, if
    \[
        \phi(x) = \phi(x').
    \]
    The relation $\sim_{\phi}$ is an equivalence relation on $\mathcal{X}$ and induces a partition of $\mathcal{X}$ into equivalence classes.
\end{Definition}

In a deep learning context, the input space $\mathcal{X}$ can be, for instance, a vector space $\mathbb{R}^d$, a set of structured objects (images, graphs, point clouds), a set of functions or a set of parameters, while $\mathcal{Y}$ can be a boolean in $\{0,1\}$ (binary classification), a label in $\{1,\dots,C\}$ (multiclass classification), a scalar (regression) or another vector space $\mathbb{R}^e$ (feature representation).

\begin{Example}[Image segmentation]\label{ex:image-segmentation}
    If $\mathcal{X}=\mathbb{R}^{H\times W\times 3}$ is the set of RGB images, $\mathcal{Y}=\{1,\dots,C\}^{H\times W}$ the segmentation maps assigning a class to each pixel and $\phi$ the map assigning its segmentation map to each image, then the equivalence classes of $\sim_\phi$ are the sets of images with the same segmentation map.
    If instead $\mathcal{X}=H\times W$ is the set of pixel slots, $\mathcal{Y} = \{1,\dots,C\}$ and $\phi$ labels each pixel, then the equivalence classes of $\sim_\phi$ are the segments.
\end{Example}

\begin{Example}[Image classification]\label{ex:image-classification}
    If $\mathcal{X}$ is a \emph{set of images}, $\mathcal{Y} = \{1,\dots,C\}$ the class labels and $\phi$ the map assigning a label to each image, then the equivalence classes of $\sim_\phi$ are the sets of images of the same class.
\end{Example}

\begin{Example}[Loss level sets]
    If $\mathcal{X}=\Theta$ is a parameter space, $\mathcal{Y} = \mathbb{R}$ and $\phi=\mathcal{L}$ a loss function, then the equivalence classes of $\sim_{\mathcal{L}}$ are the level sets of the loss function.
\end{Example}

\begin{Example}[Observationally equivalent parameters]\label{ex:equivalence_observational}
    If $\mathcal{X}=\Theta$ is a parameter space, $\mathcal{Y}=\mathcal{F}$ a function space and $\phi=\varphi:\theta\mapsto f_\theta$ the implementation map sending parameters to the function the network computes (\Cref{eq:implementation-map}), then the equivalence classes of $\sim_\varphi$ are the sets of observationally equivalent parameters (see \Cref{sec:parameter-space-symmetries}).
\end{Example}

In the following, we use the concepts of group and group action loosely.
They are made precise in \Cref{def:group} and \Cref{def:group-action}.
Here it suffices to understand $g\in G$ as an operation moving a point $x\in\mathcal{X}$ to another point $g\cdot x\in\mathcal{X}$.

\begin{Definition}[Invariance]\label{def:invariance}
    Let $G$ be a group acting on a set $\mathcal{X}$.
    A function $\phi : \mathcal{X} \to \mathcal{Y}$ is $G$-invariant if $g\cdot x$ is $\phi$-equivalent to $x$ for all $g\in G$.
    That is, 
    \[
        \phi(g\cdot x)=\phi(x)\qquad\forall g\in G,\ x\in\mathcal{X}.
    \]
\end{Definition}

In \Cref{ex:image-classification}, $G$ acts on images and could contain translations, rotations or reflections, which leave the class unchanged, while in \Cref{ex:equivalence_observational}, $G$ acts on parameters and contains transformations that leave the implemented function unchanged such as within-layer neuron permutations for MLPs.
In \Cref{ex:image-segmentation}, by contrast, translating the image does not leave the segmentation unchanged in general.
The output is however not arbitrary: it is translated accordingly, which is captured by the notion of equivariance.

\begin{Definition}[Equivariance]\label{def:equivariance}
    Let $G$ be a group acting on sets $\mathcal{X}$ and $\mathcal{Y}$. 
    A function $f : \mathcal{X} \to \mathcal{Y}$ is $G$-equivariant if
    \[
        f(g \cdot x) = g \cdot f(x)
    \]
    for all $g \in G$ and $x \in \mathcal{X}$.
\end{Definition}

Now, $G$ acts on both input and output, possibly in the same way: translating an image translates its segmentation map, and rotating a molecule rotates the forces acting on it \citep{geiger2022,bogatskiy2020p3icml}.
However, equivariance does not require that the group $G$ acts in the same way on input and output \citep{cohen2016a,bogatskiy2020p3icml}: a single abstract element $g\in G$ may take one form on $\mathcal{X}$ and a different form on $\mathcal{Y}$.
To make the distinction explicit, two maps (\emph{actions}, see \Cref{def:group-action})
\[
    \rho_X : G \to \operatorname{Sym}(\mathcal{X}), \qquad \rho_Y : G \to \operatorname{Sym}(\mathcal{Y})
\]
are introduced to keep track of the form $g$ takes on each space.
Here $\operatorname{Sym}(\mathcal{S})$ denotes the group of bijections of a set $\mathcal{S}$.
The only requirement is that each form respects the group structure, i.e.\ that $\rho_X$ and $\rho_Y$ are group homomorphisms:
\begin{equation}\label{eq:representation-homomorphism}
    \rho_X(gh)=\rho_X(g)\circ\rho_X(h), \qquad \rho_Y(gh)=\rho_Y(g)\circ\rho_Y(h) \qquad \forall g, h \in G.
\end{equation}
The equivariance condition of \Cref{def:equivariance} then reads (omitting parentheses):
\begin{equation}\label{eq:general-equivariance}
    f\big(\rho_X(g)\,x\big)=\rho_Y(g)\,f(x) \qquad \forall g\in G,\ x\in\mathcal{X}.
\end{equation}
Invariance is the special case $\rho_Y(g)=\operatorname{id}_Y$ for all $g\in G$.
The elements $g$ satisfying \Cref{eq:general-equivariance} form the \emph{symmetry group} $\operatorname{Sym}(f)=\{g\in G: f\circ\rho_X(g)=\rho_Y(g)\circ f\}$ of $f$, a subgroup of $G$ (\Cref{sec:groups}), and $f$ is equivariant if and only if $\operatorname{Sym}(f)=G$.

\begin{Example}[One rotation, several forms]\label{ex:one-rotation-several-forms}
    Let $G=C_4=\{e,r,r^2,r^3\}$ be the group of quarter rotations acting on $\mathcal{X}=\mathbb R^2$ as $\rho_X(r)=R_{\frac{\pi}{2}}$.
    Depending on the task, input and output relate differently and the representation of the same element $r$ should reflect the particular nature of this relation:
    \begin{enumerate}
        \item if $f$ predicts an orientation $\hat \alpha\in\mathcal{Y}=[0,2\pi)$, a shift of angle $\rho_Y(r)\,\hat \alpha=\hat \alpha+\pi/2 \bmod 2\pi$
        \item if $f$ predicts an axis $\hat \alpha\in\mathcal{Y}=[0,\pi)$, a different shift $\rho_Y(r)\,\hat \alpha=\hat \alpha+\pi/2 \bmod \pi$
        \item if $f$ predicts an invariant class label, trivially $\rho_Y(r)=\operatorname{id}_Y$.
    \end{enumerate}
    In each case $\rho_Y(r)^4=\operatorname{id}_Y$, so $\rho_Y$ satisfies \Cref{eq:representation-homomorphism}, but the image $\rho_Y(G)$ varies (respectively $C_4$, $C_2$ and trivial in group language).
    \Cref{fig:one-rotation} illustrates the example with a lamp aiming a ray through the origin.
    The output action should (1) also be quarter rotations if $f:\mathcal{X}\to \mathcal{Y}$ models where the ray comes from (the lamp), (2) identify opposite rays if $f$ models rays themselves, and (3) be trivial if $f$ predicts an invariant, e.g.\ the ray's color.
\end{Example}

% wrapfigure: \begin{wrapfigure}[n] spans n lines; it must start a paragraph that is not near a page bottom
\begin{wrapfigure}[13]{r}{0.55\textwidth}
    \centering
    \vspace{-7mm}
    \includegraphics[width=0.4\textwidth]{images/one-rotation.pdf}
    \caption{\centering \Cref{ex:one-rotation-several-forms}, visualized.}
    \label{fig:one-rotation}
\end{wrapfigure}

Throughout the text, we write $g\cdot x$ or the shorter $gx$ whenever the action is unambiguous, for example when $G$ acts in the same way on $\mathcal{X}$ and $\mathcal{Y}$, and revert to the explicit $\rho_X(g)$, $\rho_Y(g)$ when the two actions must be distinguished explicitly.
Following usage, and contrary to the definition given at the beginning of this section, our use of the term ``symmetry'' will be general and refer to symmetry, invariance or equivariance.

\subsection{Groups}\label{sec:groups}

\begin{Definition}[Group]\label{def:group}
    A \emph{group} is a set $G$ with a binary operation $(g,h)\mapsto gh$ satisfying four axioms.
    \emph{Closure}: $gh\in G$ for all $g,h\in G$.
    \emph{Associativity}: $(gh)k=g(hk)$ for all $g,h,k\in G$.
    \emph{Identity}: there is $e\in G$ such that $eg=ge=g$ for all $g\in G$.
    \emph{Inverses}: every $g\in G$ has some $g^{-1}\in G$ such that $gg^{-1}=g^{-1}g=e$.
\end{Definition}

Dropping axioms leads to the weaker structures of \Cref{sec:extent-restriction}.
A subset $H\subseteq G$ is a \emph{subgroup}, written $H\le G$, if it contains $e$ and is closed under products and inverses, so that it is itself a group.
In particular, $\{e\}$ is the \emph{trivial} subgroup.
A map $\psi:G\to G'$ between groups is a \emph{homomorphism} if $\psi(gh)=\psi(g)\psi(h)$ for all $g,h\in G$, where it is left implicit that the binary operation depends on the group and can differ on both sides of the equality.
A bijective homomorphism is an \emph{isomorphism}, and isomorphic groups are written as $G\cong G'$.
An isomorphism from $G$ to itself is an \emph{automorphism}, and the automorphisms of $G$ form the group $\operatorname{Aut}(G)$ under composition.

\begin{Definition}[Direct and semidirect products]\label{def:products}
    The \emph{direct product} $G_1\times G_2$ of two groups is the set of pairs $(g_1,g_2)$ with the componentwise operation $(g_1,g_2)(g_1',g_2')=(g_1g_1',\,g_2g_2')$.
    Given groups $N$, $H$ and a homomorphism $\psi:H\to\operatorname{Aut}(N)$, the \emph{semidirect product} $N\rtimes_\psi H$ is the set of pairs $(n,h)\in N\times H$ with the operation
    \[
        (n_1,h_1)(n_2,h_2)=\big(n_1\,\psi(h_1)(n_2),\ h_1h_2\big).
    \]
\end{Definition}

The direct product is the semidirect product with trivial $\psi$, i.e.\ $\psi(h)=\operatorname{id}_N$ for all $h\in H$.
We omit $\psi$ when it is clear from context.

A \emph{Lie group} is a group that is also a smooth manifold, such that multiplication $(g,h)\mapsto gh$ and inversion $g\mapsto g^{-1}$ are smooth maps.
Its dimension is that of the manifold.
Finite and countable discrete groups are Lie groups of dimension $0$, and we call \emph{continuous} the groups of positive dimension.
A Lie group is \emph{compact} if it is compact as a topological space.

\begin{Example}[Groups in this survey]\label{ex:groups}
    \phantom{u}
    \begin{itemize}
        \item The \emph{cyclic group} $C_n=\{e,r,\dots,r^{n-1}\}$ of rotations of the plane by multiples of $2\pi/n$ is generated by one rotation $r$ with $r^n=e$, and $C_n\cong\mathbb{Z}/n\mathbb{Z}$.
        \item The \emph{general linear group} $GL(n,\mathbb{R})$, or $GL(n)$, of invertible $n\times n$ real matrices under multiplication, is an open subset of $\mathbb{R}^{n\times n}$ hence a non-compact Lie group of dimension $n^2$.
        \item The \emph{symmetric group} $S_n=\operatorname{Sym}(\{1,\dots,n\})$ of permutations under composition is finite of order $n!$. For $\sigma \in S_n$ the permutation matrix $P_\sigma$ is defined by $P_\sigma e_i=e_{\sigma(i)}$.
        The map $\sigma\mapsto P_\sigma$ is an injective homomorphism from $S_n$ to $GL(n,\mathbb{R})$, so $S_n$ is isomorphic to the group of permutation matrices, a subgroup of $GL(n,\mathbb{R})$.
        \item The \emph{orthogonal group} $O(n)=\{A\in GL(n):A^\top A=I\}$ of linear isometries of $\mathbb{R}^n$ (rotations including reflections), and its subgroup $SO(n)=\{A\in O(n):\det A=1\}$ (rotations only) are compact Lie groups of dimension $n(n-1)/2$.
        \item The \emph{Euclidean group} $E(n)=\mathbb{R}^n\rtimes O(n)$ and its subgroup $SE(n)=\mathbb{R}^n\rtimes SO(n)$ of rigid motions. With elements $(v,R)$ and $\psi(R)v=Rv$, the product is $(v_1,R_1)(v_2,R_2)=(v_1+R_1v_2,\,R_1R_2)$. $E(n)$ is the group of all isometries of $\mathbb{R}^n$ with the map $x\mapsto Rx+v$ (see actions below) and both are non-compact Lie groups of dimension $n(n+1)/2$.
        \item In parameter space (\Cref{sec:catalog-param-space-symmetries}), the permutations and positive rescalings of the $d$ hidden neurons of a ReLU layer form $(\mathbb{R}_{>0})^{d}\rtimes S_{d}$, with $S_d$ permuting the scaling factors.
        Furthermore, the groups of different layers commute and combine as a direct product.
    \end{itemize}
\end{Example}

\subsection{Actions}\label{sec:actions}

\begin{Definition}[Group action]\label{def:group-action}
    An \emph{action} of a group $G$ on a set $\mathcal{X}$ is a homomorphism $\rho:G\to\operatorname{Sym}(\mathcal{X})$, and we write $g\cdot x:=\rho(g)\,x$ or the short form $gx$.
\end{Definition}

In continuation with \Cref{ex:groups}, $E(n)$ acts on $\mathbb{R}^n$ by $(v,R)\cdot x=Rx+v$ and $S_n$ on $n$ $d$-dimensional points by $\sigma\cdot X=P_\sigma X$ with $X\in\mathbb{R}^{n\times d}$.
The same group can act in multiple ways: $S_n$ acts on graphs through the adjacency matrix $A$ by $\sigma\cdot A=P_\sigma AP_\sigma^\top$.

\begin{Definition}[Orbit and stabilizer]\label{def:orbit-stabilizer}
    Let $G$ act on $\mathcal{X}$ and $x\in\mathcal{X}$.
    The \emph{orbit} of $x$ is $G\cdot x=\{g\cdot x: g\in G\}\subseteq\mathcal{X}$ and its \emph{stabilizer} is $G_x=\operatorname{Stab}(x)=\{g\in G: g\cdot x=x\}$.
\end{Definition}

The stabilizer $G_x$ is a subgroup of $G$ and orbits partition $\mathcal{X}$ since for $x,y\in \mathcal{X}$ the relation $x\sim y$ defined by $y\in G\cdot x$ is an equivalence relation.
A function is $G$-invariant (\Cref{def:invariance}) if and only if it is constant on orbits.

\begin{Definition}[Quotient and canonicalization]\label{def:quotient}
    The \emph{quotient} $\mathcal{X}/G$ is the set of orbits, and the \emph{quotient map} $\pi:\mathcal{X}\to\mathcal{X}/G$, $x\mapsto G\cdot x$, sends a point to its orbit.
    A \emph{section} of $\pi$ is a map $s:\mathcal{X}/G\to\mathcal{X}$ with $\pi\circ s=\operatorname{id}_{X/G}$, i.e.\ it picks one representative in each orbit, and $c=s\circ\pi:\mathcal{X}\to\mathcal{X}$ is the associated \emph{canonicalization}.
\end{Definition}

The quotient map is $G$-invariant, and every $G$-invariant function $f:\mathcal{X}\to\mathcal{Y}$ factors uniquely as $f=\bar f \circ\pi$ with $\bar f:\mathcal{X}/G\to\mathcal{Y}$, that is, invariant functions on $\mathcal{X}$ are functions on $\mathcal{X}/G$.
The quotient can also be written with the orbit relation such as in parameter space $\Theta/\!\sim$ (\Cref{sec:quotient-repr-align}).
A concrete example of canonicalization is centering a point cloud, $X\mapsto X-\boldsymbol{1}\bar{x}^\top$ with $\bar{x}$ its centroid.

\begin{Theorem}[Orbit--stabilizer]\label{thm:orbit-stabilizer}
    Let $G$ act on $\mathcal{X}$ and $x\in\mathcal{X}$.
    The map $gG_x\mapsto g\cdot x$ is a well-defined bijection from $G/G_x$ onto the orbit $G\cdot x$.
\end{Theorem}

A larger stabilizer thus means a smaller orbit.
If $G$ is finite, $|G\cdot x|=|G|/|G_x|$.
If it is a Lie group, we have instead $\dim(G\cdot x)=\dim G-\dim G_x$.
For example, $SO(3)$ acts on $\mathbb{R}^3$ by rotations.
For $x\neq0$, the orbit is the sphere of radius $\lVert x\rVert$, formally $G\cdot x\cong S^2$. 
The stabilizer $G_x\cong SO(2)$ contains the rotations about the axis of $x$ and we have $\dim S^2=3-1$.
The origin is special: it has orbit $\{0\}$ and stabilizer the full $SO(3)$ group.

\begin{Definition}[Free, faithful and transitive actions]\label{def:free-transitive}
    An action $\rho$ of $G$ on $\mathcal{X}$ is
    \begin{itemize}
        \item \emph{free} if all stabilizers are trivial,
        \item \emph{faithful} if $\rho$ is injective, i.e.\ only $e$ fixes every point,
        \item \emph{transitive} if for any $x\in \mathcal{X}$ the orbit is $\mathcal{X}$, and $\mathcal{X}$ is said to be a \emph{homogeneous space} of $G$.
    \end{itemize}
\end{Definition}

A free action on a nonempty set is faithful, and a faithful action identifies $G$ with the subgroup $\rho(G)\le\operatorname{Sym}(\mathcal{X})$.
By \Cref{thm:orbit-stabilizer}, each orbit of a free action is in bijection with $G$.
For instance, $E(n)$ acts transitively but not freely on $\mathbb{R}^n$, while $S_n$, by permuting coordinates, acts faithfully on $\mathbb{R}^n$ but neither freely nor transitively.\newline
A group acts transitively and freely on itself by left multiplication and transitively on $G/H$ by $g\cdot(kH)=(gk)H$.
For a homogeneous space $\mathcal{X}$ and any $x_0\in \mathcal{X}$, $\mathcal{X}\cong G/G_{x_0}$, for example $S^2\cong SO(3)/SO(2)$ and $\mathbb{R}^n\cong E(n)/O(n)$.

\subsection{Representations}\label{sec:representations}

The maps $\rho_X,\rho_Y$ of \Cref{eq:representation-homomorphism} are sometimes loosely called representations.
A representation is almost always intended as a linear action.
This is the natural setting for networks whose layers are linear maps between feature spaces, which we use as illustration in this part.

\begin{Definition}[Linear representation]\label{def:representation}
    A \emph{(linear) representation} of a group $G$ on a real vector space $V$ is a homomorphism $\rho:G\to GL(V)$, i.e.\ an action of $G$ on $V$ (\Cref{def:group-action}) by invertible linear maps.
\end{Definition}

Given representations $\rho_V$ and $\rho_W$ of $G$ on $V$ and $W$, a linear map $\Lambda:V\to W$ is $G$-equivariant (\Cref{def:equivariance}) if and only if
\begin{equation}\label{eq:linear-equivariance}
    \Lambda\,\rho_V(g)=\rho_W(g)\,\Lambda\qquad\forall g\in G.
\end{equation}
This constraint is linear in $\Lambda$, so equivariant linear maps form a vector subspace whose dimension is the number of free parameters of an equivariant linear layer.

\begin{Example}[Trivial, permutation and regular representations]\label{ex:representations-recap}
    \phantom{u}
    \begin{itemize}
        \item With the \emph{trivial representation} $\rho_W(g)=\operatorname{id}_W$, \Cref{eq:linear-equivariance} states that $\Lambda$ is invariant.
        \item With the \emph{permutation representation} $\rho_V(g)=\rho_W(g)=P_{\sigma_g}$, \Cref{eq:linear-equivariance} reads $\Lambda_{\sigma_g(i)\sigma_g(j)}=\Lambda_{ij}$ and weights are tied along the orbits of $G$ on index pairs $(i,j)$.
        \item The \emph{regular representation} of a finite group $G$ is the permutation representation associated with left multiplication of $G$ on itself: on the space $\mathbb{R}^G$ of functions $f:G\to\mathbb{R}$, $(\rho_{\mathrm{reg}}(g)f)(h)=f(g^{-1}h)$. The $G$-convolution of \Cref{sec:edl-achieving-equiv} satisfies \Cref{eq:linear-equivariance} for it.
    \end{itemize}
\end{Example}

\begin{Definition}[Irreducibility]\label{def:irreducible}
    Let $\rho$ be a representation on a vector space $V\neq\{0\}$.
    A subspace $U\subseteq V$ is \emph{invariant} if $\rho(g)U\subseteq U$ for all $g\in G$, and $\rho$ then restricts to a representation on $U$.
    If its only invariant subspaces are $\{0\}$ and $V$, $\rho$ is said to be \emph{irreducible}.
\end{Definition}

Every finite-dimensional continuous representation of a compact group is a direct sum of irreducible ones, which thus serve as building blocks.
For example, the permutation representation of $S_n$ splits into the irreducible subspaces $\operatorname{span}(\boldsymbol{1})$ and $\{x\in\mathbb{R}^n:\sum_ix_i=0\}$.

\emph{Steerable} networks \citep{cohen2016a,weiler2019anips} extend linear equivariance to typed features transforming according to a chosen representation.
For the domain $\Omega=\mathbb{R}^d$ and $G=\mathbb{R}^d\rtimes H$ with $H\le O(d)$, a \emph{feature field of type} $\rho:H\to GL(C)$ is a map $f:\mathbb{R}^d\to\mathbb{R}^C$ on which $G$ acts by
\begin{equation}\label{eq:feature-field}
    \big((v,h)\cdot f\big)(u)=\rho(h)\,f\big(h^{-1}(u-v)\big).
\end{equation}
Points $u\in\Omega$ move as $u\mapsto hu+v$ (as Euclidean actions in \Cref{sec:actions}) and the $C$ channel values are mixed by $\rho(h)$.
The trivial type gives scalar fields (e.g.\ grayscale images), the type $\rho(h)=h$ gives vector fields (e.g.\ forces), and for finite $H$ the regular type recovers $G$-CNNs.

A convolution $f_{\mathrm{out}}(u)=\int_{\mathbb{R}^d}K(u-u')\,f_{\mathrm{in}}(u')\,du'$ with kernel $K:\mathbb{R}^d\to\mathbb{R}^{C_{\mathrm{out}}\times C_{\mathrm{in}}}$ is translation equivariant, and it maps fields of type $\rho_{\mathrm{in}}$ to fields of type $\rho_{\mathrm{out}}$ $G$-equivariantly if and only if $K$ satisfies the \emph{kernel constraint} \citep{weiler2018anips,cesa2021}
\begin{equation}\label{eq:steerable-constraint}
    K(hy)=\rho_{\mathrm{out}}(h)\,K(y)\,\rho_{\mathrm{in}}(h)^{-1}\qquad\forall h\in H,\ y\in\mathbb{R}^d,
\end{equation}
i.e.\ $K$ is itself $H$-equivariant, with $H$ acting on matrices by $M\mapsto\rho_{\mathrm{out}}(h)M\rho_{\mathrm{in}}(h)^{-1}$.
Ordinary CNNs, whose kernels are unconstrained, correspond to the case $H=\{e\}$, while for $H=SO(d)$ and scalar fields the constraint $K(hy)=K(y)$ must be respected, leading to isotropic kernels.
Being linear in $K$, \Cref{eq:steerable-constraint} is solved once for a basis $(K_k)_k$ and only the coefficients of $K=\sum_kw_kK_k$ are learned.
Decomposing $\rho_{\mathrm{in}}$ and $\rho_{\mathrm{out}}$ into irreducible representations splits the constraint into small independent blocks \citep{weiler2019anips,cesa2021}.
The gauge-equivariant convolution of \Cref{sec:edl-achieving-equiv} obeys the same constraint, with $H$ the structure group.

\subsection{Lie groups and Lie algebras}\label{sec:lie-algebra}

A Lie group (\Cref{sec:groups}) is described to first order by its tangent space at the identity, called its Lie algebra.
We restrict to matrix Lie groups $G\le GL(n,\mathbb{R})$, which include those of \Cref{ex:groups}.

\begin{Definition}[Lie algebra and exponential map]\label{def:lie-algebra}
    The \emph{Lie algebra} $\mathfrak{g}\subseteq\mathbb{R}^{n\times n}$ of $G$ is the vector space of velocities $\gamma'(0)$ of smooth curves $\gamma$ in $G$ with $\gamma(0)=I$, and it is closed under the \emph{Lie bracket} $[A,B]=AB-BA$.
    The \emph{exponential map} $\exp:\mathfrak{g}\to G$ is the matrix exponential $\exp(A)=\sum_{k\ge0}A^k/k!$, and $t\mapsto\exp(tA)$ is the \emph{one-parameter subgroup} generated by $A$: $\exp((s+t)A)=\exp(sA)\exp(tA)$, with velocity $A$ at $t=0$.
\end{Definition}

For example, differentiating $\gamma(t)^\top\gamma(t)=I$ at $t=0$ gives $A^\top+A=0$: the Lie algebra $\mathfrak{so}(n)$ of $SO(n)$ is the space of antisymmetric matrices, of dimension $n(n-1)/2$, the same as $SO(n)$.
For $n=2$ it is spanned by $J=\left(\begin{smallmatrix}0&-1\\1&0\end{smallmatrix}\right)$, and $\exp(\alpha J)$ is the rotation by $\alpha$.
Near $0$, $\exp$ is a diffeomorphism onto a neighborhood of $I$, with local inverse $\log$.
It is surjective for compact connected $G$ but not in general.

\begin{Definition}[Infinitesimal generators]\label{def:infinitesimal-generators}
    The \emph{infinitesimal representation} of a smooth representation $\rho:G\to GL(V)$ is the linear map $d\rho:\mathfrak{g}\to\operatorname{End}(V)$,
    \[
        d\rho(A)=\frac{d}{dt}\rho\big(\exp(tA)\big)\Big\vert_{t=0},
    \]
    which satisfies $\rho(\exp A)=\exp(d\rho(A))$ and $d\rho([A,B])=[d\rho(A),d\rho(B)]$.
    For a basis $A_1,\dots,A_k$ of $\mathfrak{g}$, the $A_i$, or their images $d\rho(A_i)$, are called \emph{infinitesimal generators}.
\end{Definition}

Since $\exp(\mathfrak{g})$ generates $G$ when $G$ is connected, the constraint over $G$ reduces to $k$ linear ones: a linear map $\Lambda$ satisfies \Cref{eq:linear-equivariance} if and only if
\begin{equation}\label{eq:lie-algebra-equivariance}
    d\rho_W(A_i)\,\Lambda=\Lambda\,d\rho_V(A_i),\qquad i=1,\dots,k.
\end{equation}
% Necessity follows by differentiating $\rho_W(\exp tA)\Lambda=\Lambda\rho_V(\exp tA)$ at $t=0$, sufficiency by exponentiating.
A disconnected group needs additional constraints, e.g.\ a reflection for $O(n)$.

For a nonlinear map $f:\mathcal{X}\to\mathcal{Y}$, the first-order counterpart is the \emph{Lie derivative}
\[
    (\mathcal{L}_Af)(x)=\frac{d}{dt}\,\rho_Y\big(\exp(tA)\big)^{-1}f\big(\rho_X(\exp(tA))\,x\big)\Big\vert_{t=0},
\]
equal to $\partial f(x)\,d\rho_X(A)\,x-d\rho_Y(A)\,f(x)$ for linear actions, with $\partial f$ the Jacobian.
For connected $G$, $f$ is equivariant if and only if $\mathcal{L}_{A_i}f=0$ for all $i$, so the Lie derivative measures local equivariance errors (\Cref{sec:sym-quantification}).
%  and penalizes relaxed models (\Cref{sec:approx-relaxed}).
% Learning generators rather than the group underlies continuous symmetry discovery (\Cref{sec:sym-disc-from-data,sec:disco-identification}), and one-parameter subgroups acting on parameters yield conservation laws (\Cref{sec:conservation-laws}).

\section{Symmetry as a Design Principle}\label{sec:design-principle}

This section gathers the methods in which a symmetry, often \emph{known} in advance, is built into a learning system.
Here symmetry is a muse: it is recognized and taken into account to guide design, out of the desire to ``do things right''.
The purpose therefore tends to be prescriptive, as best illustrated by \Cref{sec:gdl} on Equivariant Deep Learning, which proposes ways to distill a ground-truth problem symmetry into the hypothesis class.
This contrasts with \Cref{sec:parameter-space-symmetries}, where symmetries are rather unanticipated artifacts emerging from a parametrization.

The section has three subsections.
The first two, Equivariant Deep Learning (EDL, \Cref{sec:gdl}) and Data Augmentation (DA, \Cref{sec:data-augmentation}), both seek to distill problem symmetries.
They roughly split along the explicitness axis, with EDL tending to be more formalized and group theoretic than DA.
\emph{Equivariant architectures} encode the symmetry \emph{explicitly} in the hypothesis class, so that invariance or equivariance holds by construction for every parameter setting.
Architecture-agnostic, DA instead enforces the symmetry \emph{implicitly} through input-space transformations $\mathcal X \to \mathcal X$, encouraging but not guaranteeing invariance through orbit-averaged training.
Note that these transformations can be expressed (implemented) explicitly by the designer, but this differs from the intended meaning of explicitness in \Cref{tab:taxonomy}.
Finally, \emph{framework symmetries} (\Cref{sec:framework-symmetries}) differ on the object axis: they concern the training procedure itself, involving several forward calls (joint embeddings, \Cref{sec:joint-embedding}), relying on assumptions about how the data stream is presented (data ordering, \Cref{sec:data-ordering}), or arising from the way the learning framework is set up rather than from a property of the problem (criterion symmetries, \Cref{sec:criterion-symmetries}).


\subsection{Geometric and Equivariant Deep Learning}\label{sec:gdl}
% dehmamy2021anips: achieve equivariance through Lie algebra
% figure: draft/orbit-stability.png to reproduce or illustrate conceptually (equivariant models guarantee results don't change over an orbit)
% Graphs, huang2023: distinguishes active and passive graph symmetries in constrat to CNNs
\emph{Geometric Deep Learning} (GDL) \citep{bronstein2017ispm,bronstein2021} starts with a domain $\Omega$ and signals $x\in\mathcal{X}(\Omega,\mathbb{R}^C)$. A group $G$ acting on $\Omega$ by $u\mapsto g\cdot u$ acts on signals by $(g\cdot x)(u)=x(g^{-1}\cdot u)$.
The structure of $\Omega$, in particular its symmetries, is analyzed following the legacy of Klein's Erlangen Program.
A central idea is that an appropriate inductive bias for a learning problem comes from respecting the symmetries of $\Omega$.
Symmetries, together with a scale-separation principle, are GDL's answer to the curse of dimensionality.
The scope of GDL also includes unifying architectures in a common blueprint, e.g.\ Transformers as graph neural networks \citep{joshi2025}, and extending non-Euclidean geometry from the input domain to every stage of deep learning, e.g.\ non-Euclidean representation spaces \citep{mettes2024ijcv,saezdeocarizborde2025iclr}.
In this survey, we focus on its \emph{Equivariant Deep Learning} core, which is the most concerned with symmetries.
Standard GDL mostly assumes exact symmetry, \textit{i.e.}\ equivariance to every $g\in G$ acting on $\Omega$.
The deformation stability perspective \citep[Sec.~3.3]{bronstein2021} extends control to deformations $\tau\in\operatorname{Diff}(\Omega)$ near $G$, close to the setting of \Cref{sec:extent-restriction}.

The promise of EDL is to design principled architectures that are sample-efficient and generalize well.
Equivariant networks typically benefit from fewer learnable parameters and better interpretability, \textit{e.g.}\ CNN filters.
They also offer robustness guarantees, since their response to a transformed input is dictated by the group: $f(\rho_X(g)\,x)=\rho_Y(g)\,f(x)$ for all $g\in G$.
As for generalization, an out-of-distribution input $\rho_X(g)\,x$, with $x$ in-distribution, is handled as reliably as $x$ itself once the symmetry is imposed.
Resorting to EDL is not automatic, however: \Cref{sec:edl-tradeoffs} discusses the trade-offs between EDL and DA, in particular the recurring question of whether the cost of exact equivariance is justified over a simpler model trained with augmentation.
Recommended resources include surveys \citep{rath2024air,basheer2025ijdsa}, foundations \citep{borde2025,gerken2021,carlsson2025} and a historical account of the field \citep{bronstein2023tg}.

\subsubsection{Achieving equivariance}\label{sec:edl-achieving-equiv}

Methods to design equivariant architectures generally intertwine equivariant linear maps with suitable non-linearities.

\emph{Group Convolutional Neural Networks} ($G$-CNNs) \citep{cohen2016} generalize traditional CNNs from translations to discrete groups (\textit{e.g.}\ $90^\circ$ rotations and reflections), and continuous groups must be discretized.
The processing follows three steps.
First, a lifting convolution maps a function on the domain $\Omega$ to a function on $G$.
Then, iterated $G$-convolutions $(f\star\kappa)(g)=\sum_{h\in G}f(h)\,\kappa(g^{-1}h)$ compute feature maps $f^{(l)}:G\to\mathbb{R}^{c_l}$.
Finally, the result is aggregated: pooled over $G$ for invariance, or projected back to $\Omega$ for equivariance.
\citet{cohen2019anips} unify $G$-CNNs on homogeneous spaces.
$G$-CNNs have been extended to Lie groups: \citet{finzi2020p3icml} parametrize kernels on the Lie algebra, covering any Lie group with a surjective exponential map and arbitrary continuous data, while \citet{bekkers2021} expand kernels in B-splines via the logarithmic map for grid data.
Scale-translation groups are also handled by $G$-CNNs \citep{sosnovik2020,romero2023tmlr}.

\emph{Steerable CNNs} \citep{cohen2016a} generalize $G$-CNNs, which they recover with regular representations.
Feature maps stay on the spatial domain as typed fields, transforming under a representation $\rho$ of the stabilizer $H$.
Kernels solve the steerable constraint of \Cref{eq:steerable-constraint} in a truncated equivariant basis, avoiding discretization of continuous groups.
\citet{cesa2021} extend them to $\mathbb{R}^d\rtimes H$ for any compact $H\le O(d)$.

Relatedly, \citet{finzi2021p3icml} numerically solve the linear equivariance constraint with an SVD decomposition to build equivariant MLPs for arbitrary matrix groups.
For finite groups acting by permutations, equivariance amounts to weight-tying: a layer is $G$-equivariant if $G$ preserves its parameter-sharing pattern \citep{ravanbakhsh2017p3icml}.
\emph{Clifford (geometric) algebras} yield equivariant models such as the $E(n)$-equivariant CGENN \citep{ruhe2023} and the $E(3)$-equivariant transformer GATr \citep{brehmer2023}.

Other approaches include \emph{invariant-feature construction} from distances or inner products, \textit{e.g.}\ EGNN \citep{satorras2021p3icml} and ``scalars are universal'' \citep{villar2021anips}.
\emph{Loss regularization} keeps the group but loosens the constraint, yielding only approximate (soft) equivariance (\Cref{sec:soft-equivariance}).
\emph{Symmetrization} wraps an arbitrary backbone $f$, \textit{e.g.}\ by canonicalization \citep{kaba2023p4icml} or group/frame averaging \citep{puny2021} with the Reynolds operator $(\mathcal{R}f)(x)=\frac{1}{|G|}\sum_{g\in G}\rho_Y(g)^{-1}f(\rho_X(g)\,x)$, which maps any $f$ to an equivariant one (an integral replaces the sum for compact $G$).

\paragraph{Euclidean symmetries}
Euclidean symmetries are important symmetries of the physical world, consisting of translations and rotations (with or without reflections).
Mathematically, they form the Euclidean group $E(3)$ and its subgroup $SE(3)$ of rigid motions (\Cref{ex:groups}).
Tools to treat these symmetries include irreducible representations, Fourier analysis, and spherical harmonics (functions on $S^2$ that generalize the Fourier basis on $S^1$).
We proceed from coarse to fine rotations, in 2D then 3D.

In 2D, seminal $G$-CNNs \citep{cohen2016} already treat $90^\circ$ rotations and reflections.
\citet{hoogeboom2018} reach $60^\circ$ rotations without interpolation on hexagonal grids.
\citet{weiler20182iccvpr} use finer rotation discretizations, bridging $G$-CNNs and steerable networks.
Continuous rotations are handled with circular harmonics by \citet{worrall2017}, and the full group $E(2)$ by \citet{weiler2019anips}.
In 3D, discrete rotation groups act on voxels, \textit{e.g.}\ for biomedical data \citep{worrall2018cv-e21ecmgs82ppv,winkels2019mia}.
\emph{Spherical CNNs} \citep{cohen2018} handle continuous rotations of signals on $S^2$ using a generalized Fast Fourier Transform.
\citet{monti20172iccvprc} sidestep extrinsic rotations via intrinsic surface coordinates.

Continuous $SE(3)$-equivariance mostly relies on irreducible representations (irreps) of $SO(3)$, an idea explored early by \citet{cohen2014} for commutative groups.
\citet{weiler2018anips} build $SE(3)$-equivariant convolutions on voxels with spherical-harmonic kernels.
\emph{Tensor field networks} \citep{thomas2018} extend this to point clouds, combining irreps features with spherical-harmonic filters via tensor products.
The e3nn library \citep{geiger2022} implements these ideas.
\emph{Equiformer} is an irreps-based transformer for atomistic graphs \citep{liao2022,liao2023}.

\paragraph{Permutation: sets, point clouds and graphs}\label{sec:edl-permutation}
When data is unstructured (sets), spatially unorganized (point clouds), or defined on a graph, models must be invariant or equivariant to the order of input elements.
We start with sets and point clouds since they have weaker structure.
Stacking $n$ elements as the rows of $X\in\mathbb{R}^{n\times d}$, $\sigma\in S_n$ acts by $X\mapsto P_\sigma X$, with $P_\sigma$ the permutation matrix: set-level predictions must be invariant, $f(P_\sigma X)=f(X)$, and per-element ones equivariant, $f(P_\sigma X)=P_\sigma f(X)$.
\citet{zaheer2018} introduced \emph{Deep Sets}, proving that invariant functions on a countable domain decompose as $f(X)=\psi\big(\sum_{x\in X}\phi(x)\big)$ and characterizing equivariant linear layers (\Cref{ex:representations-recap}).
The \emph{Set Transformer} \citep{lee2019p3icml} models element interactions with self-attention (equivariant without positional encodings) followed by invariant attention-based pooling.
\citet{wang2025} make a more abstract use of $S_n$-equivariance in multi-view 3D reconstruction, where treating views as a set removes the bias of an arbitrary reference view.
PointNet \citep{qi2017a} pioneered deep learning directly on raw point clouds, though only permutation invariance is built in: translation invariance comes from normalization, and approximate rotation invariance from learned canonicalization plus data augmentation. 
PointNet++ \citep{qi2017} applies PointNet hierarchically to local neighborhoods to capture multi-scale structure.
\citet{wang2019} bridge point clouds and graphs: their DGCNN recomputes at each layer a $k$-nearest-neighbor graph in feature space and updates each point as $x_i'=\max_{j\in\mathcal{N}_k(i)}h_\theta(x_i,x_j-x_i)$.
Since this graph is built from features rather than indices, the layer is permutation-equivariant.
It is a message passing layer, the central tool on graphs.

The graph domain and \emph{Graph Neural Networks} (GNNs) are a prolific spearhead of EDL, especially for molecular applications.
On graphs, $S_n$ relabels nodes, acting on adjacency and node features as $\sigma\cdot(A,X)=(P_\sigma AP_\sigma^\top,P_\sigma X)$.
The message passing framework \citep{gilmer2017} unified earlier graph models \citep{duvenaud2015,kipf2017,bruna2014} into a single formalism of message, aggregation and update functions: aggregating messages invariantly makes each layer equivariant, and an invariant readout yields graph-level predictions.
\citet{velickovic2018} incorporated attention into this scheme while preserving permutation equivariance.
For \emph{geometric graphs}, where nodes carry coordinates in $\mathbb{R}^3$, permutation equivariance must be combined with Euclidean symmetries (see above).
\citet{satorras2021p3icml} achieve joint $E(n)$ and permutation equivariance through invariant edge features based on interpoint distances, while \citet{fuchs2020} attain $SE(3)$-equivariance via self-attention on steerable features, at higher computational cost but with greater expressivity.

Message passing propagates information along edges, i.e.\ between pairs of nodes.
\emph{Topological Deep Learning} (TDL) \citep{papillon2023,papamarkou2024} extends this permutation-equivariant design to higher-order relational structures (simplicial, cell and combinatorial complexes, hypergraphs), passing messages between cells of different dimensions while remaining equivariant to cell relabeling.
The geometric and topological extensions compose: \citet{battiloro2024} construct $E(n)$-equivariant topological neural networks.

\paragraph{Gauges and local symmetries}
% Introduction, intuition and motivation for why gauges
Not all symmetries arise from isometries of the space $\mathcal X$: some come from the way we represent computations.
In general, the tangent spaces of a $d$-dimensional manifold $M$ do not come with a canonical basis.
Computing a convolution in coordinates requires fixing one, a \emph{frame}, at each point: an arbitrary choice to which we want to be equivariant.
This motivation has led to \emph{coordinate-independent CNNs}, ``generalizations of steerable CNNs from Euclidean (or homogeneous) spaces to Riemannian manifolds'' \citep{weiler2021}, which are equivariant to these passive and local transformations: no point moves, only the frame changes, independently at each point.

% Formal introduction
Formally, a \emph{gauge} on an open set $U\subseteq M$ is a smooth family of frames $w_p:\mathbb{R}^d\to T_pM$, $p\in U$, and a \emph{gauge transformation} $\gamma:U\to GL(d)$ changes it into $w_p\,\gamma(p)$, so that coordinates of tangent vectors change as $v\mapsto\gamma(p)^{-1}v$.
Restricting to \emph{allowed frames} (e.g.\ orthonormal ones) restricts $\gamma$ to a \emph{structure group} $H\leq GL(d)$ (e.g.\ $O(d)$).
The smaller $H$, the more structure is fixed, down to a global frame field for $H=\{e\}$, which topology may forbid (e.g.\ on $S^2$).
Features of type $\rho:H\to GL(C)$, where $C$ and $\rho$ depend on the feature (scalar, vector, \dots), have coefficients transforming as $f(p)\mapsto\rho(\gamma(p))^{-1}f(p)$.
As features at different points live in unrelated frames, a convolution first parallel transports $f(q)$ to $p$, through $\rho(g_{p\leftarrow q})$ with $g_{p\leftarrow q}\in H$, then applies a kernel $K:\mathbb{R}^d\to\mathbb{R}^{C_{\mathrm{out}}\times C_{\mathrm{in}}}$ shared across $M$:
\[
    (K\star f)(p)=\int_{\mathbb{R}^d}K(v)\,\rho_{\mathrm{in}}(g_{p\leftarrow q_v})\,f(q_v)\,dv,\qquad q_v=\exp_p(w_pv).
\]
Here $\exp_p$ is the exponential map at $p$: $q_v$ is reached by following the geodesic from $p$ with initial velocity $w_pv$.
For $H\leq O(d)$, this convolution is gauge equivariant if and only if $K(h^{-1}v)=\rho_{\mathrm{out}}(h)^{-1}K(v)\,\rho_{\mathrm{in}}(h)$ for all $h\in H$, equivalent to the kernel constraint of \Cref{eq:steerable-constraint}.
% Resources
Detailed material can be found in the survey by \citet{gerken2021} and the comprehensive work by \citet{weiler2021,weiler2026}.

% Research
The need for gauge equivariance already surfaces with global symmetries.
Consider a scalar field on $S^2$, and the problem of $SO(3)$-equivariance.
\citet{esteves2018} stay on $S^2$ and perform \emph{spherical convolutions} with filters $\kappa:S^2\to\mathbb{R}$:
\[
    (f\star \kappa)(x)=\int_{SO(3)} f(g\eta)\,\kappa(g^{-1}x)\,dg,
\]
where $\eta$ is a reference point (the north pole).
The filters are, however, isotropic (\emph{zonal} in their terminology).
From the gauge perspective, $H=SO(2)$ with scalar features, and the kernel constraint reduces to $K(hv)=K(v)$ for all $h\in SO(2)$ and $v\in\mathbb{R}^2$: kernels are invariant under rotations about their center and cannot detect orientation (e.g.\ no oriented edge detectors).
Instead, \citet{cohen2018} lift from $S^2$ to $SO(3)$:
\[
    [\psi \star f](R) = \langle L_R \psi, f \rangle = \int_{S^2} \sum_{k=1}^{C} \psi_k(R^{-1}x)\, f_k(x)\, dx,
\]
avoiding the restriction at a higher computational cost.
The follow-up of \citet{esteves2020anips} obtains anisotropic kernels while staying on $S^2$ with spin-weighted spherical functions, a spectral realization of gauge-equivariant CNNs on $S^2$ (without the gauge formalism).
The seminal work of \citet{cohen2019p3icml} introduced the theory of gauge-equivariant CNNs on manifolds, with the Icosahedral CNN ($H=C_6$) as a concrete instantiation.
Still with anisotropy in mind, \citet{haan2020} generalize Graph Convolutional Networks (GCNs) to meshes.
A standard GCN layer $(K\star f)_p=K_{\mathrm{self}}f_p+\sum_{q\in\mathcal{N}_p}K_{\mathrm{neigh}}f_q$ is isotropic, with constant matrices $K_{\mathrm{self}},K_{\mathrm{neigh}}$.
Their \emph{Gauge Equivariant Mesh CNNs} generalize the convolution above as
\[
    (K \star f)_p = K_{\mathrm{self}}\, f_p + \sum_{q \in \mathcal{N}_p} K_{\mathrm{neigh}}(\alpha_{pq})\, \rho_{\mathrm{in}}(g_{p\leftarrow q})\, f_q,
\]
where $\alpha_{pq}$ is the angle between neighbor $q$ and an arbitrary reference neighbor (the gauge), and $g_{p\leftarrow q}\in SO(2)$ is the discrete parallel transporter.
Parametrizing $SO(2)$ by angles $h$, the kernel constraint reads $K_{\mathrm{neigh}}(\alpha-h)=\rho_{\mathrm{out}}(-h)\,K_{\mathrm{neigh}}(\alpha)\,\rho_{\mathrm{in}}(h)$ and $K_{\mathrm{self}}=\rho_{\mathrm{out}}(-h)\,K_{\mathrm{self}}\,\rho_{\mathrm{in}}(h)$.
\citet{weiler2021} (see also \citep{weiler2026}) deepen the theory of coordinate independent CNNs: they prove that coordinate independence and weight sharing imply gauge equivariance for a general $H\leq GL(d)$, and that the resulting convolutions are equivariant to the isometries preserving the $H$-structure (a subgroup of the space symmetries), linking local and global transformations.
They also extensively review the literature and classify previous works within their formalism.
Reinterpreting isomorphisms of graphs and neighborhoods as gauge transformations, \citet{dehaan2020} keep only the minimal constraint on GNN layers: equivariance to isomorphisms $\phi : \mathcal{V}(\Gamma) \to \mathcal{V}(\Gamma')$ between the nodes of graphs $\Gamma$ and $\Gamma'$ (called \emph{naturality}).
On a single graph, this relaxes the usual node-permutation equivariance from $S_n$ to its subgroup $\operatorname{Aut}(\Gamma)$.
Finally, \citet{lippmann2025iclr} focus on frames and handle $O(d)$-equivariance with \emph{local canonicalization}.
For a point cloud $x_i \in \mathbb E^d$, a global frame is given by an origin $o$ with an orthonormal basis $(e_1, \dots, e_d)$.
They additionally fix a local frame (i.e.\ a gauge) $R_i\in O(d)$ at each point $x_i$, which transforms as $R_i\mapsto R_i\hat R^{-1}$ under a global rotation or reflection $\hat R \in O(d)$.
This local canonicalization is partly learned: $R_i$ is obtained by orthonormalizing the vectors
\[
    v_{i,k} = \sum_{j \in \mathcal{N}(i)} \alpha_{ij,k} \, \frac{x_i - x_j}{\lVert x_i - x_j\rVert},
\]
where the invariant weights $\alpha_{ij,k} \in \mathbb R$ are computed by an MLP from invariant quantities (distances, scalar features).
Feature coordinates $f_i = \rho(R_i)F_i$ are then invariant to the active global $\hat R$, and can be processed by arbitrary MLPs.
Moreover, $R_i$ only depends on a small neighborhood of $x_i$, so similar local structures are oriented similarly and get similar coordinates $f_i$: the downstream network is presented with aligned inputs.

\subsubsection{Expressivity analysis}
% maron2018
% Dunford-Pettis theorem
% weiler2026: Part I seems to contain "necessity theorems" e.g. translation-equivariant linear layers are exactly convolutions
Equivariant architectures guarantee by construction that $\mathcal{H}\subseteq C_G(\mathcal{X},\mathcal{Y})$, the space of continuous $G$-equivariant maps.
Expressivity is about the reverse inclusion.
We report three answers from finest to coarsest.
\emph{Representation theorems} describe $C_G(\mathcal{X},\mathcal{Y})$, or a subclass of it such as linear maps, through an explicit form.
\emph{Universality} asks if $\overline{\mathcal{H}}=C_G(\mathcal{X},\mathcal{Y})$.
\emph{Separation power} cares about orbits, asking which inputs in distinct orbits, $x'\notin G\cdot x$, $\mathcal H$ can tell apart, i.e.\ if there is an $f \in\mathcal H$ with $f(x)\neq f(x')$.

\paragraph{Representation theorems and the particular role of the convolution}
Representation theorems give an explicit form that every equivariant map must take.
We give two examples.
First, for linear maps, they lead to ``convolution is all you need'' statements.
\citet{kondor2018p3icml} prove that, for a compact group $G$, linear equivariant maps between feature spaces over homogeneous spaces of $G$ are exactly generalized convolutions: convolution is not only sufficient for equivariance but also necessary.
\citet{weiler2018anips} prove the analogous statement for $SE(3)$ and feature fields over $\mathbb{R}^3$, and \citet{cohen2019anips} extend it to general homogeneous spaces from a fiber bundle perspective.
Second, on sets, \citet{zaheer2018} show that every permutation-invariant function on a countable domain decomposes as $f(X)=\psi\big(\sum_{x\in X}\phi(x)\big)$.

Such characterizations guide the design of $\mathcal{H}$: \textit{e.g.}\ approximating $\phi$ and $\psi$ by MLPs makes Deep Sets a universal approximator, a notion we now address.

\paragraph{Universality}
% pacini2026: something about \overline H not being too far from H
Universality is a strong requirement: $\overline{\mathcal{H}} = C_G(\mathcal{X},\mathcal{Y})$, typically for uniform convergence on compact sets.
For finite groups, \citet{ravanbakhsh2020p3icml} prove that equivariant MLPs with a single hidden layer on which $G$ acts regularly are universal.
For continuous groups, \citet{yarotsky2018r} covers invariant maps under compact groups and equivariant maps under translations and $SE(2)$, \citet{bogatskiy2020p3icml} tensor-product networks for $SO(3)$ and the Lorentz group, and \citet{vadgama2025} their $SE(3)$-equivariant network Rapidash.
\citet{elbaz2025} show that FS-KANs, permutation-equivariant Kolmogorov--Arnold networks, match the expressive power of parameter-sharing MLPs, so any universality statement transfers.

\paragraph{Separation power}
Universality for invariant functions implies that $\mathcal{H}$ separates every pair of orbits that continuous invariant functions separate, hence separation power is the weaker notion.
GNNs have been extensively studied from this angle through the \emph{WL hierarchy} of Weisfeiler--Lehman graph isomorphism tests, which measures the ability of architectures to distinguish non-isomorphic graphs.
\citet{xu2019,morris2019pacai} showed that message passing GNNs (MPNNs) are at most as powerful as $1$-WL, i.e.\ they cannot tell apart two graphs, or two nodes, that $1$-WL cannot.
Research then sought to surpass this barrier by adding subgraph counts as structural features \citep{bouritsas2021}, going to higher orders \citep{morris2019pacai,maron2020}, or lifting graphs to cell complexes \citep{bodnar2021a}.
\citet{maron2020} showed that $k$-th order invariant networks are at least as powerful as $k$-WL, and \citet{bodnar2021a} that cell complex networks are strictly more powerful than $1$-WL, a motivation for TDL (\Cref{sec:edl-permutation}).
However, the quest for expressivity need not be an arms race towards higher orders: \citet{hordan2024} showed that $2$-WL on complete weighted graphs is already \emph{complete} for position--velocity point clouds, i.e.\ separates all orbits under permutations, rotations and translations.

\subsubsection{Applications}
EDL principles have permeated a wide range of domains.
In atomistic modeling, \emph{neural network interatomic potentials} (NNIPs) serve as cheap surrogates of \textit{ab initio} methods such as density functional theory, predicting energies and per-atom forces from atomic configurations with $E(3)$-invariant \citep{schutt2017anips} or $E(3)$-equivariant \citep{batzner2022nc,liao2022,liao2023,passaro2023} architectures.
Equivariant GNNs also simulate $n$-body dynamics \citep{satorras2021p3icml} and generate 3D molecules via equivariant diffusion \citep{cornet2024}.
In structural biology, geometric inductive biases underpin protein structure prediction \citep{jumper2021na,baek2021s} and RNA structure scoring \citep{townshend2021s}.
Further scientific uses include Lorentz-equivariant jet tagging in particle physics \citep{bogatskiy2020p3icml} and 3D G-CNNs for pulmonary nodule detection in medical imaging \citep{winkels2019mia}.
In computer vision, spherical CNNs handle omnidirectional images and global climate patterns \citep{cohen2018,cohen2019p3icml}, while rotation and scale equivariance benefit 3D object detection \citep{yu2022} and object tracking \citep{sosnovik2021}, respectively.
Beyond vision, scale--translation equivariance has been applied to raw audio \citep{romero2023tmlr}, permutation equivariance to compositional generalization in language \citep{gordon2020}, state--action symmetries to reinforcement learning \citep{vanderpol2020anips}, and Euclidean equivariance to robotics and 3D scene understanding \citep{seo2025a,fei2024air}.



\subsubsection{Tradeoffs of equivariance}\label{sec:edl-tradeoffs}
% last par of [vadgama2025 sec 4]
Restricting the hypothesis class $\mathcal{H}$ to $G$-equivariant maps is a double-edged sword: it can improve sample efficiency and generalization, but it costs speed and design effort, can hinder optimization and can impose the wrong symmetry.
We review both sides and when equivariance is worth its cost.

\paragraph{Advantages of EDL}
Better sample complexity, i.e.\ learning from fewer examples, is put forward as a benefit of EDL.
\citet{tahmasebi2023} predict the gain in effective dataset size from the symmetry group (invariance scaling laws for kernel ridge regression).
Their follow-up \citep{tahmasebi2024p4icml} shows that a $G$-invariant distribution $\mu$ on a compact manifold $\mathcal{X}$ can be estimated from $n$ samples by some $\tilde{\mu}_n$ with, in $1$-Wasserstein distance,
\[
    \mathbb{E}\big[W_1(\mu,\tilde{\mu}_n)\big] \lesssim \Big(\frac{\operatorname{vol}(\mathcal{X}/G)}{n}\Big)^{1/d}, \qquad d = \dim(\mathcal{X}/G) \geq 3,
\]
versus $(\operatorname{vol}(\mathcal{X})/n)^{1/\dim\mathcal{X}}$ without invariance.
For a finite group acting faithfully, $\operatorname{vol}(\mathcal{X}/G) = \operatorname{vol}(\mathcal{X})/|G|$, i.e.\ each sample counts $|G|$ times, while groups of positive dimension also improve the exponent ($d < \dim\mathcal{X}$).
Concretely, these gains can translate into needing less data in applications, e.g.\ an order of magnitude less in medical imaging \citep{winkels2019mia}.

Another benefit is generalization.
In early work, \citet{shawe-taylor1991anips} derived a PAC bound in which the number of individual parameters is replaced by the number of weight-tying orbits.
\citet{elesedy2021p3icml} prove that, for linear models and an equivariant target, an equivariant model can generalize strictly better than its unconstrained counterpart.
\citet{sokolic2017p2icais} show that the generalization error of an invariant classifier is upper bounded in terms of the covering complexity of the base space $\mathcal{X}/G$, rather than that of $\mathcal{X}$.
By design, equivariant models are stable over $G$-orbits, which is not always the case for non-equivariant baselines, even with data augmentation \citep{lafarge2021mia}: this is a form of robustness.

\paragraph{Disadvantages of EDL}
The first drawback is speed.
Equivariant models can have a larger memory footprint, scale worse with the input size \citep{bogatskiy2020p3icml}, and have higher latency and lower training throughput (steps/sec) \citep{gerken2022p3icml,elhag2026}.
They also perform more FLOPs per parameter with worse GPU utilization \citep{brehmer2025}, and $E(3)$-equivariant models built on irreducible representations remain computationally heavy despite recent tensor-product speed-ups \citep{passaro2023,luo2024a}.
These costs, however, depend strongly on hardware and software: equivariant models cannot always tap into the most optimized libraries \citep{qu2024anips,shen2026}.

Second, equivariant models can be harder to optimize, i.e.\ have a lower \emph{optimization ease} \citep{shen2026}, the loss reduction per training step $\mathcal{L}(\theta_t)-\mathcal{L}(\theta_{t+1})$, with $\theta_t$ the parameters at step $t$.
Multiplied by throughput, it gives the loss reduction per second, whose deficit \citet{shen2026} call the \emph{efficiency gap}.
\citet{qu2024anips} hypothesize that symmetry constraints impede optimization as models scale, and \citet{pertigkiozoglou2024} conjecture that the equivariant solution space is too constrained, hence relax equivariance during training and project back onto that space at test time.
\citet[app.~A]{elhag2026} give qualitative support: GATr's loss landscape near its minimum is rougher than a standard Transformer's.

Third, misspecification: baking in the wrong prior can degrade performance \citep{yang2023p4icml,finzi2021anips,benton2020anips,miao2023p4icml}.
For instance, \citet{weiler2019anips} report, as expected, that equivariance to the dihedral group $D_n$ (rotations and reflections) performs worse than $C_n$ on rotated (unreflected) MNIST.
\citet{wang2023a} prove an error lower bound for incorrect equivariance, whereas \citet{wang2022b} report that equivariant models can still outperform unconstrained ones when the imposed symmetry is \emph{extrinsic} (\Cref{sec:symmetry-breaking-dl}).

\paragraph{The symmetry bottleneck}\label{par:sym-bottleneck}
Equivariant neural network methods demand \textit{a priori} knowledge of the data's symmetries before they can be hard coded to the architecture.
This creates a bottleneck for three reasons.
\begin{enumerate}
    \item \textbf{Unknown symmetries}: in real datasets symmetries are often unspecified or difficult to define analytically.
    \item \textbf{Manual sophistication}: understanding which symmetry group acts on data or designing a novel architecture requires domain expertise and mathematical sophistication.
    \item \textbf{Inflexibility}: once a symmetry is hard-coded into the architecture, the model cannot adapt if the assumptions are wrong.
    Many real-world datasets only approximately satisfy a given symmetry, for instance due to sensor noise, occlusions, or genuine symmetry breaking in the underlying phenomenon, and in that case strictly equivariant models may be over-constrained.
\end{enumerate}
These hurdles are addressed, respectively, by (1) symmetry discovery (\Cref{sec:sym-disc-from-data}), (2) general recipes addressing multiple groups at once \citep{ravanbakhsh2017p3icml,finzi2020p3icml,finzi2021p3icml,pearce-crump2023p4icml}, foundational approaches \citep{goel2025} or symmetrization of a backbone \citep{puny2021,kaba2023p4icml}, and (3) approximate soft-equivariance methods (\Cref{sec:soft-equivariance}).

\paragraph{To be or not to be}
A practical question is whether the design cost, potential optimization difficulties and hardware inefficiency of less mature implementations are justified compared with a simpler model, possibly trained with data augmentation.

\citet{nordenfors2024} compare equivariant models with augmented vanilla ones: under a compatibility condition, both share the same equivariant stationary points, but some are stable only for the equivariant model.
The relative benefit of equivariance also seems to depend on how well the data satisfy the symmetry.
\citet{wang2023a} show that, depending on levels of \emph{correct}, \emph{incorrect} and \emph{extrinsic equivariance} in the data (\Cref{sec:symmetry-breaking-dl}), either the equivariant or the unconstrained model has the advantage.
Similarly, \citet{lawrence2026} propose a decision rule based on the amount of \emph{distributional symmetry breaking} (\Cref{sec:distributional-sb}) in the dataset and its correlation with labels.
\citet{rygiel2026} reach a similar conclusion by comparing equivariant and vanilla models on datasets with varying degrees of canonical alignment.
Using Rapidash, an $SE(3)$-equivariant model whose equivariance can be selectively broken, \citet{vadgama2025} find that the advantage of well-aligned constraints is not closed by scaling capacity.
A notable, if isolated, example of equivariance receding is AlphaFold~3 \citep{abramson2024n}, which dropped the $SE(3)$-equivariant structure module from its earlier version \citep{jumper2021na} to simplify the architecture at no accuracy cost.

Beyond direct comparisons, \citet{brehmer2025} study scale: at matched compute ($10^{16}$--$10^{19}$ FLOPs), an $E(3)$-equivariant transformer reaches a roughly $2\times$ lower test loss than a non-equivariant one, with compatible power-law exponents.
Equivariance thus remains worthwhile at scale, provided its implementations match standard FLOP throughput.
Conversely, \citet{wang2024p4icmla} (conformer generation) and \citet{qu2024anips} (interatomic potentials) argue that scale compensates for the absence of built-in equivariance.
However, they compare models at similar parameter count or wall-clock time rather than at matched FLOPs, so the two views are less opposed than they appear.
\citet{shen2026} find that, for 3D-rotation equivariant tasks, augmented unconstrained models reduce their equivariance error to below 2\% of the total loss within thousands of steps.
The learnability of a symmetry should therefore also enter the cost-benefit analysis.

While there is no denying that equivariant models are scientifically well-founded, one may wonder whether this will be another occurrence of the bitter lesson~\citep{sutton2019}.

\subsection{Data augmentation}\label{sec:data-augmentation}
% add adversarial training ? [Goodfellow, Explaining and harnessing adversarial examples] and [Madry, Towards deep learning models resistant to adversarial attacks]
% shorten2019jbd and mumuni2022a: surveys on DA
% miao2023p4icml: instance-dependent DA
\begin{figure}[!tbp]
    \centering
    \includegraphics[width=\textwidth]{images/five-routes.pdf}
    \caption{Five routes to invariance on a single pipeline: each row colors the block where $G$ enters.
    Model-level routes are exact for every $\theta$, parameter-level ones learned and approximate.}
    \label{fig:five-routes}
\end{figure}

Data augmentation (DA) shares the goal of EDL, distilling a symmetry of the ground truth $f^*$, but from the \emph{problem} side: rather than constraining $\mathcal{H}$, it enriches the training set by applying transformations known (or assumed) to leave the target label unchanged.
Formally, letting $G$ act on the input space $\mathcal{X}$, the augmented risk replaces the standard empirical risk by an average over orbits:
$$
\hat R_{\mathrm{aug}}(\theta)=\frac{1}{|G|}\sum_{g\in G}\hat R(\theta;\,g\cdot \mathcal{D}),\qquad g\cdot\mathcal{D}=\{(g\cdot x_i,\,y_i)\}_{i},
$$
and for a continuous group, an expectation is taken over a distribution on $G$.
The hope is to encourage $f_\theta$ to be approximately invariant to those transformations.
Since the hypothesis class is unconstrained, the symmetry is encouraged for a specific parameter $\theta$, which corresponds to the parameter locus of \Cref{fig:sym-loci}.
\Cref{fig:five-routes} contrasts the entry point of the symmetry group $G$ in DA with equivariant layers and symmetrization methods (\Cref{sec:edl-achieving-equiv}), and soft penalties (\Cref{sec:approx-soft}).
DA is declarative: it states which transformations should not matter, not how $f_\theta$ realizes the invariance, accommodating transformations lacking an algebraic description (\textit{e.g.}\ illumination, viewpoint or background changes).
The symmetry often stays \emph{implicit}, and even \emph{unknown a priori} when automated strategies (\Cref{par:auto-da}) learn the transformations.
DA has further practical advantages over EDL.
It applies to any pretrained model without architectural changes, which suits fine-tuning.
It requires no group-theoretic implementation, which eases rapid prototyping.
It also seamlessly handles extent-restricted symmetries (\Cref{sec:extent-restriction}), \textit{e.g.}\ non-invertible crops or small-angle rotations on MNIST (avoiding $6\leftrightarrow 9$), where EDL mostly assumes exactness.
On the downside, a continuous group cannot be covered exhaustively, and materialized augmentations multiply storage and compute by the number of samples drawn per orbit.

\subsubsection{Theoretical approaches}
\citet{chen2020} propose a group-theoretic framework for DA and formally connect orbit averaging with lower generalization error.
Variance reduction results can also be obtained \citep{chen2020,dao2019p3icml}.
Orbit averaging is adapted and practically implemented by \citet{benton2020anips}.
For infinite-width and kernel-classifier settings, DA is shown to be equivalent to symmetrization of the model \citep{maass2024,dao2019p3icml}.
\citet{lyle2020} extend PAC-Bayes bounds to DA and obtain variance reduction for the empirical risk.
The value of DA may be context dependent, the relevant aspects being the learning regime (over- or underparametrized) and the task (classification or regression).
A spectral point of view is given by \citet{lin2024}, while \citet{lawrence2026} further study the role of the symmetry distribution.

\subsubsection{Automated DA strategies}\label{par:auto-da}
% automated DA vs sym discovery: automated DA has the goal of producing new augmentations, it can be transductive whereas sym discovery is more inductive and aims at recognizing the group and its actions (which can be an intermediate step before generating augmentations anyways)
Rather than hand-crafting augmentation policies, several works have proposed to search for them automatically.
\citet{cubuk20192iccvprc} introduce AutoAugment, which uses reinforcement learning to find a policy composed of image-processing operations such as rotations, translations, and shearing, optimized for validation accuracy.
\citet{lim2019anips} reduced the search time by leveraging Bayesian optimization over policies and using \textit{density matching} as a criterion, i.e.\ minimizing the divergence between the augmented and training distributions.
\citet{hataya2020cv-e2} proposed a differentiable approach, training a generative adversarial network (GAN) to output augmentation parameters that minimize a similar density-matching objective, which further accelerates the search.
Another strategy is to transfer augmentations across domains: \citet{antoniou2018} learn a GAN on a data-rich source domain and replay the learned transformations on a data-scarce target domain.
For this symmetry transfer to be effective, the transformations learned on the source domain must also be symmetries of the target domain, i.e.\ the transferred group must be a subgroup of the target's symmetry group. 
Symmetry discovery can be operationalized in a straightforward two-step pipeline, for instance by inferring a Lie algebra basis and using augmentation in a (soft) inductive bias \citep{escriche2025,yang2023p4icml} (see \Cref{sec:sym-disc-from-data} for more on symmetry discovery).
A practical advantage shared by all these approaches is interpretability: learned augmentations produce concrete transformed examples that can be inspected directly.
One limitation, however, is that materializing augmentations requires an explicit parametric form for the transformations.
\citet{shoshan2021} partly circumvent this in the face-generation setting by showing that a scalar attribute predictor (\textit{e.g.}\ age) provides sufficient supervision for a GAN to learn to manipulate that attribute in a controllable way without a closed-form group action.

\subsubsection{DA in relation to Feature Averaging}\label{sec:da-fa}
Feature Averaging (FA) is another strategy for incorporating symmetries: it averages model outputs over a set of transformed inputs, effectively symmetrizing the model.
Theoretically, \citet{dao2019p3icml} show that DA can be approximated to first order by FA, while \citet{maass2024} establish, in the infinite-width limit, that DA and FA induce identical gradient flow dynamics when the activation functions respect the group action.
Despite this asymptotic equivalence, \citet{lyle2020} prove that, for convex losses, FA has a risk no larger than DA (by Jensen's inequality) and tighter PAC-Bayes generalization bounds, suggesting FA is preferable when applicable.
Similarly to DA, FA can be applied beyond strict group-invariance settings, accommodating approximate or partial symmetries.
Another natural comparison is with equivariant models, covered in \Cref{sec:edl-tradeoffs}.

\subsection{Framework Symmetries}\label{sec:framework-symmetries}
Unlike EDL and DA, which distill a problem symmetry into $\mathcal{H}$ or $f_\theta$, \emph{framework} symmetries are carried by the training pipeline, across several loci or forward calls (\Cref{fig:sym-loci}).
Symmetries arise in the way data is piped through during training (\Cref{sec:data-ordering,sec:joint-embedding}) or reflect a limitation in downstream processing rather than a true problem constraint: criterion symmetries (\Cref{sec:criterion-symmetries}) are the irreducible loss symmetries that cannot be addressed through parameter or architectural adjustments.
From a taxonomic viewpoint, data ordering (\Cref{sec:data-ordering}) is inherently \emph{dynamic}, as its symmetry permutes the data stream over training, while joint embeddings (\Cref{sec:joint-embedding}) and criterion symmetries (\Cref{sec:criterion-symmetries}) are both \emph{prescriptive}, designing how to train.
These symmetries can be incomplete (\Cref{sec:incomplete-symmetries}), \textit{e.g.}\ full-batch GD is $S_n$-invariant but mini-batch SGD only under a subgroup or only approximately. The works are often \emph{implicit} and the overall symmetry reading of joint embedding (\Cref{sec:joint-embedding}) is ours.
Compared to the symmetries of $\mathcal{L}$ on $\Theta$ (\Cref{sec:parameter-space-symmetries}), criterion symmetries act on $\mathcal{Y}$, and the corresponding works lean towards the implicit and the prescriptive.

\subsubsection{Data Ordering}\label{sec:data-ordering}
We start with the symmetries baked into the training procedure by the data stream.
Formally, view a learning algorithm as a map
$$
\mathcal{A} : \left ( \{\mathcal{D}_e\}_{e\in [E]},\, \theta_0 \right ) \longmapsto \theta_T,
$$
sending per-epoch datasets $\mathcal{D}_e$ of $n_e$ points, with $E$ the number of epochs, and an initialization $\theta_0$ to trained parameters $\theta_T$ after $T$ training steps.
Here we do not consider other sources of randomness such as augmentation draws or dropout masks.
Materialized as a sequence, each $\mathcal{D}_e$ is acted on by $S_{n_e}$ through sample permutation, and one may ask whether $\mathcal{A}$ is invariant: $\mathcal{A}(\{\sigma_e\cdot\mathcal{D}_e\}_{e\in [E]},\theta_0)=\mathcal{A}(\{\mathcal{D}_e\}_{e\in [E]},\theta_0)$ for all $\sigma_e\in S_{n_e}$.
Over one epoch of $n=n_e$ samples, full-batch gradient descent (GD) sums the gradients of all $n$ samples before updating, and is therefore $S_n$-invariant.
Mini-batch GD with a fixed partition into $k$ batches of size $B=n/k$ (assumed integer for simplicity) sums gradients within each batch but applies the $k$ updates sequentially, so it is invariant only under the subgroup $\prod_{i=1}^{k} S_B \leq S_n$ of within-batch permutations, and the extreme case ($B=1$) retains no invariance.
The size ratio at fixed $k$ between these symmetry groups is
$$
\frac{n!}{(B!)^{k}} \;\underset{n\to\infty}{\sim}\; (2\pi n)^{\frac{1-k}{2}}\, k^{\,n+\frac{k}{2}}
$$
by Stirling's formula.
As $n$ grows, GD is invariant to an ever larger set of data orderings than its mini-batch counterparts, and over multiple epochs the gap widens further when the data are reshuffled at each epoch (see below).

\citet{zhuang2022pmls} distinguish \emph{algorithmic} sources of stochasticity (data augmentation, data shuffling) from \emph{implementation} ones (floating-point effects) and quantify their impact on training through three metrics: (i) the \emph{churn} $C(f_1,f_2)=\mathbb{E}_x\big[1_{\operatorname{pred}(f_1,x)\neq \operatorname{pred}(f_2,x)}\big]$, the fraction of test points on which two independently trained models disagree, (ii) the distance $\|\theta_1-\theta_2\|_2$ between trained parameters, an imperfect metric because of the parameter-space symmetries of \Cref{sec:parameter-space-symmetries}, and (iii) the standard deviation of accuracy on subgroups of the data, quantifying local variability in input space.
While global test accuracy is constant across runs, these three metrics fluctuate substantially: churn ranges from $5\%$ to over $20\%$.
For data ordering specifically, SGD exhibits higher variability than GD, consistent with its smaller invariance group.

\paragraph{Shuffling schemes and interaction with Batch Normalization}
In practice, there are multiple ways to implement mini-batch training.
\emph{Stochastic sampling} draws each batch with replacement and thus involves no permutation.
\emph{Single shuffle} (SS) samples one permutation $\sigma\in S_n$ once and reuses it at every epoch.
\emph{Random reshuffling} (RR) resamples a new permutation at each epoch.
Randomized ordering is not optional: mini-batch SGD requires it to converge well \citep{xu2024vj}, and, perhaps surprisingly, the permutation-based schemes SS and RR outperform with-replacement sampling \citep{koloskova2024p4icml}.
The distinction between SS and RR is of practical importance because globally reshuffling a dataset can be a hardware bottleneck, which at scale makes RR impractical \citep{nguyen20222iipdpsi,koloskova2024p4icml,xu2024vj}.
Batch normalization (BN) introduces an additional dependency on data ordering since batch composition determines the batch statistics.
\citet{wu2023p4icml} show that standard GD, SS and RR then optimize \emph{three different objectives}.
They study linear models with BN applied to the input, which take the form $x \mapsto W\Gamma\operatorname{BN}(x)$ with $\Gamma$ a diagonal matrix carrying BN's learnable scaling parameters.
Writing $k$ for the number of batches and $(\sigma X)_i$ for the $i$-th batch under the ordering $\sigma$, the objectives read
$$
\mathcal{L}_{\mathrm{GD}} = L\big(W\Gamma\operatorname{BN}(X),\, Y\big), \qquad
\mathcal{L}_{\sigma} = \sum_{i=1}^{k} L\big(W\Gamma\operatorname{BN}((\sigma X)_i),\, (\sigma Y)_i\big), \qquad
\mathcal{L}_{\mathrm{RR}} = \frac{1}{n!}\sum_{\sigma\in S_n} \mathcal{L}_{\sigma},
$$
where SS minimizes $\mathcal{L}_\sigma$ for a single fixed random $\sigma$.
Since $\sigma$ affects the per-batch normalization statistics, $\mathcal{L}_\sigma \neq \mathcal{L}_{\mathrm{GD}}$ in general.
\citet{wu2023p4icml} further show that SS can suffer from training instabilities that RR avoids and, consistently, \citet{zhuang2022pmls} observe empirically that with BN, RR stabilizes training against data-ordering variability.

\begin{Remark}[Non-invariance due to implementation]
    Floating-point addition is not associative, so even the gradient sum of full-batch GD is not permutation-invariant at the numerical level: \citet{zhuang2022pmls} measure non-zero churn for GD under data reordering alone.
\end{Remark}

\paragraph{Exchangeability and invariance in law}\label{par:exchangeability}
The invariance lost by SGD can be restored in distribution, by the data or by the algorithm.
A sequence $\boldsymbol{X}_n=(X_1,\dots,X_n)$ is \emph{exchangeable} (a generalization of the i.i.d.\ assumption) if $\sigma\boldsymbol{X}_n \overset{d}{=} \boldsymbol{X}_n$ for every $\sigma\in S_n$.
Restricting to a single epoch $\mathcal{D}$ for simplicity, if $\mathcal{D}$ is exchangeable, then $(\sigma\cdot\mathcal{D},\theta_0)\overset{d}{=}(\mathcal{D},\theta_0)$ and
$$
\operatorname{Law}\big(\mathcal{A}(\sigma\cdot\mathcal{D},\,\theta_0)\big) = \operatorname{Law}\big(\mathcal{A}(\mathcal{D},\,\theta_0)\big) \qquad \text{for all } \sigma\in S_n,
$$
even though each realization of $\mathcal{A}$ is not $S_n$-invariant.
\citet{bloem-reddy2020} characterize a stronger, conditional invariance through the \emph{empirical measure} $\mathbb{M}_{\boldsymbol{X}_n}(\cdot):=\sum_{i=1}^{n}\delta_{X_i}(\cdot)$, which represents the data but forgets its ordering: for exchangeable $\boldsymbol{X}_n$, $(Y \mid \boldsymbol{X}_n)\overset{d}{=}(Y \mid \sigma\boldsymbol{X}_n)$ for all $\sigma\in S_n$ if and only if $(\boldsymbol{X}_n, Y)\overset{a.s.}{=}\big(\boldsymbol{X}_n,\, f(\eta, \mathbb{M}_{\boldsymbol{X}_n})\big)$ for some measurable $f$ and noise $\eta$ independent of $\boldsymbol{X}_n$.
Shuffling (RR, SS) yields this factorization even if $\mathcal{D}$ is not exchangeable.
If $\sigma'\in S_n$ is drawn uniformly and independently,
$$
\theta_T=\mathcal{A}(\sigma'\cdot\mathcal{D},\theta_0)=f(\eta_\sigma,\mathbb{M}_{\mathcal{D}})
$$
and the trained parameters depend on the data only through their empirical measure, up to independent noise.

\begin{Remark}[Examples of non-exchangeable data]
    Not all scenarios can be modeled with exchangeable data.
    In \emph{curriculum learning} \citep{bengio2009p2aicml}, examples are deliberately ordered, often from easy to hard, or progressing from primitive to compositional tasks \citep{lee2024}.
    This can be framed either as a single sorted dataset ($E=1$) or as a progression over per-epoch datasets ($\mathcal{D}_e\neq \mathcal{D}_{e'}$ for $e\neq e' \in [E]$).
    Ordering biases can also be more fortuitous: data stored in database systems can be clustered by attribute, producing a non-exchangeable stream \citep{xu2024vj}.
\end{Remark}

\subsubsection{Joint-embedding architectures}\label{sec:joint-embedding}

Joint-embedding (or Siamese) architectures learn representations by training a network to closely embed two related views rather than by predicting a label or reconstructing the input.
We give here a symmetry-centered reading of the development of this family of self-supervised methods, and refer to \citet{ilina2022sa} for a more in-depth account.
In this part, invariance prevails, but equivariance-based methods have also been explored~\citep{gupta2023a}.

In the simplest setting, for views $x_1$ and $x_2$ of the same underlying object $x$ and encoders $f_{\theta_1}$ and $f_{\theta_2}$, a loss 
\begin{equation}\label{eq:joint-embeddings}
    L\big(f_{\theta_1}(x_1),\,f_{\theta_2}(x_2)\big)
\end{equation}
computing a similarity between embeddings is minimized.
A common way to generate views is through data augmentation: for $x\sim\mathcal{D}$, two transformations $t_1,t_2$ are drawn i.i.d.\ from an augmentation family $\mathcal{T}$ to form the views $x_1=t_1\cdot x$ and $x_2=t_2\cdot x$.
This setup carries a double (often inexact) symmetry:
\begin{enumerate}
    \item A target invariance $f_{\theta_i}(t_1x)=f_{\theta_i}(t_2x)$ of the implemented encoder functions (parameter-level in \Cref{fig:sym-loci}).
    Involving data transformations, it is conceptually close to DA (\Cref{sec:data-augmentation}) and soft penalties (\Cref{sec:soft-equivariance}).
    In contrast with DA, which is typically a supervised approach focusing on the invariance of the $(x,y)$ mapping, this symmetry focuses on the invariance of the representation.
    This symmetry is the end goal: ultimately we want $f_\theta$ to be invariant to $\mathcal{T}$.
    \item A \emph{branch symmetry} in the training protocol: the objective is invariant under swapping the two branches $(\theta_1,t_1)\leftrightarrow(\theta_2,t_2)$.
    With shared weights $\theta_1=\theta_2$ and a symmetric $L$, this is an exact $C_2$ symmetry, but it is often deliberately broken, and much of the design space below is organized by how and where it is broken (\Cref{fig:joint-embedding}).
    It is transparent at inference time and is only a means to train and achieve the above target invariance.
\end{enumerate}

\noindent \emph{Representational collapse} is the failure mode where the encoder is constant, while the subtler \emph{dimensional collapse} occurs when the embedding space becomes rank-deficient.
Methods can be understood by how far they depart from exact symmetries toward carefully engineered asymmetry designed to learn and avoid collapse.

\begin{figure}[!tbp]
    \centering
    \includegraphics[width=\textwidth]{images/joint-embedding.pdf}
    \caption{Joint-embedding architectures. (\textbf{a}) Fully symmetric: $z_1=z_2$, so $L$ is uninformative. (\textbf{b}) Design space: carefully chosen asymmetries departing from (\textbf{a}).}
    \label{fig:joint-embedding}
\end{figure}

\paragraph{Avoiding representational collapse}
\emph{Contrastive methods} keep branch-swap symmetry (shared weights) exactly and instead prevent collapse by adding negative pairs to the objective, balancing agreement between augmented views of the same $x$ against disagreement between different inputs.
Early metric learning methods proposed a symmetric loss $\ell(x_1,x_2,\theta)$ trained on matched/mismatched pairs \citep{chopra20052icsccvprc,hadsell20062icsccvpr-v2c}.
\citet{schroff20152iccvprc} extend it to a triplet loss $\ell(x,x_p,x_n,\theta)$ contrasting one positive against one negative, and \citet{oord2018a} generalize further to an arbitrary number of negatives with the InfoNCE loss.

\citet{chen2020p3icml} (SimCLR) apply strong compositional data augmentation and a projection head $\operatorname{proj}_\theta$, still shared between branches, on top of which the loss is computed:
\begin{equation}
    \ell_{i,j}=-\log\frac{\exp\big(\operatorname{CoSim}(z_i,z_j)/\tau\big)}{\sum_{k=1}^{2N} 1_{[k\neq i]}\exp\big(\operatorname{CoSim}(z_i,z_k)/\tau\big)},\qquad z_k=\operatorname{proj}_\theta(f_\theta(\tilde x_k)),
\end{equation}
where $\tilde x_1,\dots,\tilde x_{2N}$ are the paired views of each of the $N$ inputs of a batch, $(i,j)$ indexes the two views of the same input, $\operatorname{CoSim}$ is the cosine similarity and $\tau$ a temperature. The full loss averages $\ell_{i,j}$ over all positive pairs.
An ablation shows that the representation $f_\theta(x)$ before the projector is better than $\operatorname{proj}_\theta(f_\theta(x))$ after it: $\operatorname{proj}_\theta$ is trained to be augmentation-invariant, which it achieves by discarding useful representational information.
The authors confirm this empirically by measuring how much augmentation-related information (\textit{e.g.}\ color) is decodable at each stage.
\citet{he2020} (MoCo) relax branch-swap symmetry itself: the second branch (the ``momentum encoder'') tracks the first through an Exponential Moving Average (EMA) of its weights, $\xi\leftarrow m\,\xi+(1-m)\,\theta$, rather than sharing its weights outright. The non-contrastive methods below push this decoupling of the branches further.

\emph{Non-contrastive methods} dispense with negative pairs altogether and instead prevent collapse either by decoupling the branches more strongly or by regularizing the embedding statistics directly.
Among the former, \citet{grill2020} (BYOL) equip only the ``student'' branch with an extra prediction head $q_\theta$, computing $q_\theta(\operatorname{proj}_\theta(f_\theta(t\cdot x)))$ against the ``teacher'' target $\operatorname{proj}_\xi(f_\xi(t'\cdot x))$, where $\xi$ is an EMA of $\theta$: the student must \emph{predict} rather than match the teacher, turning the symmetric matching objective into an asymmetric one.
\citet{caron2021} (DINO) instead keeps a single set of weights (EMA aside) but breaks symmetry in the output distribution with centering and temperature.
The student computes
\[
    \operatorname{Softmax}_{\tau_s}(\operatorname{proj}_{\theta}(f_\theta(t\cdot x))),
\]
while the teacher computes
\[
    \operatorname{Softmax}_{\tau_t}
    \bigl(\operatorname{Center}(\operatorname{proj}_{\xi}(f_\xi(t'\cdot x)))\bigr),
\]
with $\xi$ an EMA of $\theta$ and $\tau_t<\tau_s$: the lower temperature artificially sharpens the teacher's target.
\citet{chen2020a} (SimSiam) pulls back toward full symmetry: weights are entirely shared and the two branches alternate between student and teacher roles within a single symmetrized loss.
With $z=\operatorname{proj}_\theta(f_\theta(t\cdot x))$ and $z'=\operatorname{proj}_\theta(f_\theta(t'\cdot x))$,
\[
\mathcal{L}=-\operatorname{CoSim}\big(q_\theta(z),\,\operatorname{sg}(z')\big)-\operatorname{CoSim}\big(q_\theta(z'),\,\operatorname{sg}(z)\big),
\]
where $q_\theta$ is a prediction head as in BYOL and $\operatorname{sg}$ the stop-gradient operator.
Branch-swap symmetry is fully retained, and collapse is instead prevented by the loss circuitry and the stop-gradient, an asymmetry introduced in the \emph{optimization} rather than in the architecture.

A second strand of non-contrastive methods regularizes the batch statistics of the embeddings, keeping branch-swap symmetry exact and using no stop-gradient.
\citet{zbontar2021p3icml} (Barlow Twins) drive the cross-correlation matrix $\mathcal{C}$ between the two branches' embeddings toward the identity with a loss
\[
\mathcal{L}=\big\|[\,I-\mathcal{C}\odot I\,]\big\|_F^2+\lambda\big\|\mathcal{C}\odot(\boldsymbol 1\boldsymbol 1^\top-I)\big\|_F^2\quad ,
\]
where $\boldsymbol{1}$ is the all-ones vector, $\odot$ is the Hadamard product, and $\lambda$ a hyperparameter.
The first term enforces invariance (at a minimum, branches agree on each embedding coordinate) and the second term decorrelates coordinates across branches, preventing dimensional collapse.
\citet{bardes2022} (VICReg) optimize similar objectives but reintroduce a direct invariance penalty $\sum_i \lVert z_i-z_i'\rVert_2^2$ computed on embeddings $z_i$ and $z_i'$. 

\paragraph{Invariance across modalities}
\citet{radford2021p3icml} (CLIP) apply the same branch-swap logic across modalities: image and text encoders embed into a shared space, cosine similarities $s_{ij}=\operatorname{CoSim}(z_i^{\mathrm{txt}},z_j^{\mathrm{img}})$ populate a matrix $S$, and a symmetric cross-entropy
\[
-\left[\sum_i\log\frac{\exp(s_{ii})}{\sum_j\exp(s_{ij})}+\sum_j\log\frac{\exp(s_{jj})}{\sum_i\exp(s_{ij})}\right]
\]
treats both modalities identically, echoing SimSiam's~\citep{chen2020a} symmetric treatment.
Here, however, the two encoders cannot share weights, and the symmetry is only present at the block level and in the loss.
CLIP also uses comparatively little data augmentation, presumably because the image/text modality gap already supplies ample view diversity.
Multimodal joint-embeddings highlight the question of a representational hierarchy: \citet{desai2023p4icml} replace the Euclidean embedding space with a hyperbolic one paired with an entailment loss, so that the difference in level of detail between an image and the text describing it is represented geometrically.

\subsubsection{Criterion Symmetries}\label{sec:criterion-symmetries}
% wang2025: quotient loss for SE(3) + orbit loss for scale = Sim(3) overall
The differentiable function $L:\mathcal{Y}\times\mathcal{Y}\to\mathbb{R}$ used to compare predictions $\hat y$ and ground truth $y$, called here the \emph{criterion}, is distinct from the loss $\mathcal{L}(\theta)=\frac{1}{n}\sum_{i=1}^{n} L(f_\theta(x_i),y_i)$ on $\Theta$: its symmetries act on predictions and targets in $\mathcal{Y}$, not on parameters.
The two are related when a transformation $g$ of $\mathcal{Y}$ is realizable in parameter space, $f_{\tilde g\cdot\theta}=g\cdot f_\theta$ for some $\tilde g$: a criterion symmetry on predictions then induces a type of loss symmetry $\mathcal{L}(\tilde g\cdot\theta)=\mathcal{L}(\theta)$ with $f_{\tilde g\cdot\theta}\neq f_\theta$.
We say $L$ is \emph{invariant} under a group $G$ acting on $\mathcal{Y}$ when it depends only on the target or prediction orbit,
$$
L(\hat y,\, y) = L(\hat y,\, g\cdot y) \quad \text{ or } \quad L(\hat y,\, y) = L(g\cdot \hat y,\, y)
$$
for all $g\in G$.
\citet{kaba2025anips} additionally require symmetry between the two arguments and call \emph{invariant loss function} a loss criterion satisfying
\begin{equation}\label{eq:invariant-loss}
    L(g\cdot\hat y,\, y) = L(\hat y,\, g\cdot y) = L(\hat y,\, y) \qquad \forall g\in G.
\end{equation}

\begin{Remark}[Diagonal invariance]
    \Cref{eq:invariant-loss} is invariance under $G\times G$ acting on $(\hat y,y)$ componentwise. The weaker invariance under its diagonal subgroup $\Delta G=\{(g,g)\}\le G\times G$, i.e.\ $L(g\cdot\hat y,\; g\cdot y)=L(\hat y,\; y)$ for $g\in G$ (e.g.\ MSE for orthogonal $\rho_Y$), also appears in the literature~\citep{nordenfors2024,lawrence2026,shen2026}, for instance as a natural assumption on $L$ when training an equivariant model on an equivariant task: requiring the transformed pair $(gx,gy)$ to incur the same loss as $(x,y)$, $L(f(gx),gy)=L(f(x),y)$, is exactly diagonal invariance since $f(gx)=gf(x)$.
\end{Remark}

A typical example is when $y$ is a vector whose ranges of indices represent the elements of a set.
Concretely, assume $y=(u,v)\in \mathbb R^{2n}$ is composed of two objects identical up to properties we want to predict: $u\in\mathbb{R}^n$ represents the properties of the first object (such as a spatial position for $n=3$) and $v\in\mathbb{R}^n$ those of the second.
A desideratum for $L$ is then invariance under the action of $S_2$ permuting the two objects: $g\cdot(u,v)=(v,u)$.
Subtler cases exist, as shown in the next examples.

\begin{Example}[Dependent group action]
    The MSE criterion $L(\hat y,y)=\|\hat y-y\|_2^2$ is \emph{not} invariant under a translation $y\mapsto y+t$, since $\|\hat y-(y+t)\|_2^2\neq\|\hat y-y\|_2^2$ in general.
    However, it carries a \emph{prediction-dependent} orthogonal symmetry: for $R\in O(d)$ and the action $g\cdot y:=R(y-\hat y)+\hat y$,
    $$
    \|g\cdot y-\hat y\|_2^2 = \|R(y-\hat y)\|_2^2 = \|y-\hat y\|_2^2,
    $$
    so MSE is an invariant loss for this particular ($\hat y$-dependent) action.
    Analogously, a ground-truth-dependent action can be defined on the target.
\end{Example}

\begin{Example}[False-class symmetry]\label{ex:false-class-symmetries}
    Because the criterion is blind to part of the prediction, a symmetry arises in cross-entropy $\operatorname{CE}(\hat y, y)=-\sum_i y_i\log\hat y_i$.
    When labels are pure, i.e.\ $y=e_c$ for the true class $c$, one has $\operatorname{CE}(\hat y,e_c)=-\log \hat y_c$, which depends on $\hat y$ only through its $c$-th coordinate.
    Hence
    \begin{equation*}
        \operatorname{CE}(P\hat y,\, y)=\operatorname{CE}(\hat y,\, y),
    \end{equation*}
    for any permutation matrix $P=P_\sigma$ with $\sigma(c)=c$, i.e.\ $\sigma$ in the stabilizer of the true class (\Cref{def:orbit-stabilizer}).
    Note that here $\hat y$ and $y$ do not play the same role.
\end{Example}

When the loss is not invariant but the task is, several targets $\{g\cdot y:g\in G\}$ are equally valid, and training against a fixed representative can be harmful.
The following works address this problem.

\paragraph{Orbit and quotient losses}
A sensible approach is to treat the loss as a function on the quotient $\mathcal{Y}/G$ using an \emph{orbit loss} to minimize over the group:
$$
L_{\mathrm{orbit}}(\hat y,\, y) = \inf_{g\in G}\, L(\hat y,\, g\cdot y).
$$
Such a loss is invariant by construction but raises an optimization problem that can be non-trivial (see \Cref{sec:framework-criterion-alignment}).
\citet{mehr20182iccvpr} instantiate this for an autoencoder on point clouds in $\mathbb{R}^3$ with $G=SO(3)$.
Given an input $x$, they sample its orbit $(g_1 x,\dots,g_M x)$, encode each element, and max-pool across the orbit.
This \emph{orbit-pooling} layer produces an (approximately) invariant latent vector $z(x)\approx z(g\cdot x)$, which is then decoded back to input space.
The optimization is conducted with a \emph{quotient loss} $$\mathcal{L}(\theta, \eta)=\inf_{g\in G} d\big(g\cdot x,\,\operatorname{dec}_\eta(\operatorname{enc}_\theta(x))\big),$$ whose infimum is found either by discretizing $SO(3)$ or numerically.

\paragraph{Permutation losses for unordered outputs}
A pervasive case is $G=S_n$ acting on an unordered collection of outputs, the \emph{label-ambiguity} or \emph{permutation problem}.
In source separation, the ``cocktail-party'' task of recovering $K$ speakers $x_k(t)$ from their mixture $\sum_{k=1}^K x_k(t)$, the speakers are unordered.
\citet{yu20172iicasspi} fold the assignment into the objective through \emph{permutation-invariant training} (PIT),
$$
L_{\mathrm{PIT}} = \min_{\sigma\in S_K} L_\sigma, \qquad L_\sigma \propto \sum_{k=1}^K \big\|\hat X_k - X_{\sigma(k)}\big\|_2^2,
$$
over predicted and ground-truth short-time Fourier transforms $\hat X_k, X_{\sigma(k)}$.
The same symmetry drives deep clustering for self-supervised learning, where arbitrary cluster indices are defined only up to a permutation and serve as pseudo-labels to learn representations.
\citet{caron2018,caron2019} recompute the clustering offline every few epochs, which tolerates a relabeling between rounds but incurs stale gradients, whereas \citet{caron2020anips} move to online, batch-level pseudo-labels.
In these settings, gradients stay fresh but labels are possibly permuted at every step, rendering optimization impossible with the previous approaches \citep{caron2018,caron2019}.
The issue is bypassed without leveraging group theory by maintaining a jointly-learned set of prototypes $(c_1,\dots,c_K)$ that act as persistent centroids.
Only the initial ordering of prototypes remains arbitrary, which does not prevent optimization.

In object-centric tasks, the output is a set of unordered objects and permutation invariance can only be enforced in the downstream loss, typically by optimal assignment.
\citet{locatello2020anips} (\emph{Slot Attention}) map CNN features $X\in\mathbb{R}^{N\times D_{\mathrm{in}}}$ to $K$ object slots $\{s_1,\dots,s_K\}$ through a recurrent module $S^{(t+1)}=\operatorname{SlotAttention}(X,S^{(t)})$ that is permutation-equivariant in the slots and permutation-invariant in the features,
$$
\operatorname{SlotAttention}(\sigma_i X,\, \sigma_s S) = \sigma_s \operatorname{SlotAttention}(X,\, S),
$$
the slot symmetry stemming from the set nature of the output, not from the input.
Because the output order is genuinely random, so that a consistent labeling in the manner of \citet{caron2020anips} is impossible here, targets and predictions are matched by the Hungarian algorithm before the loss is evaluated.
\citet{carion2020cv-e2} (DETR) and \citet{zhu2021} (Deformable DETR) apply the same device to bounding-box object detection.
For ground-truth objects $y=(y_1,\dots,y_M)$ with $y_i=(c_i,b_i)$ (class and bounding box), the model emits $N>M$ predictions $\hat y_j=(\hat p_j,\hat b_j)$, where $\hat p_j$ is a distribution over classes, and a \emph{Hungarian loss} matches them by minimizing
$$
L_{\mathrm{match}}(y_i,\hat y_j) = -\hat p_j(c_i) + L_{\mathrm{box}}(b_i,\hat b_j),
$$
trading classification confidence against a bounding-box overlap term $L_{\mathrm{box}}$ (e.g.\ IoU).

\paragraph{Set-distance losses for point clouds}
Principled comparison of point clouds can be done via a permutation-invariant set distance.
For 3D reconstruction, \citet{fan2017} predict a point cloud $\hat S_i = f_\theta(I_i, r)$ from an image $I_i$ and noise $r$, and train with a set distance $d$ summed over examples, $\mathcal{L}(\theta)=\sum_i d(\hat S_i, S_i^{\mathrm{gt}})$, proposing two choices.
The first is the \emph{Chamfer distance}
$$
d_{\mathrm{CD}}(S_1,S_2) = \sum_{x\in S_1}\min_{y\in S_2}\|x-y\|_2^2 + \sum_{y\in S_2}\min_{x\in S_1}\|x-y\|_2^2,
$$
invariant due to the nearest-neighbor correspondence.
The second, the \emph{Earth Mover's distance}, is a matching loss minimizing over bijections $\sigma$ (which requires $|S_1|=|S_2|$, although relaxations exist):
$$
d_{\mathrm{EMD}}(S_1,S_2) = \min_{\sigma:S_1\to S_2}\sum_{x\in S_1}\|x-\sigma(x)\|_2 .
$$
\citet{lin2023} likewise adopt the Chamfer distance for \emph{point-cloud completion}, the task of inferring a full cloud from a partial one (as arises under occlusion or limited sensor resolution).

\paragraph{Alignment}\label{sec:framework-criterion-alignment}
In the above examples, minimization problems are only a means to an end: the loss is evaluated on the best-aligned pair of outputs and targets.
Recovering the minimizing group element can be a challenge in itself.
\citet{alvarez-melis2018p2cemnlp} align point clouds living in \emph{distinct} metric spaces (e.g.\ word embeddings of two languages) via the Gromov--Wasserstein distance: from within-domain similarity matrices $C_{ij}=d(i,j)$ and $C'_{kl}=d'(k,l)$, they solve
$$
\min_{\Gamma}\ \sum_{ijkl} L\big(C_{ij}, C'_{kl}\big)\,\Gamma_{ik}\Gamma_{jl}
$$
over transport plans $\Gamma$, transporting and comparing \emph{distances} rather than points.
The two clouds may differ in size, making the setup more general.
\citet{even2024} study a \emph{planted} version: given a point cloud $X=(x_1,\dots,x_n)$ and a noisy transformed copy $y_i=Q^* x_{\sigma^*(i)}+\nu z_i$ with $Q^*\in O(d)$ and $z_i\sim\mathcal{N}(0,I_d)$, the goal is to recover the pair $(Q^*,\sigma^*)$.
Their \emph{Ping-Pong} algorithm alternates between solving for $Q$ with $\sigma$ fixed and for $\sigma$ with $Q$ fixed, two tractable subproblems.

\begin{Remark}[A transparent permutation symmetry]
Not every permutation symmetry in a loss calls for special treatment.
Passing from a triplet loss to the InfoNCE / multi-negative losses \citep{oord2018a} introduces a permutation symmetry over the negative samples, but the loss is manifestly symmetric in them (a plain sum over negatives) and hence requires no alignment or assignment step.
\end{Remark}

\section{Loss Landscape and Parameter Space: Symmetries Puzzles}\label{sec:parameter-space-symmetries}
% zhao2025: indicate overlaps
In \Cref{sec:design-principle}, a symmetry carried by the problem or the training framework is deliberately built into the learning system.
Here, the group acts on the parameter space $\Theta$, and symmetries are rather unanticipated artifacts of the parametrization, often a source of confusion.
Neural networks are parametrized functions: the \emph{implementation} (or \emph{realization}) map
\begin{equation}\label{eq:implementation-map}
    \varphi:\Theta\to\mathcal{F}, \qquad \varphi(\theta)=f_\theta,
\end{equation}
relates $\Theta$ and the function space $\mathcal F$, through an architecture (left implicit) that takes $\theta$ as input.
Frequently $\varphi$ is not injective, i.e.\ \emph{non-identifiable}: distinct parameters can yield the same function.
A transformation $g$ of $\Theta$ can then be a symmetry of three loci (\Cref{fig:sym-loci}).
The most studied is the implementation map, $\varphi(g\cdot\theta)=\varphi(\theta)$ (model).
The second is the loss or the optimizer, $\mathcal L(g\cdot\theta)=\mathcal L(\theta)$ or $U(g\cdot\theta)=g\cdot U(\theta)$ for $U$ the update map giving $\theta_{t+1}=U(\theta_t)$ (optimization).
The third is a single parameter, $g\cdot\theta=\theta$ (symmetric state).
Parameter and functional spaces are jointly visualized in \Cref{fig:param-space-symmetries} b, in particular through the orbits, fibers and level sets they contain.
\Cref{sec:classifications-param-space-symmetries} breaks down the object axis of the taxonomy (\Cref{tab:taxonomy}) specifically for the parameter space.
\Cref{sec:descr-param-space-symmetries} is also \emph{descriptive}: it catalogs the symmetries and their consequences (identifiability, conservation laws, learning dynamics).
\Cref{sec:bug-or-feature} is, in contrast, \emph{prescriptive}: it gathers guidance to account for, remove or exploit them.
\Cref{sec:adjacent-fields} closes with weight-space learning and singular learning theory.
Almost all works treat these symmetries \emph{explicitly}.

\subsection{Classifications}\label{sec:classifications-param-space-symmetries}
Before describing the actual transformations and their consequences (\Cref{sec:descr-param-space-symmetries}), we organize their definitions along two axes, partly overlapping with the classification of \citet{zhao2025}.
The first axis forms a hierarchy and follows a chain of implications between loci of \Cref{fig:sym-loci}: parameter (symmetric state) $\to$ model (implementation map) $\to$ optimization (loss).
The second axis concerns scope: what may enter the definition of a transformation, and on which objects (parameters, data, architecture) conditions may be imposed.
It consists of variants of the functional symmetry of $\varphi$, which become \emph{extent-restricted} (\Cref{tab:taxonomy}) when restricted to a subset of $\Theta$ or $\mathcal X$.
\Cref{tab:param-double-entry} references works discussing parameter-space symmetries and their source, i.e.\ the ingredient (architecture, layer, activation, loss or optimizer) producing them.
One source can produce multiple symmetries, and the same symmetry can arise from different sources.

We propose in \Cref{fig:param-space-symmetries} a schematic description of the two axes, as transformations leaving invariant: (1) the parameter $\theta$ itself, (2) the function $f_\theta=\varphi(\theta)$, (3) a loss depending on $\theta$ only through $f_\theta$, (4) a general loss e.g.\ with regularizer $\mathcal{L}(\theta)=\mathcal{L}_F(f_\theta)+\mathcal{L}_R(\theta)$, (5) the restriction $f_\theta|_{\mathcal{X}_0}$ for a subset $\mathcal{X}_0\subset\mathcal{X}$ of the data, and (6) $f_\theta$ for $\theta$ in a subset $\Theta_0\subset\Theta$.

\subsubsection{Hierarchical Axis}\label{sec:hierarchy-param-space-symmetries}

We present the definitions following a natural inclusion order.

\paragraph{(SS) Symmetric States}
Let $G$ be a group of symmetries of $\varphi$ or $\mathcal{L}$ acting faithfully on $\Theta$.
A parameter $\theta$ is in a \emph{symmetric state} \citep{simsek2021,ziyin2025} if its stabilizer (\Cref{def:orbit-stabilizer}) is nontrivial, $G_\theta=\{g\in G : g\cdot\theta=\theta\}\neq\{e\}$.
Since $G_{g\cdot\theta}=gG_\theta g^{-1}$, being symmetric is a property of the orbit, and a larger stabilizer means a smaller orbit $G\cdot\theta\cong G/G_\theta$ (\Cref{thm:orbit-stabilizer}), as depicted in \Cref{fig:param-space-symmetries} b: the orbit of the more symmetric $\tilde \theta$ is smaller than the orbit of $\theta$.
For example, with the group of permutations, $\theta$ is unchanged when permuting neurons $i$ and $j$ exactly when the two neurons carry identical incoming and outgoing weights.

\paragraph{(PS) Parameter-Space Symmetries}\label{sec:PS_param-space-sym}
% a symmetric state for conjugation appears in [ravanbakhsh2017p3icml] see thread at LLM_saved-outputs/responses-saved/weight-tying_sym-state.md
$G$ is a group of parameter-space symmetries if $f_{g\cdot\theta}=f_\theta$ for all $g\in G$, i.e.\ if
\begin{equation*}
    g\cdot\theta \in \varphi^{-1}\big(\{\varphi(\theta)\}\big),
\end{equation*}
the set $\varphi^{-1}(\{\varphi(\theta)\})$ being called the \emph{fiber} of $\theta$: orbits remain within fibers.
This notion is sometimes referred to as \emph{observational equivalence}, \emph{non-identifiability} in statistics or \emph{reparametrization} invariance.
When the action of $G$ is transitive on fibers, capturing all the redundancy of the parametrization, $G$ is called \emph{complete} and $\theta$ is identifiable up to $G$ \citep{zhao2025}.

\paragraph{(LS) Loss-Level Symmetries}
$G$ is a loss symmetry if orbits remain within level sets of the loss,
\begin{equation*}
    g\cdot\theta \in \mathcal{L}^{-1}\big(\{\mathcal{L}(\theta)\}\big) \qquad \forall g \in G.
\end{equation*} 
Regularization and criterion symmetries (\Cref{sec:criterion-symmetries}) are two ways in which loss-level symmetries can exceed parameter-space symmetries, and the \emph{Q-symmetries} of \citet{ziyin2024a} formalize this notion.

\vspace{3mm}

The three conditions form a chain of implications for $g$,
\begin{equation}\label{eq:param-symmetry-hierarchy}
    \underbrace{g\cdot\theta=\theta}_{(SS)} \implies \underbrace{f_{g\cdot\theta}=f_\theta}_{(PS)} \overset{\star}{\implies} \underbrace{\mathcal{L}(g\cdot\theta)=\mathcal{L}(\theta)}_{(LS)},
\end{equation}
where the first implication is immediate and the second ($\star$) holds whenever the objective factors through the implementation map, $\mathcal{L}=\mathcal{L}_F\circ\varphi$, which requires in particular an unregularized loss.
Neither implication reverses.

\begin{Example}[Loss-level symmetry by regularization]\label{ex:loss-level-symmetry-regularization}
    A level set of $\mathcal{L}$ can carry non-constant $f_\theta$ when $\mathcal{L}$ does not factor through $\varphi$.
    For example, with a regularizer $\mathcal{L}(\theta)=\mathcal{L}_F(f_\theta)+\mathcal{L}_R(\theta)$, the second implication of \Cref{eq:param-symmetry-hierarchy} itself fails: a functional symmetry has no reason to preserve $\mathcal{L}_R$ (\Cref{sec:dynamical-sb}) and there might be a direction $v$ where the two parts compensate, i.e.\ $D_v\mathcal{L}_F=-D_v\mathcal{L}_R\neq 0$, with $D_v$ the directional derivative.
    Moving infinitesimally along $v$ stays in the level set yet changes the function.
    As a further illustration, following from $\theta$ a continuous functional symmetry ((2) in \Cref{fig:param-space-symmetries}) that decreases $\mathcal{L}_R$ yields $\theta'$ with $f_{\theta'}=f_\theta$ and $\mathcal{L}(\theta')<\mathcal{L}(\theta)$.
    Then, degrading the model by gradient ascent until the loss climbs back to $\mathcal{L}(\theta)$ produces $\theta''$ with $\mathcal{L}(\theta'')=\mathcal{L}(\theta)$ but $f_{\theta''}\neq f_\theta$.
\end{Example}

\begin{Remark}[False-class symmetry]
    Even when $\mathcal{L}$ factors through $\varphi$, a loss-level symmetry can arise, as in the cross-entropy case of \Cref{ex:false-class-symmetries}.
    In an overparametrized setup where the model has enough flexibility to change per-example false-class values while fixing the true class probability $\hat y_c$, multiple parameters whose predictions differ by such rearrangements are loss-equivalent without being functionally equivalent.
    Interestingly, these false-class logits have been used to compare models: \citet{zhuang2022pmls} measure model instability across retrainings as the fraction of test points on which two models disagree (\Cref{sec:data-ordering}), and \citet{godinot2024} use this principle as a fingerprinting technique.
    The distribution of false-class logits is not directly optimized by training, so the influence of initialization might be perceptible even after training.
\end{Remark}


\subsubsection{Scope Axis}\label{sec:scope-param-space-symmetries}
The scope axis starts from the functional symmetry (PS) of $\varphi$ (model locus, \Cref{sec:PS_param-space-sym}) and varies where the invariance must hold and what the transformation may act on: holding only on a subset of $\Theta$ or $\mathcal X$ makes the symmetry \emph{extent-restricted} (\Cref{tab:taxonomy}).

\paragraph{General PS}
To be unequivocal, we call \emph{general} a PS symmetry that holds for every $\theta\in\Theta$ of a fixed architecture.
The positive rescaling of ReLU networks, $(w_{\mathrm{in}}, w_{\mathrm{out}})\mapsto(\lambda w_{\mathrm{in}}, w_{\mathrm{out}}/\lambda)$ with $\lambda>0$ at each hidden neuron, is the classical example.

\paragraph{Instance-Level PS}
An \emph{instance-level} or \emph{hidden} \citep{grigsby2023p4icml} symmetry holds on a subset $\Theta_0\subset\Theta$.
These are local obstructions to the completeness of the group $G$.
An example is a neuron in a second hidden layer of a ReLU network whose incoming weights and bias are all non-positive.
Its preactivation is thus non-positive as well, so its incoming parameters can be perturbed (staying non-positive) and its outgoing weights modified arbitrarily without changing $f_\theta$.

\paragraph{Data-Dependent PS}
A \emph{data-dependent} symmetry preserves the function on a subset of the data: $f_{g\cdot\theta}|_{\mathcal{X}_0}=f_\theta|_{\mathcal{X}_0}$ for some $\mathcal{X}_0\subset\mathcal{X}$.
We give an example based on activation patterns.
\begin{Example}\label{ex:data-PS-GL-activation}
    For a shallow ReLU network $f_{(W_1,W_2)}(x)=W_2\sigma(W_1x)$ with $W_1\in\mathbb{R}^{h\times n}$ and $W_2\in\mathbb{R}^{m\times h}$, fix an activation pattern $A\subseteq[h]$ of the hidden layer and let $C_A:=\{z\in\mathbb{R}^h : z_i>0 \text{ for } i\in A\}$.
    Given $\mathcal{X}_0\subset\mathcal{X}$ with $W_1x\in C_A$ for all $x\in \mathcal{X}_0$, the set
    \begin{equation*}
        \Omega(\mathcal{X}_0,W_1):=\big\{M\in GL(|A|,\mathbb{R}) : MW_{1,A}x\in\mathbb{R}^{|A|}_{>0}\ \forall x\in \mathcal{X}_0\big\}
    \end{equation*}
    acts on $(W_1,W_2)$ by $(W_{1,A}, W_{2,A})\mapsto(MW_{1,A}, W_{2,A}M^{-1})$, where $W_{1,A}$ (resp.\ $W_{2,A}$) selects the rows of $W_1$ (resp.\ the columns of $W_2$) indexed by $A$ and the complementary blocks are untouched.
    Note that $\Omega(\mathcal{X}_0,W_1)$ is not a group: it depends on $W_1$ and collects the transformations admissible at $(W_1,W_2)$ in a groupoid.
\end{Example}
\citet{zhao2023a} provide another example: a \emph{nonlinear} data-dependent action $G\times(\Theta\times\mathcal{X})\to\Theta\times\mathcal{X}$ in which the transformation of the weights is itself a function of the input $x$.
The output is preserved at $x$ only (batch size one), but the action is defined for the whole $GL(h,\mathbb{R})$.

\paragraph{Architectural PS}\label{sec:architectural-sym}
An \emph{architectural} symmetry is allowed to change the architecture: if the parameter space is $\mathbb{R}^n$, we admit function-preserving maps $\mathbb{R}^n\to\mathbb{R}^m$ with $m\neq n$ into a different architecture (\Cref{fig:param-space-symmetries} b).
These are no longer group actions (see the groupoids of \Cref{sec:extent-restriction}).
Neuron splitting is the prototype: duplicating a neuron and splitting its outgoing weights between the two copies maps a width-$r$ network to a functionally identical width-$(r+1)$ one, creating a symmetric state in the enlarged architecture.

\paragraph{Extensions}\label{sec:param-catalog-extension}
A transformation may act on more than the weights: the quiver representation formalism $(W,f)$ of \citet{armenta2021a} includes both weights and activation functions.
This allows transformations of activation functions to be incorporated in the symmetry definition.
For example, rescaling the incoming weights of a ReLU neuron $v$ by a negative number $\tau_v<0$ and its outgoing weights by $1/\tau_v$ is a symmetry if the activation $\max (0, x)$ becomes $\tau_v\max(0, x/\tau_v)=\min(0,x)$.
Moreover, on the hierarchical axis, the classification error can be argued to form a further level below the loss.
Since the prediction is an argmax over the outputs of $f_\theta$, any transformation preserving the top class at each input preserves the error while altering the loss.
The error landscape studied by \citet{pittorino2022jsm} possesses more symmetries than the loss landscape.

\begin{figure}[!tbp]
\begin{adjustwidth}{-\extralength}{0cm}
    \centering
    \includegraphics[width=\linewidth]{images/param-space-symmetries.pdf}
\end{adjustwidth}
    \caption{\textbf{(a)} Inclusion of transformations applicable to (1) symmetric states $\left \{g: g\theta=\theta\right \}$, (2) general functional symmetry $\left \{g: f_\theta = f_{g\theta} \right \}$, (3) loss symmetry when the loss factors through the implementation map $\left \{g:\mathcal{L}(f_{g\theta})=\mathcal{L}(f_{\theta})\right \}$, (4) general loss symmetry $\left \{g:\mathcal{L}(g\theta)=\mathcal{L}(\theta) \right \}$, (5) data-dependent functional symmetry $\left \{g: {(f_{g\theta})}_{\vert \mathcal{X}_0}={(f_{\theta})}_{\vert \mathcal{X}_0}\right \}$ for $\mathcal{X}_0\subset \mathcal{X}$, and (6) instance-dependent functional symmetry $\left \{g: \theta\in \Theta_0 \Rightarrow f_{g\theta}=f_\theta \right \}$ for $\Theta_0\subset \Theta$.
    Each class is illustrated on an exemplar network $\mathcal{N}$ whose setup is shown at the top of the panel.
    Note that (3) and (4) can intersect (e.g.\ neuron permutations in an MLP with $\ell_p$ regularization for (4)), but none is included in the other.
    The illustration is not exhaustive: data-dependent and instance symmetries can act at the loss level or be combined.
    \textbf{(b)} Spaces and maps.
    The map $\varphi$ sends $\Theta$, and the parameter space $\Theta'$ of a larger architecture, to $\mathcal{F}$.
    The map $i$ restricts functions to $\mathcal{X}_0$, and $\operatorname{AGG}$ combines $\mathcal{L}_F$ and $\mathcal{L}_R$ into $\mathcal{L}$.
    In $\Theta$ and $\Theta'$, colors (2)--(5) mark the preimage, along the arrows, of the dot of the same color, i.e.\ the set reached from $\theta$ by the corresponding class of (a).
    Class (6) is not shown.
    For (1), the orbit $G\theta$ has two components, whereas that of the more symmetric $\tilde{\theta}$ has one.}
    \label{fig:param-space-symmetries}
\end{figure}

\subsection{Description of Parameter-Space Symmetries}\label{sec:descr-param-space-symmetries}
After categorizing parameter-space symmetries along two axes in the previous part, we look at concrete examples, starting with a catalog (\Cref{sec:catalog-param-space-symmetries}).
We then investigate their consequences, which split along the state axis of \Cref{tab:taxonomy}.
Parameter identifiability (\Cref{sec:parameter-identifiability}) is \emph{static}: it seeks to make the catalog complete, i.e.\ asks whether fibers of the implementation map $\varphi$ coincide with orbits.
Conservation laws (\Cref{sec:conservation-laws}) and learning dynamics (\Cref{sec:learning-dynamics}) are \emph{dynamic}: symmetries of the loss $\mathcal L$ (optimization) constrain training trajectories.
Conservation laws are quantities kept constant along trajectories by continuous symmetries.
One branch treats them as \emph{unknowns} and seeks all laws of an architecture.
Learning dynamics goes further, including discrete symmetries and symmetric states (parameter), usually starting from a \emph{known} symmetry.

This part focuses on understanding the symmetries themselves: its purpose is \emph{descriptive}, while \Cref{sec:bug-or-feature} will discuss how to account for or leverage them.

\subsubsection{Catalog}\label{sec:catalog-param-space-symmetries}
We list the main parameter-space symmetries, starting with discrete then continuous ones, and finally report approaches unifying multiple symmetries.
\Cref{tab:param-double-entry} collects symmetries and their source in a double-entry table where cells contain references.

\paragraph{Permutation}
The permutation of units within a layer is a ubiquitous parameter-space symmetry.
For a layer $x\mapsto W_2\sigma(W_1x)$ with pointwise nonlinearity $\sigma$,
\begin{equation*}
    W_2\sigma(W_1x)=W_2P^{\top}\sigma(PW_1x)
\end{equation*}
for any permutation matrix $P$: relabeling the hidden neurons, i.e.\ permuting the rows of $W_1$ (with their biases) and the columns of $W_2$ accordingly, leaves $f_\theta$ unchanged.
The same principle applies to filters in CNNs and to heads in multi-head attention~\citep{tran2025iclr}.
\citet{hecht-nielsen1990anc} gave an early group-theoretic account: for an MLP with hidden widths $d_1,\dots,d_L$, the permutation group has order $\prod_{l=1}^{L}d_l!$.

The symmetry also appears one level up.
In Mixture-of-Experts \citep{shazeer2017}, a gating score $s$ with components $s_i(x)=w_i^{\top}x+b_i$ estimates the contribution of expert $\mathcal{E}_i:\mathbb{R}^d\to\mathcal{Y}$.
The output, the weighted sum $\operatorname{Softmax}(s(x))^{\top}\mathcal{E}(x)$, is invariant to the simultaneous permutation of the $n$ experts and their gating rows.
Combined with the translation symmetry of the softmax (see translation symmetry below), \citet{tran2025anips} prove that the resulting group $G(n)=\mathbb{R}^{d+1}\times S_n$ is \emph{complete} in the sense of \Cref{sec:hierarchy-param-space-symmetries} for the mixture level.

\paragraph{Sign Flip}
Odd activations, $\sigma(-z)=-\sigma(z)$ such as $\tanh$, admit a discrete flip symmetry,
\begin{equation*}
    W_2\tanh(W_1x)=(-W_2)\tanh(-W_1x),
\end{equation*}
which holds neuron-wise: negating the incoming weights, bias, and outgoing weights of any single hidden unit preserves the function.
Each hidden layer thus carries a group $\{\pm 1\}^{d_l}$, which combines with the permutations above to give a discrete symmetry group of order $\prod_{l=1}^{L}2^{d_l}d_l!$ for a $\tanh$ MLP \citep{hecht-nielsen1990anc}.

\paragraph{Translation}
Following \citet{kunin2020}, a subset $\mathcal{A}$ of the parameters, with indicator vector $\boldsymbol{1}_{\mathcal{A}}$, possesses a translation symmetry if
\begin{equation*}
    \mathcal{L}(\theta)=\mathcal{L}(\theta+\alpha\boldsymbol{1}_{\mathcal{A}}) \qquad \forall\alpha\in\mathbb{R},
\end{equation*}
a one-parameter action of the noncompact additive group $(\mathbb{R},+)$.
The prototype is the softmax.
Since $\operatorname{Softmax}(z+\alpha\boldsymbol{1})=\operatorname{Softmax}(z)$, any parameter subset that shifts all logits equally (biases or a single column of the weight matrix in the final layer) can be translated freely.
The gating translation of Mixture-of-Experts encountered above is the same phenomenon.

\paragraph{Scale}
A subset $\mathcal{A}$ possesses a scale symmetry if \citep{kunin2020}
\begin{equation*}
    \mathcal{L}(\theta)=\mathcal{L}(\alpha_{\mathcal{A}}\odot\theta) \quad \forall\alpha>0, \qquad \text{where}\quad \alpha_{\mathcal{A}}:=\alpha\boldsymbol{1}_{\mathcal{A}}+\boldsymbol{1}_{\mathcal{A}^{c}}.
\end{equation*}
Normalization layers are a main source.
Batch normalization \citep{ioffe2015p3icml} standardizes each preactivation with batch statistics, $\operatorname{BN}_{\gamma,\beta}(x_i)=\gamma\odot\frac{x_i-\mu_B}{\sqrt{\sigma_B^2+\epsilon}}+\beta$ with $\mu_B=\frac{1}{m}\sum_{i=1}^{m}x_i$ and $\sigma_B^2=\frac{1}{m}\sum_{i=1}^{m}(x_i-\mu_B)^2$, hence $\operatorname{BN}(\lambda Wx)=\operatorname{BN}(Wx)$ for $\lambda>0$ (up to the stabilizing constant $\epsilon$).
Variants differ by the index set $\mathcal{S}_i$ over which the statistics are computed, and the granularity of $\mathcal{S}_i$ dictates the granularity of the symmetry.
Instance normalization \citep{ulyanov2017} normalizes per channel and per sample, yielding per-unit rescalings as in BN.
Layer normalization \citep{ba2016} pools all channels, leaving only the coarsest per-layer scalar rescaling.
Group normalization \citep{wu2018a} interpolates by splitting the $C$ channels into $M$ groups, $\mathcal{S}_i=\big\{k \mid k_N=i_N,\ \lfloor k_C/(C/M)\rfloor=\lfloor i_C/(C/M)\rfloor\big\}$ for indices $(i_N,i_C,i_H,i_W)$, giving one rescaling per group.
Weight normalization \citep{salimans2016anips} uses a radial-spherical reparametrization $w=r\frac{v}{\lVert v\rVert}$ of each unit's input weights, trained in the new variables $(r,v)$ with claimed benefits such as gradient magnitude stabilization.
Since $v$ and $\lambda v$ implement the same $w$ for $\lambda>0$, this reparametrization introduces a scale symmetry.

\paragraph{Rescale}
Two subsets $\mathcal{A}_1,\mathcal{A}_2$ possess a rescale symmetry if \citep{kunin2020}
\begin{equation*}
    \mathcal{L}(\theta)=\mathcal{L}\big(\alpha_{\mathcal{A}_1}\odot\alpha_{\mathcal{A}_2}^{-1}\odot\theta\big) \qquad \forall\alpha>0.
\end{equation*}
It occurs at every hidden neuron with $\mathcal{A}_1$ the incoming and $\mathcal{A}_2$ the outgoing weights for a positively homogeneous activation ($\mathrm{ReLU}$, leaky $\mathrm{ReLU}$, identity): $(w_{\mathrm{in}},w_{\mathrm{out}})\mapsto(\lambda w_{\mathrm{in}},w_{\mathrm{out}}/\lambda)$ (the classical example of \Cref{sec:scope-param-space-symmetries}).
It gives a group $(\mathbb{R}_{>0})^{d_l}$ per layer.
Degree-$r$ homogeneity skews the pairing: for a monomial activation $z^r$, the symmetry becomes $(w_{\mathrm{in}},w_{\mathrm{out}})\mapsto(\lambda w_{\mathrm{in}},\lambda^{-r}w_{\mathrm{out}})$, i.e.\ $\mathcal{L}(\theta)=\mathcal{L}(\alpha_{\mathcal{A}_1}\odot\alpha_{\mathcal{A}_2}^{-r}\odot\theta)$.
The symmetry is also found in lightning self-attention, which satisfies $(K^{\top}Q,V)\mapsto(\lambda K^{\top}Q,\lambda^{-1}V)$ in addition to its other symmetries \citep{henry2025}.

Although rescaling preserves $f_\theta$, it alters internal quantities such as feature maps, gradients, and measures of flatness \citep{armenta2021b}, a point we return to in \Cref{sec:learning-dynamics,sec:bug-or-feature}.

\paragraph{General Linear Group}
Whenever linear maps compose without an intervening nonlinearity, the factorization is defined only up to $AB=(AM^{-1})(MB)$ for $M$ invertible.
In deep linear networks $f_\theta(x)=W_L\cdots W_2W_1x$ \citep{saxe2014piclr2}, every intermediate width carries this symmetry,
\begin{equation*}
    (W_L,\dots,W_{l+1}M^{-1},\,MW_l,\dots,W_1), \qquad M\in GL(d_l,\mathbb{R}).
\end{equation*}
Transformers contain the same structure inside each attention layer: the products $W_K^{\top}W_Q$ and $W_OW_V$ (with $W_O$ the output projection) are invariant under a $GL(d_k,\mathbb{R})$ action on the factors \citep{kobayashi2024anips}.
Low-rank adapters \citep[LoRA,][]{hu2021} inherit it as well: the update $\Delta W=AB^T$ with $A\in \mathbb R^{m\times r}, B\in \mathbb R^{n\times r}$ is invariant under $(A, B)\mapsto(AM, BM^{-T})$ for $M\in GL(r,\mathbb{R})$ \citep{yen2025iclr}.
This symmetry, or its orthogonal restriction, is also called \emph{double rotation}.

\paragraph{Orthogonal and Rotation}
Orthogonal symmetries often arise as the residue of a broken $GL$ symmetry: an $L_2$ penalty is not invariant under the full factorization group, but its subgroup preserving norms survives.
\citet{kunin2019p3icml} work this out for linear autoencoders, with encoder $W_1\in\mathbb{R}^{k\times d}$ and decoder $W_2\in\mathbb{R}^{d\times k}$ trained on $\mathcal{L}(W_1,W_2)=\lVert X-W_2W_1X\rVert_F^2$: the unregularized problem has symmetry group $GL(k,\mathbb{R})$, but adding $\lambda(\lVert W_1\rVert_F^2+\lVert W_2\rVert_F^2)$ breaks the scaling part and reduces the group to $O(k)$.
\citet{kobayashi2024anips} observe the same reduction $GL(d_k,\mathbb{R})\downarrow O(d_k)$ for weight-decayed attention layers.
\citet{bamler2018p3icml} exhibit the phenomenon in dynamic word embeddings: a per-time-step objective satisfies $\ell(X_t,R_tZ_t)=\ell(X_t,Z_t)$ for rotations $R_t$ of the embeddings $Z_t$, but a coupling regularizer $\psi(Z)=\frac{\lambda}{2}\sum_{i=1}^{N}\sum_{t=1}^{T-1}\lVert z_{t+1,i}-z_{t,i}\rVert_2^2$ in the total loss $\mathcal{L}(Z)=\sum_{t=1}^{T}\ell(X_t,Z_t)+\psi(Z)$ breaks the independent tuple $(R_1,\dots,R_T)$ down to a single global rotation $R_1=\dots=R_T$.

Rotations also index dynamics: the saddle-to-saddle trajectory of deep linear networks is characterized up to rotation \citep{jacot2022}.
Finally, an orthogonal symmetry can hold exactly, without regularization, for non-pointwise activations: the radial rescaling activation $\sigma(z)=f(\lVert z\rVert)z$ satisfies $\sigma(gz)=g\sigma(z)$, so $W_2g^{-1}\sigma(gW_1x)=W_2\sigma(W_1x)$ for all $g\in O(h)$ \citep{zhao2023a}.

\paragraph{Overparametrization}
Symmetries appear when a group of neurons can be emulated by a smaller subset.
The converse classical notion is \emph{minimality} \citep{sussmann1992nn,fukumizu1996nn,brea2019,simsek2021,grillo2026,gegenfurtner2026}: a network $f_\theta$ with $H$ hidden units is minimal if for any network $\tilde f$ with $\tilde{H}<H$ hidden units there exists $x$ with $f_\theta(x)\neq \tilde f(x)$.
\citet{fukumizu1996nn} showed that, for sigmoidal activations, the Fisher information matrix is positive definite if and only if the network is minimal, tying overparametrization to degenerate local geometry.
At non-minimal parameters the fiber grows beyond the orbit of the general symmetry group.
\citet{sareen2026} describe overparametrization mechanisms as width-changing transformations: in linear networks the factorization symmetry extends to $(W_L,\dots,W_{l+1}M^{\dagger},MW_l,\dots,W_1)$ for a full-column-rank $M\in\mathbb{R}^{r\times d_l}$ with $r>d_l$ and $M^{\dagger}$ a left inverse.
For general activations, possible width-changing transformations are \emph{duplication} (post-activation splitting), where $M$ duplicates rows of the identity and $M^{\dagger}$ recombines them with coefficients summing to $1$, and \emph{0-type neurons}, groups of neurons implementing the zero function.
Homogeneous activations additionally allow \emph{pre-activation splitting}.
These are architectural symmetries, forming a groupoid (\Cref{sec:extent-restriction}) rather than a group: expansion is always possible but merging is only defined on duplicated configurations.
Non-minimality can also induce instance symmetries, in particular at symmetric states.
When two neurons have identical incoming weights, as for permutation symmetric states, a $d_{l+1}$-dimensional additive group acts freely ($d_{l+1}$ the width of the next layer): the paired outgoing weights can be varied arbitrarily as long as their sums stay constant \citep{brea2019}.
At a rescale-symmetric state ($\theta_v=0$ for a neuron $v$), additive groups of dimensions $d_{l+1}$ and $d_{l-1}$ separately act freely.
From the second hidden layer of a ReLU MLP on, a unit $v$ whose incoming parameters $\theta^+_v$ are all non-positive causes non-minimality since it does not contribute to the implemented function.
The scaling $D\theta^+_v$, with $D$ a diagonal matrix with strictly positive diagonal, constitutes a symmetry.

\begin{Remark}[Overparametrization by width and depth]
    Here, we discuss overparametrization by width.
    Overparametrization by depth also exists, for example in linear networks.
    However, depth alone cannot cause overparametrization for a generic architecture: \citet{grillo2026} show that in ReLU networks additional layers keep increasing expressivity.
\end{Remark}

\begin{sidewaystable}[p]
    \centering
    \scriptsize
    \setlength{\tabcolsep}{2pt}
    \renewcommand{\arraystretch}{1.2}
    \begin{tabularx}{\textheight}{|>{\raggedright\arraybackslash\bfseries}p{2.0cm}|*{10}{>{\raggedright\arraybackslash}X|}}
        \hline
        \textbf{Source} & \textbf{Permutation} & \textbf{Rescaling} & \textbf{Sign flip} & \textbf{Scale} & \textbf{Translation} & \textbf{Orthogonal} & \textbf{$GL(n)$} & \textbf{Cross-\-architectural} & \textbf{Data-\-dependent} & \textbf{Instance / hidden} \\
        \hline
        Feedforward networks & \citep{hecht-nielsen1990anc,brea2019,tatro2020,entezari2022,ainsworth2023,jordan2023,navon2023p4icml,chen2024jsm,lim2024,nurisso2024a,kalogeropoulos2024,simsek2021,zhang2025p4icml,ziyin2025,godfrey2022,ziyin2024} & \citep{ziyin2024a,ziyin2025a,ziyin2025} & \citep{ziyin2025,ziyin2024} & \citep{ziyin2024a,ziyin2025a,ziyin2025} & \citep{ziyin2025} & \citep{ziyin2024a,ziyin2025a,ziyin2025,ziyin2024} &  & \citep{fukumizu1996nn} & \citep{zhao2023a} &  \\
        \hline
        Transformers & \citep{zhang2025p4icml,tran2025iclr} &  &  &  &  & \citep{zhang2025p4icml,kobayashi2024anips} & \citep{henry2025,kobayashi2024anips} &  &  &  \\
        \hline
        Multi-head attention & \citep{tran2025iclr} &  &  &  &  & \citep{zhang2025p4icml} (per head) & \citep{tran2025iclr} &  &  &  \\
        \hline
        Self-attention &  & \citep{zhang2025p4icml} &  &  &  & \citep{zhang2025p4icml,kobayashi2024anips,ziyin2024a,ziyin2025} (double rotation) & \citep{kobayashi2024anips} &  &  &  \\
        \hline
        Lightning self-attention &  & \citep{henry2025} (within-layer) &  &  &  &  & \citep{henry2025} (between- and within-layer) &  &  &  \\
        \hline
        LoRA &  &  &  &  &  &  & \citep{yen2025iclr} &  &  &  \\
        \hline
        BatchNorm &  &  &  & \citep{badrinarayanan2015,huang2017,kunin2020,li2020anips,armenta2021b,tanaka2021,yi2022ptcuai,jordan2023,lyu2023,ziyin2024a} &  &  &  &  &  &  \\
        \hline
        BatchNorm variants &  &  &  & \citep{wu2018a} (GN, LN, IN), \citep{kunin2020} (LN, WN), \citep{lyu2023} (LN), \citep{ziyin2024a} (LN), \citep{li2020anips} (GN, WN), \citep{salimans2016anips} (WN), \citep{tanaka2021} (GN, LN) &  &  &  &  &  &  \\
        \hline
        Softmax / cross entropy &  &  &  &  & \citep{kunin2020,tanaka2021,tran2025anips,ziyin2025} &  &  &  &  &  \\
        \hline
        Linear networks &  &  &  & \citep{saxe2014piclr2} &  & \citep{jacot2022,saxe2014piclr2,li2018,kunin2019p3icml,ziyin2024a} & \citep{godfrey2022,saxe2014piclr2,arora2018p3icml,kunin2019p3icml,zhao2025a} & \citep{jacot2022} &  &  \\
        \hline
        ReLU networks & \citep{godfrey2022,phuong2019,bona-pellissier2023,freeman2017,pittorino2022jsm,lim2024,nurisso2024a,elisendagrigsby2025aim,grillo2026,gegenfurtner2026,rolnick2020p3icml} & \citep{rolnick2020p3icml,godfrey2022,phuong2019,bona-pellissier2023,neyshabur2015anips,dinh2017,freeman2017,huang2017,badrinarayanan2015,liang2019pticais,meng2019,tsuzuku2019,yi2019,kunin2020,armenta2021b,kwon2021p3icml,petzka2021anips,pittorino2022jsm,saul2023,zhao2023,kalogeropoulos2024,lim2024,nurisso2024a,zhao2024,elisendagrigsby2025aim,zhao2025a,nurisso2025,ziyin2025,grillo2026,gegenfurtner2026} &  &  &  &  &  &  &  & \citep{rolnick2020p3icml,grigsby2023p4icml,gegenfurtner2026} \\
        \hline
        Homogeneous activation & \citep{godfrey2022} (leaky ReLU), \citep{tanaka2021} & \citep{godfrey2022} (leaky ReLU), \citep{kunin2020,saul2023,zhao2023,zhao2024,kalogeropoulos2024,nurisso2025}, \citep{usevich2025} ($\sigma(\alpha x)=\alpha^{r}\sigma(x)$) &  &  &  &  &  &  &  &  \\
        \hline
        Polynomial activation & \citep{godfrey2022,usevich2025} & \citep{godfrey2022,usevich2025} &  & \citep{henry2025} &  &  & \citep{henry2025} &  &  &  \\
        \hline
        Tanh, odd activation & \citep{fefferman1993anips} &  & \citep{fefferman1993anips,kalogeropoulos2024} &  &  &  &  &  &  &  \\
        \hline
        Radial activation & \citep{godfrey2022} &  & \citep{godfrey2022} (diagonal with $\pm1$ entries) &  &  & \citep{zhao2023a} &  &  &  &  \\
        \hline
        Mixture-of-experts & \citep{tran2025anips} (experts) &  &  &  & \citep{tran2025anips} (gating params) &  &  &  &  &  \\
        \hline
        Quiver NN & \citep{armenta2021a} & \citep{armenta2021a}, \citep{saul2023} (DAG), \citep{nurisso2025} (DAG) &  & \citep{armenta2021a} &  &  &  & \citep{armenta2021a} &  &  \\
        \hline
        Word embedding &  &  &  &  &  & \citep{bamler2018p3icml,tanaka2021} &  &  &  &  \\
        \hline
        Over\-para\-metri\-zation by width &  &  &  &  & \citep{brea2019,simsek2021} &  &  & \citep{fukumizu1996nn,jacot2022,simsek2021,sareen2026} &  &  \\
        \hline
        CNN filters & \citep{godfrey2022,entezari2022,pittorino2022jsm,ainsworth2023,jordan2023,chen2024jsm} & \citep{pittorino2022jsm} &  &  &  &  &  &  &  &  \\
        \hline
        Classification error (argmax) &  &  &  & \citep{pittorino2022jsm} & \citep{pittorino2022jsm} &  &  &  &  &  \\
        \hline
        Optimizers & \citep{yang2025} & \citep{neyshabur2015anips,meng2019,huang2017} &  & \citep{zhuang2022tmlr,badrinarayanan2015,kingma2015,yi2022ptcuai} &  & \citep{bamler2018p3icml} & \citep{amari1998nc,zhao2023,zhao2024,yen2025iclr,kristiadi2023anips} &  &  &  \\
        \hline
    \end{tabularx}
    \caption{Double-entry table of parameter-space symmetries: source (rows: architecture, layer, activation, loss or optimizer) against type of symmetry (columns). Most sources yield symmetries of $\varphi$ (model locus of \Cref{fig:sym-loci}), while losses and optimizers yield symmetries of $\mathcal L$ or $U$ (optimization locus).}
    \label{tab:param-double-entry}
\end{sidewaystable}

\paragraph{Unifying Works}\label{sec:unifying-works-param-space-symmetries}
Rather than treating the symmetries above as a list of unrelated case studies, several works view them as instances of a single formalism.

Unification at the level of the network is proposed by \citet{armenta2021a}, with the quiver formalism of \Cref{sec:extent-restriction,sec:param-catalog-extension}, in which thin quiver isomorphisms $\tau: (W,f)\to (V,g)$ define isomorphisms of neural networks.
\citet{godfrey2022} work at the level of a single nonlinear layer, through the \emph{intertwiner group} of an activation $\sigma_n:\mathbb{R}^n\to\mathbb{R}^n$,
\begin{equation*}
    G_{\sigma_n}=\big\{A\in GL(n,\mathbb{R}) : \exists B\in GL(n,\mathbb{R}),\ \sigma_n\circ A=B\circ\sigma_n\big\},
\end{equation*}
the largest group of linear transformations that can be undone after the nonlinearity.
They report the intertwiner groups of $6$ activations (identity, sigmoid, ReLU, LeakyReLU, RBF and monomial) and the corresponding matrices $B$.

At the level of the loss, \citet{kunin2020} frame differentiable symmetries in the form $\mathcal{L}(\theta)=\mathcal{L}(\psi(\theta,\alpha))$ for a one-parameter transformation $\psi$ with $\psi(\theta,0)=\theta$.
Translation, scale, and rescale are recovered by taking $\psi$ additive, multiplicative, or paired multiplicative-inverse on a subset $\mathcal{A}$ of parameter coordinates.
For $O$ with orthonormal columns and reflection $R=I-2OO^{\top}$, \citet{ziyin2024} define the \emph{O-mirror symmetry} as $\mathcal{L}(\theta)=\mathcal{L}(R\theta)$ for all $\theta\in\Theta$.
Rotation, permutation and sign flip all imply O-mirror symmetries for a suitable choice of $O$.
\citet{ziyin2024a} formalize a general notion (reused in \citep{ziyin2025a}) of \emph{Q-symmetry} as the case $\ell(x,\theta)=\ell(x,Q_s(\theta))$ for a one-parameter family $Q_s$ with $Q_0=\operatorname{id}$.
Within this class, the \emph{exponential symmetries} are those with a linear infinitesimal generator $Q_s=I+sA+o(s)$ for symmetric $A$.
The group element is given by exponentiation, $e^{\lambda A}:=\lim_{s\to0}(I+sA+o(s))^{\lambda/s}$, and acts on $\theta$ by matrix-vector multiplication.
Rescale, scale, and matrix factorization are all exponential symmetries for a suitable $A$.

\subsubsection{Parameter Identifiability}\label{sec:parameter-identifiability}
The catalog of \Cref{sec:catalog-param-space-symmetries} raises a question: is it complete?
A group $G$ of parameter-space symmetries only guarantees the inclusion of orbits into fibers (\Cref{sec:hierarchy-param-space-symmetries}), and $\theta$ is \emph{identifiable up to $G$} when the inclusion is an equality,
\begin{equation*}
    \varphi^{-1}\big(\{\varphi(\theta)\}\big)=G\cdot\theta.
\end{equation*}
An identifiability result thus certifies that a list of known symmetries is exhaustive.
Instance-level symmetries of \Cref{sec:scope-param-space-symmetries} are where exhaustivity can fail locally.
Identifiability underlies what can in principle be recovered about the weights by querying a black-box model.
See \citet{zhao2025} for a more focused account.

Early results concern sigmoidal networks and are quite rigid.
For smooth nonlinear activations $\sigma$, works by \citet{albertini1993,fefferman1993anips} revealed that the architecture and weights are generically determined up to permutations (and sign flips for odd activations such as $\mathrm{tanh}$).
A local result was also given by \citet{fukumizu1996nn}: by relating the Fisher information matrix to minimality (\Cref{sec:catalog-param-space-symmetries}, overparametrization), they gave a criterion guaranteeing identifiability in a small neighborhood around the parameter.

\paragraph{Linear Region Boundaries of ReLU Networks}
For ReLU networks, the function $f_\theta$ is continuous piecewise linear, and its non-differentiability locus (the boundaries between linear regions) is a fingerprint of the parameters: each hidden neuron contributes a ``bent hyperplane'' (the zero set of its preactivation), whose direction and bendings characterize $f_\theta$.
\citet{rolnick2020p3icml} turn this observation into a recoverability result: when boundaries are connected and each intersects the boundaries of all previous layers, the architecture, weights, and biases can be reconstructed up to permutation and rescaling.
They further show that boundaries can be approximated by querying the network, and algorithmically reverse-engineer ReLU networks from output queries alone.
\citet{phuong2019} prove identifiability for non-increasing architectures: fixing $(d_0,\dots,d_L)$ with $d_0\ge d_1\ge\dots\ge d_{L-1}\ge 2$ and $d_L=1$, if $\varphi(\theta)=\varphi(\eta)$ for general $\theta$ and $\eta$ satisfying non-degeneracy conditions, then there exist permutation matrices $P_l$ and positive diagonal matrices $D_l$ transforming $\theta$ into $\eta$, e.g.\ on a hidden layer $1<l<L$,
\begin{equation*}
    W_l^{\eta}=D_lP_l\,W_l^{\theta}\,P_{l-1}^{-1}D_{l-1}^{-1}:
\end{equation*}
permutations and rescalings are the only general symmetries.
\citet{bona-pellissier2023} reach the same conclusion under weaker and different conditions.

\paragraph{Instance-Level Symmetries of ReLU Networks}
In general, fibers can strictly exceed orbits.
\citet{grigsby2023p4icml} identify four mechanisms producing instance-level (\emph{hidden}) symmetries.
For the example of ``never co-active pairs'', let $H_v:=\{x\in\mathbb{R}^{d_0} : z_v(x)>0\}$ denote the active region in input space of a neuron $v$.
If $u$ in layer $l-1$ and $v$ in layer $l$ satisfy $H_u\cap H_v=\emptyset$, the weight $w_{uv}$ never contributes to the function and can generically be perturbed (locally) without changing it.
Empirically, they observe that hidden symmetries become more frequent with depth and rarer with width.
\citet{grigsby2023p4icml} also prove that for architectures at least as wide as the input ($d_0\le d_l$ for all hidden $l$), a positive-measure subset of $\Theta$ has no hidden symmetries.
Together with the complementary result of \citet{phuong2019}, this shows that monotone architectures admit large sets of identifiable parameters.
For three-layer bottleneck architectures $(d_0,d_1,d_2,d_3)$ with $d_1\le d_0$, \citet{gegenfurtner2026} characterize the fibers completely.
Since the first-layer map is then generically surjective onto the nonnegative orthant $\mathbb{R}^{d_1}_{\ge0}$, the analysis reduces to one-hidden-layer networks on this domain.
They encode such a function $f_\theta: \mathbb R^{d_1}_{\geq 0}\to \mathbb R^{d_2}$ by its \emph{hyperplane representation}, a weighted hyperplane arrangement together with a reference affine map such that $f_\theta=f_\eta$ if and only if the representations coincide.
Partitioning hidden neurons into those acting \emph{linearly} and \emph{nonlinearly} on the orthant, they uncover new instance-level symmetries created by linear neurons, and derive conditions for generic bottleneck network identifiability.
They also notably show that symmetries can appear from the composition of layers, making non-identifiability non-localizable.

\paragraph{Functional Dimension}
A quantitative measure of identifiability is the \emph{functional dimension} of \citet{elisendagrigsby2025aim}.
For a batch $Z=(z_1,\dots,z_k)$ of inputs, the \emph{evaluation map} $E_Z:\Theta\to\mathbb{R}^{kd_L}$, $\theta\mapsto(f_\theta(z_1),\dots,f_\theta(z_k))$, probes the function space through finitely many queries, and
\begin{equation*}
    \dim_{\mathrm{fun}}(\theta):=\sup_{Z}\ \operatorname{rank}\,\partial E_Z(\theta)
\end{equation*}
where $\partial E_Z$ denotes the Jacobian of the evaluation map, counts degrees of freedom of $f_\theta$ under infinitesimal perturbations of $\theta$.
For ReLU networks, the functional dimension is upper bounded by the parameter count minus one degree of freedom per hidden neuron, accounting for the rescaling symmetries.
Crucially, the functional dimension is inhomogeneous across parameter space, dropping with instance-level symmetries, which inflate the fiber.
\citet{grillo2026} further define its upper bound as the \emph{expected functional dimension}, $\dim(\Theta/\!\sim)$, the dimension of the parameter space quotiented by permutations and rescalings, and prove that every ReLU architecture (with input and hidden widths at least $2$) contains an open set of parameters in $\Theta$ attaining it.

\paragraph{Polynomial Networks and Attention}
\citet{henry2025} study \emph{lightning self-attention}, $X\mapsto VXX^{\top}K^{\top}QX$, a normalization-free variant of self-attention that is polynomial in its weights.
For a single layer, the generic fiber is exactly the orbit of two symmetries: the group $GL(d_k,\mathbb{R})$ acting on the factors of the attention matrix $A=K^{\top}Q$ where $Q, K\in \mathbb R^{d_k\times d}$ (typically $d_k<d$), and the one-parameter rescaling $(A,V)\mapsto(\lambda A,\lambda^{-1}V)$, $\lambda\neq0$ (deep networks carry an additional between-layer $GL$ symmetry).
The dimension of the \emph{neuromanifold} $\mathcal{M}=\varphi(\Theta)$, the set of realizable functions, is then computed from the dimension of the generic fiber.
With the normalization restored, single-layer fibers are generically singletons. The authors conjecture, with numerical support, generic identifiability in the deep case.
For networks with monomial activation $z\mapsto z^{r}$, \citet{usevich2025} study \emph{Finite Identifiability} (FI) at generic $\theta$: whether the fiber $\{\eta\in\Theta : f_\eta=f_\theta\}$ consists of finitely many orbits of permutations and rescalings.
They prove a local-to-global theorem: if every shallow sub-architecture $\big(\min\{d_0,\dots,d_{l-1}\},d_l,d_{l+1}\big)$ is FI, so is the full network.
Thus, pyramidal architectures ($d_0\ge\dots\ge d_{L-1}\ge2$) and networks whose degree grows linearly with the width ($r_l\ge 2d_l-1$) are FI.
\citet{shahverdi2026} treat generic polynomial activations, for which the rescaling symmetry disappears, and their relation to the neuromanifold $\mathcal{M}$: generic MLP fibers are finite and therefore $\dim\mathcal{M}=\dim\Theta$.
Moreover, generic CNN parametrizations are one-to-one.
Non-identifiability is carried by \emph{subnetworks}, parameters where the outgoing weights of some neurons vanish ($|A_i|$ deactivated neurons in layer $i$): strict subnetworks (with vanishing incoming weights too) have fibers of dimension at least $\sum_{i=1}^{L-1}|A_i|d_{i-1}$ and generate the singular locus of the neuromanifold, which could be an attractive region of training~\citep{marchetti2025a}.

\subsubsection{Conservation Laws (CL) and Conserved Quantities (CQ)}\label{sec:conservation-laws}

The continuous symmetries of \Cref{sec:catalog-param-space-symmetries} do more than shape the loss landscape: they constrain the training trajectories themselves.
Since the loss is constant along a symmetry orbit, there are directions along which the gradient has null projection.
In analogy with Noether's theorem in mechanics, this integrates into a \emph{conservation law}: a function of the parameters whose value, the \emph{conserved quantity}, is fixed at initialization and retained throughout training.
We briefly explain the mechanism by which CL emerge, then review the literature.

\paragraph{From Continuous Symmetry to Conservation Law}
A function $h:\Theta\to\mathbb{R}$ is \emph{conserved} by a flow $\dot\theta=\chi(\theta)$ if $h(\theta(t))=h(\theta(0))$ for every trajectory and every $t$, or equivalently if $\langle\nabla h(\theta),\chi(\theta)\rangle=0$ for all $\theta$.
In the simplest case, the flow is the usual gradient flow $\chi=-\nabla\mathcal{L}$.
Let $(\psi_\alpha)_\alpha$ be a one-parameter symmetry, defined for $\alpha$ in a neighborhood of $0$ with $\psi_0=\operatorname{id}$, i.e.\ $\mathcal{L}(\psi_\alpha(\theta))=\mathcal{L}(\theta)$, and let
\begin{equation*}
    \xi(\theta):=\frac{d}{d\alpha}\psi_\alpha(\theta)\Big\vert_{\alpha=0}
\end{equation*}
be its \emph{generator}, the vector field giving at each $\theta$ the local direction of the symmetry curve.
Differentiating the invariance at $\alpha=0$ yields $\langle\nabla\mathcal{L}(\theta),\xi(\theta)\rangle=0$, hence along the gradient flow
\begin{equation*}
    \langle \xi(\theta(t)), \dot\theta(t)\rangle=0 \qquad \forall t:
\end{equation*}
the trajectory never moves along the symmetry direction.
This orthogonality becomes a conservation law when the generator is a gradient field, $\xi=\nabla C$, since then $\frac{d}{dt}C(\theta(t))=\langle\nabla C(\theta(t)),\dot\theta(t)\rangle=0$.
The existence of $C$ can be characterized by the exactness of the differential $1$-form $\xi^\flat$.
The Poincar\'e lemma states that, on a simply connected domain such as a standard parameter space $\Theta=\mathbb{R}^n$, every closed $1$-form is exact:
\begin{equation*}
    \xi^\flat \text{ closed} \iff \exists\, C:\Theta\to\mathbb{R} \text{ such that } \nabla C=\xi,
\end{equation*}
and by properties of the exterior derivative, closedness reduces to the symmetry of the Jacobian $\partial\xi(\theta)$.
Taking concrete examples from \citep{kunin2020}, the translation generator $\xi=\boldsymbol{1}_{\mathcal{A}}$ integrates to $C(\theta)=\langle\theta,\boldsymbol{1}_{\mathcal{A}}\rangle$, the scale generator to $\lVert\theta_{\mathcal{A}}\rVert^2$, and the rescale generator to $\lVert\theta_{\mathcal{A}_1}\rVert^2-\lVert\theta_{\mathcal{A}_2}\rVert^2$.
See \citep{kunin2020,zhao2023a,marcotte2023,ziyin2024a} for more details and \citet[Prop.~52]{gegenfurtner2026} for an example of a continuous symmetry that induces no conserved quantity.

\paragraph{Conserved Quantities and Convergence}
Conservation laws entered deep learning theory as an analytical device for linear networks.
\citet{saxe2014piclr2} analyze the nonlinear learning dynamics of deep linear networks: after an SVD-aligned change of variables, the gradient flow decouples into independent modes $(a,b)$, constrained by the conserved imbalance $a^2-b^2$ to a hyperbola in the $(a,b)$ plane.
\citet{tarmoun2021p3icml} quantify how conserved quantities control convergence speed: for asymmetric matrix factorization with a spectrally aligned initialization, the $i$-th eigencomponent of the product matrix converges at rate $O(e^{-t\sqrt{4\sigma_i^2+\lambda_i^2}})$, where $\sigma_i$ is the $i$-th singular value of the ground truth and $\lambda_i>0$ is a conserved quantity growing with the imbalance at initialization.
Strong imbalance accelerates convergence, and the order in which eigencomponents are learned also depends on $\lambda_i$.
\citet{zhao2023a} similarly show that the loss trajectory depends explicitly on the conserved quantity $Q$, i.e.\ $\partial_Q\mathcal{L}(\theta(t))\neq0$: the value of the loss at a given time depends on $Q$.
\citet{du2018anips} exploit approximate conservation under gradient \emph{descent}.
A challenge in low-rank asymmetric matrix factorization is that $(A,B) \mapsto \lVert AB^\top-M^*\rVert_F^2$ has no globally Lipschitz gradient due to the rescale symmetry.
However, with stepsizes decreasing sufficiently fast, the balancedness quantities remain bounded, which allows convergence to a bounded global optimum to be proved.
The conserved quantities for homogeneous networks are also derived, namely for ReLU networks: $\frac{d}{dt} \left( \|W^{(h)}[i, :]\|^2 - \|W^{(h+1)}[:, i]\|^2 \right) = 0$.

\paragraph{Finding All Conservation Laws}
Given an architecture, one may ask whether the known conserved quantities are the only ones.
\citet{marcotte2023} formalize the question through the \emph{flow field} $W(\theta):=\operatorname{span}\{\nabla\mathcal{L}(\theta;X,Y),\ \forall X,Y\}$, the space of all possible gradients: $h$ is conserved if and only if $\nabla h(\theta)\perp W(\theta)$ for all $\theta$.
They propose two algorithms to find all polynomial conservation laws and count the exact number of independent laws through the Lie algebra generated by the flow field, certifying that for shallow linear and ReLU networks no conservation law exists beyond the known ones.
\citet{marcotte2025p4icml} extend to ResNets and Transformers: for shallow architectures, the number of laws is governed by the number of hidden components (channels or heads), and residual connections $x\mapsto f_\theta(x)+x$ do not alter them.
For stacked architectures, they frame conservation laws by blocks $h_k:\Theta_k\subset\Theta\to\mathbb{R}$ depending only on the parameters of one block, and prove that these are exactly the laws of the building block in isolation.

\paragraph{Beyond Idealized Gradient Flow}
Training heuristics perturb the conservation mechanism.
\citet{kunin2020} derive a modified gradient flow accounting for finite learning rate, momentum and weight decay, under which CL break.
With weight decay $\lambda$, for instance, the translation law decays as $\langle\theta(t),\boldsymbol{1}_{\mathcal{A}}\rangle=e^{-\lambda t}\,c$ with $c=\langle\theta(0),\boldsymbol{1}_{\mathcal{A}}\rangle$ the quantity conserved when $\lambda=0$ (see \Cref{sec:dynamical-sb}).
\citet{marcotte2025p4icml} prove that weight decay does not create new laws, by showing that any law under weight decay maps to a law without it.
They also prove that Adam-like optimizers typically break conservation laws, and that discretization induces a drift of order $O(\eta^2 k)$ after $k$ steps of size $\eta$, so laws hold approximately for small learning rates.
Instead, \citet{saul2023} restore CL: they derive a \emph{weight-balancing flow} under which a balanced network stays exactly balanced while (under assumptions) still descending the regularized loss (see \Cref{sec:quotient-repr-align}).
Gradient noise breaks conservation in a more structured way.
For exponential symmetries, whose CQ takes the quadratic form $C(\theta)=\theta^\top A\theta$, \citet{ziyin2024a} show a drift under SGD noise that defines a flow on the symmetry orbit, admitting (under conditions) a unique attractive fixed point.
Thus, a value of the CQ is selected independently of initialization.
For example, with scale-invariant losses, since the noise-free gradient is tangent to the sphere of radius $\lVert\theta\rVert$, the parameters are attracted to $\lVert\theta\rVert=\infty$.

\paragraph{Geometry of the Reachable Set}
CQ are fixed at initialization, so gradient flow is confined to their level set, whose geometry then delimits what training can reach.
For shallow ReLU networks, \citet{nurisso2024a} show that this \emph{invariant set} $\Theta(c)$ (the level set of the conserved vector $c$) is homeomorphic to a product of hyperquadrics, one per hidden neuron,
\begin{equation*}
    \Theta(c)\cong\mathcal{Q}(c_1)\times\dots\times\mathcal{Q}(c_l),
\end{equation*}
where $c_k$ is the neuron balance (incoming minus outgoing squared norms).
For a scalar output, each negative balance disconnects the optimization space, so $\Theta(c)$ fragments into $2^{l_-}$ connected components ($l_-$ the number of neurons with $c_k<0$), potentially causing a \emph{topological obstruction}: initializing on a ``bad'' component makes the optimum unreachable, for any dataset and any loss depending on $\theta$ only through $f_\theta$.
\citet{nurisso2025} extend the analysis to general DAG architectures: collecting the balances $c:=[c_v]_{v\in\tilde{V}}$ over hidden neurons, the connectedness of the invariant set $\Theta_\Gamma(c)$ of a DAG $\Gamma$ is characterized by sign conditions on sums of conserved quantities around \emph{bottleneck} neurons (neurons with a single incoming or outgoing edge).
Also seeing CQ as space descriptors, \citet{zhao2023a} employ them as coordinates along the symmetry orbits of minima, though they do not parametrize a level set completely in general.
Exploring these orbits, they observe variations of convergence speed and sharpness.

\paragraph{Conservation Laws from Data Symmetries}
We have discussed CL stemming from architectural symmetries holding for any data, but \citet{galley2026} ask whether they can be induced by symmetries of the data.
Following the criterion of \citet{marcotte2023}, a necessary condition is that symmetrization strictly shrinks the flow field,
\begin{equation*}
    \operatorname{span}_{(x,y)}\nabla_\theta\ell^G_{(x,y)}(\theta)\subsetneq\operatorname{span}_{(x,y)}\nabla_\theta\ell_{(x,y)}(\theta),
\end{equation*}
where $\ell_{(x,y)}$ is the per-example loss and $\ell^G_{(x,y)}$ its average over the orbit of the symmetry group $G$.
This yields a negative result: for finite groups and analytic non-polynomial losses of the form $L(\hat y,y)=f(\langle\hat y,y\rangle)$ (e.g.\ the logistic loss), no additional conservation law exists.
However, on networks whose function decouples parameters from data as $f_\theta(x)=M(\theta)T(x)$ (called tensorizable networks, e.g.\ polynomial networks), the discrete data symmetry lifts to a continuous parameter-space symmetry with MSE loss, for which a CL holds approximately under gradient descent.

\paragraph{Intrinsic Dynamics of the Lifted Parameter}
Finally, conservation laws govern when the symmetry can be factored out of the dynamics altogether.
\citet{marcotte2026} consider losses that factor through a symmetry-invariant lifting $\phi$, i.e.\ $\mathcal{L}=f\circ\phi$ (e.g.\ the end-to-end matrix $\phi(\theta)=W_L\cdots W_1$ for linear networks, the path-lifting for ReLU networks), and study the dynamics of the lifted parameter $z(t):=\phi(\theta(t))$.
\emph{Intrinsic recoverability} asks that $\theta(t)$ be recoverable from the lift and the conservation laws $h:\mathbb{R}^D\to\mathbb{R}^N$, i.e.\ that there exists $\Gamma$ with $\Gamma(\phi(\theta),h(\theta))=\theta$.
It is equivalent to
\begin{equation*}
    \ker\partial\phi(\theta)\cap\ker\partial h(\theta)=\{0\},
\end{equation*}
meaning that any infinitesimal displacement of $\theta$ must change either the lifted variable or a conserved quantity.
As a consequence, linear networks with the product lifting are not intrinsically recoverable in general (except for some initializations), but ReLU networks are.

\subsubsection{Learning Dynamics}\label{sec:learning-dynamics}
% yang2025: permutation equivariant optimizers preserve neuron parameter's manifold topology for small enough LR
The previous part on CL showed a specific way continuous symmetries constrain training trajectories.
This section widens the view to the interplay between symmetries, discrete ones included, and the learning dynamics.
We cover symmetric states (\Cref{sec:hierarchy-param-space-symmetries}), which populate the landscape with critical manifolds or attractive sets, and the interactions of symmetries with regularization, overparametrization, and discretization.

\paragraph{Symmetric States Structure the Landscape}
In the context of matrix factorization, \citet{li2018} show, for a function $f:\mathbb{R}^m\to\mathbb{R}$ invariant under a group $G$, that whenever the group directions span the whole space, $\operatorname{Span}\{g\cdot x-x \mid g\in G,\ x\in\mathbb{R}^m\}=\mathbb{R}^m$, every fixed point of the group ($g\cdot x_G=x_G$ for all $g\in G$, i.e.\ a symmetric state) is a critical point of $f$.
They also observe that at any critical point $x$, the tangent space of the orbit contains only flat directions,
\begin{equation*}
    T_x(G\cdot x)\subseteq\operatorname{Null}\big(\nabla^2 f(x)\big):
\end{equation*}
symmetry directions carry no curvature.
For neural networks, \citet{brea2019} encounter \emph{permutation points}, a type of permutation-symmetric state.
They are critical points, typically saddles of index one.
Their argument applies to permutation-symmetric states in general: each of the $d_{l+1}$ pairs of outgoing weights can be changed as long as its sum remains constant, yielding $d_{l+1}$ flat Hessian directions.
These flat directions form a $d_{l+1}$-dimensional affine plateau of critical points from which, remarkably, all $d_l!$ reparametrizations of $\theta$ due to layer-$l$ neuron permutations can be reached along equal-loss paths.
Conversely, regularization can force symmetric structure on critical points.
For the linear autoencoder with encoder $W_1\in\mathbb{R}^{k\times d}$, decoder $W_2\in\mathbb{R}^{d\times k}$ and loss $\mathcal{L}(W_1,W_2)=\lVert X-W_2W_1X\rVert_F^2$, \citet{kunin2019p3icml} show that every critical point of the $L_2$-regularized loss satisfies $W_1=W_2^{\top}$. On whitened data ($XX^\top\propto I$), the swap $(W_1,W_2)\mapsto(W_2^{\top},W_1^{\top})$ is a loss symmetry, so the configuration $W_1=W_2^{\top}$ is a symmetric state.

\paragraph{Attractivity of Symmetric States}
Several works have explored the dynamical question of whether training is drawn to symmetric states.
Since these states are related to overparametrization and subnetworks, this question is related to simplicity biases (also a theme in singular learning theory, see \Cref{sec:adjacent-fields}). 
\citet{chen2024jsm} call a set $A\subseteq\Theta$ an \emph{invariant set} of SGD if $\mathbb{P}\big(\theta^{(t)}\in A \mid \theta^{(0)}\in A\big)=1$ for all $t\in\mathbb{N}$.
For losses satisfying an \emph{approximately} $Q$-\emph{symmetric} condition, with $Q$ a symmetric or orthogonal matrix, $\{\theta : Q\theta=\theta\}$ is an invariant set.
Two examples show the connection to simpler subnetworks.
For $Q$ diagonal with $\pm1$ entries and origin-passing activations, the \emph{sign invariant sets} are characterized by disconnected neurons.
For $Q$ a neuronwise permutation matrix, the \emph{permutation invariant sets} are the symmetric states $\theta_v=\theta_w$ for two hidden neurons $v,w$.
They derive a sufficient condition for invariant sets to be \emph{stochastically attractive}, based on the comparison between the loss curvature and the gradient noise strength, and observe this stochastic attractivity in VGG-16 trained on CIFAR-10.
\citet{ziyin2024} reach a similar conclusion for the mirror symmetries of \Cref{sec:unifying-works-param-space-symmetries}: the reflection surface $\{\theta : O^{\top}\theta=0\}$ of a symmetry $\mathcal{L}(\theta)=\mathcal{L}(R\theta)$, $R=I-2OO^{\top}$, is attractive either for GD with high enough weight decay or for SGD with a high enough learning rate.
\citet{ziyin2025} nuance the picture: measuring a \emph{symmetry-breaking distance}, $\lVert\theta_i-\theta_j\rVert_2^2$ in the permutation case, unveils a nontrivial dynamics that alternates between more and less symmetric states, both through space (layers) and time (iterations).

\paragraph{Scale Invariance and Weight Decay}
The scale symmetry combined with weight decay generates a dynamics of its own.
Under GF the parameter norm decreases monotonically.
This is unsustainable for GD: scale invariance implies $\nabla^2 \mathcal{L}(\alpha\theta)=\frac{1}{\alpha^2}\nabla^2 \mathcal{L}(\theta)$, so the loss is increasingly non-smooth towards the origin and the stability condition $\eta<2/\lambda_{\mathrm{max}}$, with $\eta$ the learning rate and $\lambda_{\mathrm{max}}$ the largest Hessian eigenvalue, must eventually break.
This leads to a chaotic behavior of GD \citep{li2020anips}, which ceases to be a faithful discretization of GF.
Gradient noise can, however, stabilize the dynamics.
Projecting the dynamics onto the unit sphere, \citet{li2020anips} show that the direction $\overline{W}_t=W_t/\lVert W_t\rVert$ evolves with an \emph{effective learning rate} $\eta_{\mathrm{eff}}=\eta/\lVert W_t\rVert^2$, the angular step size of directional updates.
Under assumptions on the noise, it equilibrates at a stationary value $\propto\sqrt{\eta\lambda}$ with $\lambda$ the weight decay parameter: the dynamics does not collapse.
Adopting the same sphere-projected view, \citet{lyu2023} explain why normalization and weight decay lead GD to flatter minima (under a well-defined \emph{spherical sharpness}, see \Cref{sec:measurements-param-space}).
When the parameter reaches a local minimum after a first phase of ``standard'' descent, weight decay accounts for most of the update direction and drives the effective learning rate upwards.
The parameter then keeps oscillating while drifting towards regions of lower spherical sharpness.
Weight decay also interacts with the symmetries of multiplicative parametrizations $W=AB^{\top}$, such as the key-query products of attention layers.
Although not phrased in symmetry terms, \citet{kobayashi2024anips} prove the bound
\begin{equation*}
    \Big|\, \lVert AB^{\top}\rVert_{*}-\frac{\lVert A\rVert_F^2+\lVert B\rVert_F^2}{2} \,\Big| \;\le\; \sqrt{\lVert A^{\top}A-B^{\top}B\rVert_{*}}\ \frac{\lVert A\rVert_{*}+\lVert B\rVert_{*}}{2},
\end{equation*}
where $\lVert\cdot\rVert_{*}$ is the nuclear norm and $A^{\top}A-B^{\top}B$ is the matrix-valued conserved quantity of linear networks (\Cref{sec:conservation-laws}).
Weight decay breaks this conservation law and drives the imbalance exponentially fast to zero, so (for bounded factors) the $L_2$ and nuclear norm penalties behave similarly, which explains a bias towards low rank.
Empirically, increasing the regularization strength decreases the \emph{pseudo rank}, the smallest $k$ such that $\sum_{i=1}^{k}\sigma_i \ge 0.95 \sum_{i=1}^{n}\sigma_i$ for singular values $\sigma_1\ge\dots\ge\sigma_n$.
Finally, \citet{tanaka2021} go beyond GF and account for a finite learning rate through the modified flow
\begin{equation*}
    \frac{\eta}{2}\ddot\theta+\dot\theta=-\nabla\mathcal{L}(\theta),
\end{equation*}
whose second-order term is an \emph{implicit acceleration}.
A Lagrangian formulation reveals that symmetries of $\mathcal{L}$ do not necessarily imply symmetries of the learning dynamics, and the conserved quantities of gradient flow can have a nontrivial evolution.
Accounting for discretization can be interpreted as producing an adaptive momentum effect similar to RMSProp for normalization layers.

\paragraph{Overparametrization}
Similarly to \citet{brea2019}, \citet{simsek2021} observe that for a minimal MLP of widths $(r_1^{*},\dots,r_L^{*})$ achieving zero loss, adding a single hidden neuron to every layer connects the $\prod_l r_l^{*}!$ disconnected permutation-equivalent minima.
They count two types of subspaces: $N_{\mathrm{sym}}(r,m)$, the number of \emph{symmetry-induced critical subspaces} obtained by expanding critical points of width-$r$ sub-architectures into the width-$m$ network, and $N_{\min}(r^{*},m)$, the number of affine subspaces composing the global-minima manifold, where the minimal width $r^{*}$ measures the degree of overparametrization.
Either count can dominate, depending on the regime: under mild overparametrization, the saddle subspaces $N_{\mathrm{sym}}(r,m)$ dominate, which intuitively means more strictly-positive-loss saddle volume, hence harder optimization.
\citet{sareen2026} assess the dynamical effect of overparametrizing by neuron splitting: splitting acts as a diagonal preconditioning of the Hessian, and minima measures such as the pseudo-determinant $\prod_{i:\lambda_i>0}\lambda_i$ increase under splitting (corresponding to a volume decrease), though the overall impact on optimization remains unclear.

\subsection{Prescription for Parameter-Space Symmetries}\label{sec:bug-or-feature}

Having described parameter-space symmetries and their consequences, we turn to the \emph{prescriptive} question: how can the ``symmetry bug'' be fixed, or can it even be exploited?
A direct fix is to remove symmetries altogether, from the model \citep{lim2024} or from the loss \citep{ziyin2024b} (\Cref{sec:mode-connectivity,sec:dynamical-sb} respectively).
Beyond this radical fix, the literature below works with symmetries and, as in \Cref{sec:descr-param-space-symmetries}, splits along the state axis.
The first two parts are mostly \emph{static} and concern the model locus, i.e.\ the functional symmetry $\varphi(g\cdot\theta)=\varphi(\theta)$.
\Cref{sec:quotient-repr-align} removes the redundancy from the representation of $\theta$ through quotients $\Theta/\!\sim$, canonical representatives or pairwise alignment.
One branch embeds this in Riemannian gradient descent, yet each normalization needs only the current $\theta$, not its history.
\Cref{sec:measurements-param-space} requires measurements on $\Theta$, such as mode-connectivity barriers and flatness, to factor through $\Theta/\!\sim$.
This literature was historically driven by intuition, and symmetries can be treated more \emph{implicitly} there.
The last two parts are \emph{dynamic}.
\Cref{sec:symmetry-optimization} moves to the optimization locus, the loss $\mathcal L$ and update rule $U$, e.g.\ with equivariant optimizers.
Finally, \Cref{sec:exploit-symmetry-feature} seeks to turn the bug into a feature.

\subsubsection{Quotient, Representative and Alignment}\label{sec:quotient-repr-align}
Keeping the symmetries of $\varphi$ in the model, a branch of works takes them out of its \emph{representation} by identifying $\theta$ with every $g\cdot\theta$, i.e.\ working modulo the orbit relation $\theta\sim\theta' \iff \theta'=g\cdot\theta$ for some $g\in G$.
Working with equivalence classes or with canonical elements thereof uncovers patterns that symmetries could otherwise hide and sidesteps the optimizer mismatch of \Cref{sec:symmetry-optimization}.
We report on three strategies.
\emph{Quotienting} constructs a representation of the equivalence classes themselves: orbits become points of the quotient space $\Theta/\!\sim$.
\emph{Representative extraction} selects a particular element inside each class, that is, a section of the quotient map $\Theta\to\Theta/\!\sim$ picking one point in each orbit.
\emph{Alignment} is a pairwise version of representative extraction: the chosen representatives are not absolute but depend on the relation between two parameter instances.

\paragraph{Quotiented Parameter Spaces}
\citet{meng2019} push the path perspective of \citet{neyshabur2015anips} one step further: they optimize ReLU networks directly in a space of invariant coordinates, the path values $\pi_c(w)$ of \Cref{sec:symmetry-optimization}.
Since these values are not independent, a subset of $(m-H)$ \emph{basis paths} is first extracted, with $m$ the number of weights and $H$ the number of hidden neurons: one degree of freedom removed per hidden neuron, in agreement with the expected functional dimension of \Cref{sec:parameter-identifiability}.
For batch-normalized networks, \citet{yi2022ptcuai} constructs the quotient as a Riemannian manifold.
Collecting the incoming weight vectors $w^{(1)},\dots,w^{(m)}$ of the $m$ neurons, the \emph{Positively Scale-Invariant manifold} is the quotient $\big((\mathbb{R}^{d_1}\setminus\{0\})\times\dots\times(\mathbb{R}^{d_m}\setminus\{0\})\big)/\!\sim$ by per-neuron positive rescalings.
It is equipped with the Riemannian metric
\begin{equation*}
    \langle\xi_1,\xi_2\rangle_{w}=\sum_{i=1}^{m}\frac{\langle\xi_1^{(i)},\xi_2^{(i)}\rangle}{\lVert w^{(i)}\rVert^2}.
\end{equation*}
Riemannian gradient descent on this manifold is equivariant in the sense of \Cref{eq:equivariant-optimizer}.

Quotienting can also express the loss on another domain.
In the linear autoencoder of \Cref{sec:learning-dynamics}, where every critical point of the regularized loss satisfies $W_1=W_2^{\top}\in \mathbb R^{k\times d}$, \citet{kunin2019p3icml} reinterpret the loss as a function on the Grassmannian $\operatorname{Gr}_k(\mathbb{R}^d)$ of $k$-planes in $\mathbb{R}^d$.
Fixing a point cloud $X$, the loss value at a $k$-plane is the sum of squared distances from the points of $X$ to that plane: the residual symmetry has been quotiented out.
Finally, identifiability questions have been investigated from a quotient perspective: writing $\bar\varphi:\Theta/\!\sim\,\to\mathcal{F}$ for the implementation map induced on the quotient, \citet{grillo2026} phrase identifiability up to quotiented symmetries as the fiber $\bar\varphi^{-1}(f_\theta)$ being a singleton.

\paragraph{Representative Extraction}
\citet{hecht-nielsen1990anc} constructs a cone $H\subset\mathbb{R}^q$ in parameter space, itself a union of cones $H_k=\{\theta\in\mathbb{R}^q : |\theta_i|\le|\theta_k|\ \forall i\neq k\}$ over an architecture-dependent subset of indices $k$, guaranteed to contain at least one parameter implementing each realizable function.
Weight search can be restricted to $H$, but it remains redundant, still containing a number of permuted copies of any given network.
This motivates the notion of a \emph{minimal search set}: in modern terms, a section of the quotient.

For rescaling symmetries, a possible section is a unit-norm slice.
\citet{badrinarayanan2015} constrain each neuron of a batch-normalized network to a unit-norm incoming weight vector, $\operatorname{diag}\big(W^{(l)}{W^{(l)}}^{\top}\big)=I$ for every layer $l$, and train by Riemannian SGD on the resulting \emph{oblique manifold}.
\citet{huang2017} approximate Riemannian gradient descent by a standard gradient step $\tilde{w}^{(t+1)}=w^{(t)}-\eta\nabla\mathcal{L}(w^{(t)})$ followed by norm projection $w^{(t+1)}=\tilde{w}^{(t+1)}/\lVert\tilde{w}^{(t+1)}\rVert$ back onto the manifold.
They do not change variables, contrary to e.g.\ \citep{salimans2016anips}.
Note that the projection rescales only incoming weights without other compensation.
Thanks to a subsequent normalization layer, this still preserves the function, showcasing an interplay between the scale and rescale symmetries.
The minimal-norm representative is another canonical choice.
\citet{stock2018} extract the unique representative of minimal $\ell_p$-norm by minimizing
\begin{equation*}
    \sum_{l=1}^{L}\big\lVert D_{l}W_{k}D_{l-1}^{-1}\big\rVert_p^p
\end{equation*}
over positive diagonal matrices $D_1,\dots,D_{L-1}$ ($D_0=D_L=I$) with a coordinate descent scheme.
Applying this \emph{Equi-normalization} in-between SGD steps empirically improves performance.
\citet{saul2023} reach the minimal $\ell_{p,q}$-norm representative through per-unit closed-form multiplicative updates: the incoming weights of every hidden neuron $i$ are simultaneously multiplied (and its outgoing weights divided) by
\begin{equation*}
    a_i=\Bigg[\frac{\sum_{j}\rho_j(W)\,\pi_{ji}(W)}{\rho_i(W)}\Bigg]^{\frac{1}{4\max(p,q)}} ,
    \quad
    \pi_{ji}(W)=\frac{|W_{ji}|^p}{\sum_{k}|W_{jk}|^p},
\end{equation*}
where $\rho_i$ measures the input norm of neuron $i$.
If neuron $i$ has a smaller input norm than the weighted average of the input norms of the next-layer neurons, then $a_i>1$ and the update drives the network towards balance.

\paragraph{Alignment}
Alignment extracts paired representatives, which makes it a tool of choice for comparing or merging trained networks and a pillar of the mode connectivity literature (\Cref{sec:mode-connectivity}).
Working in activation space, \citet{lenc2014a} aligned the representations $f_\theta^{(l)}$ and $f_{\theta'}^{(l)}$ of two networks with a learned linear map $E$, $E\,f_\theta^{(l)}\approx f_{\theta'}^{(l)}$, and showed that stitching representations from one trained network to another was possible.
\citet{godfrey2022} sharpen the connection to parameter-space symmetries: for ReLU networks, restricting the stitching map from arbitrary linear layers to the intertwiner group $G_{\mathrm{ReLU}}$ of \Cref{sec:unifying-works-param-space-symmetries} suffices to preserve good performance, suggesting that parameter symmetries account for much of the misalignment between independently trained representations.
They propose feature-space dissimilarity measures invariant to $G_{\mathrm{ReLU}}$ (and to global translation): $G_{\mathrm{ReLU}}$-\emph{Procrustes} compares feature matrices $X,Y\in\mathbb{R}^{n\times d}$ through
\begin{equation*}
    \mu(X,Y):=1-\frac{1}{2\sqrt{d}}\,\min_{P}\big\lVert XD_X^{-1}-YD_Y^{-1}P\big\rVert_F,
\end{equation*}
where $P$ ranges over permutation matrices and $D_X=\operatorname{diag}\big(\lVert X_{:,1}\rVert,\dots,\lVert X_{:,d}\rVert\big)$ normalizes features.
\citet{zhang2025p4icml} align the double rotation symmetry of self-attention in transformer models.
The restriction from the full $GL(d_k,\mathbb{R})$ symmetry to rotations trades generality for a closed-form solution, but their matching algorithm separately includes a rescaling alignment of self-attention matrices and a permutation alignment in MLPs.
It improves model fusion.
Finally, for ReLU classifiers, \citet{pittorino2022jsm} combine representative extraction with alignment in two steps.
(1) Each hidden neuron $v$ is rescaled to the oblique manifold, and the output layer is globally rescaled (they work with the error landscape rather than the loss landscape).
The normalized parameter space becomes a product of spheres, including one sphere $S^{d_{l-1}}$ per hidden neuron of layer $l$.
(2) When two networks $\theta$ and $\eta$ are compared, layerwise permutations maximizing the cosine agreement $\sum_{v=1}^{H}\langle\theta_v^{+},\eta_v^{+}\rangle$ between normalized incoming weight vectors over each sphere are found by linear assignment.
On the resulting geometry, geodesics replace linear interpolation paths, and distances, barriers and flatness are closer to a functional meaning, as we see next in \Cref{sec:measurements-param-space}.

\subsubsection{Measurements in Parameter Space}\label{sec:measurements-param-space}
Evaluations of $f_\theta$, such as generalization measures, should not depend on the specific $\theta$ realizing it~\citep{dinh2017}.
Concretely, a measurement $m:\Theta^k\to\mathbb R$ taking $k$ parameters as input should factor through a well-defined induced measurement on the quotient, $(\Theta/\!\sim)^k\to\mathbb R$, or through the realization map $\varphi:\Theta\to\mathcal F$ as a measurement of functions, $\mathcal F^k\to\mathbb R$.
Parameter-space symmetries are thus an invariance desideratum for any measurement built from $\theta$, similarly to the equivariant metanetworks of \Cref{sec:adjacent-fields}.
We focus on two such measurements: barriers in (linear) mode connectivity, and flatness of the loss landscape.

\begin{Remark}
    Other measurements of $\theta$ face the same requirement.
    Classical generalization bounds scale with the dimension of $\theta$, but \emph{norm-based complexity measures} strive for a dimensionless notion of the magnitude of $\theta$, which should also be consistent with symmetries.
    \citet{liang2019pticais} propose the \emph{Fisher-Rao norm} (technically not a norm):
    \begin{equation*}
        \lVert\theta\rVert_{\mathrm{fr}}^2:=\langle\theta,I(\theta)\theta\rangle,\qquad I(\theta)=\mathbb E\big[\nabla_\theta L(f_\theta(X),Y)\otimes\nabla_\theta L(f_\theta(X),Y)\big],
    \end{equation*}
    which is invariant under the reparametrizations at stake, contrary to others such as the $\ell_2$ norm.
    One can also measure a property of the minima manifold $\varphi^{-1}(f_\theta)\subset\Theta$ as a geometric object.
    For instance, \citet{zhao2024} look at the curvature along orbits, though the measure is not invariant under reparametrization.
\end{Remark}

\paragraph{(Linear) Mode Connectivity}\label{sec:mode-connectivity}

\emph{Mode connectivity} (MC) \citep{garipov2018,draxler2018p3icml} asks whether two minima are connected by a low-loss path, the obstruction being quantified by the \emph{path height}
\begin{equation}\label{eq:path-height}
    \mathcal{B}:\Theta\times\Theta\to\mathbb{R},\qquad(\theta_0,\theta_1)\longmapsto\inf_{\gamma}\ \max_{t\in[0,1]}\mathcal{L}(\gamma(t)),
\end{equation}
over the continuous curves $\gamma:[0,1]\to\Theta$ with $\gamma(0)=\theta_0$ and $\gamma(1)=\theta_1$.
If the symmetry group $G$ of the loss is a Lie group, the path height is invariant under the action of the path component \citep{baker2003} of the identity $G^{\circ}$:
$$
    g_1, g_2\in G^{\circ} \Rightarrow \mathcal{B}(g_1\theta_1, g_2\theta_2)=\mathcal{B}(\theta_1, \theta_2)
$$
since any path between $\theta_1$ and $\theta_2$ can be extended at no cost by actions from $G^{\circ}$.
The literature therefore concentrates on disconnected symmetry groups, in particular permutations.

\citet{freeman2017} motivate the study of sublevel sets $\Theta_{\le\lambda}=\{\theta\in\Theta:\mathcal{L}(\theta)\le\lambda\}$ since their connectedness guarantees the absence of spurious local minima.
They show that overparametrization asymptotically prevents obstructions due to permutation symmetry.
For a shallow ReLU network of width $m$ and any $\theta,\theta'$ with $\max\big(\mathcal{L}(\theta),\mathcal{L}(\theta')\big)\le\lambda$, they construct a continuous path $\gamma$ from $\theta$ to $\theta'$ such that
\begin{equation*}
    \max_{t\in[0,1]}\mathcal{L}(\gamma(t))\le\max(\lambda,\epsilon),\qquad \epsilon=\mathcal{O}\big(m^{-1/d}\big),
\end{equation*}
with $d$ the input dimension.
Remarkably, the rescaling symmetry enters the proof: a continuous symmetry serves to remove the obstacle raised by a discrete one.
Rather than diluting the permutation asymptotically, other works aim at resolving it directly.
\citet{tatro2020} align the hidden neurons between identical architectures, choosing the permutation that minimizes, on the training set, the activation-based cost
\begin{equation*}
    c(x,y):=1-\operatorname{corr}(x,y)
\end{equation*}
between neurons $x$ and $y$ of the two networks.
After alignment, the path height is reduced.
\citet{brea2019} sample paths between a minimum and a permuted version under the constraint of passing through a permutation-symmetric state, which is an index-$1$ saddle.
They notice that, at this saddle, an extended loss plateau gives access to all possible symmetric states, and observe that one redundant pair of neurons suffices to make the permuted minima connected \citep{brea2019,simsek2021}.
\citet{zhao2025a} study the minima of a full-rank linear network, $\mathcal{L}^{-1}(0)$ for $\mathcal{L}(\theta)=\lVert Y-W_L\cdots W_1X\rVert_F^2$ with $X,Y\in GL(h,\mathbb{R})$ and $W_l\in\mathbb{R}^{h\times h}$.
The internal $GL(h,\mathbb{R})^{L-1}$ symmetry acts transitively on that set, and since continuous maps preserve connectedness, the topology of the group transfers to the topology of the minima: $\mathcal{L}^{-1}(0)\cong GL(h,\mathbb{R})^{L-1}$, which thus contains exactly $2^{L-1}$ components, one per sign pattern of the determinants.
This sign pattern can be chosen by permuting neurons, so MC holds up to permutation.

\emph{Linear mode connectivity} (LMC) replaces general paths with straight segments.
A \emph{barrier} can then be measured against the linear interpolation of the endpoint losses,
\begin{equation}\label{eq:linear-barrier}
    B(\theta_0,\theta_1)=\sup_{\alpha\in[0,1]}\Big[\mathcal{L}\big(\alpha\theta_0+(1-\alpha)\theta_1\big)-\big(\alpha\mathcal{L}(\theta_0)+(1-\alpha)\mathcal{L}(\theta_1)\big)\Big],
\end{equation}
and the invariance to continuous symmetries of \Cref{eq:path-height} breaks down.
To begin with, LMC can occur without alignment when the two endpoints are produced from a common starting point.
\citet{frankle2020p3icml} split training after the $k$-th iteration into two copies with distinct SGD noise.
The two outcomes have a non-negligible linear barrier only for small $k$: a short initial phase of shared training is typically sufficient for LMC.
\citet{lim2024} obtain the same effect architecturally by removing the symmetries from the model.
They build \emph{asymmetric} networks, either through an untrainability mask $W\mapsto M\odot W+(1-M)\odot W_{\mathrm{fixed}}$ whose rows are nonzero and distinct (the networks remaining universal approximators) or through a non-elementwise activation.
Independently trained networks are then forced to be compatible and interpolate linearly with low barrier, with smoother loss landscapes.
LMC up to permutation encompasses more natural settings.
\citet{entezari2022} conjecture that, once permutations are accounted for, most SGD solutions have LMC.
As empirical support, they show that, when sampling $\theta_0,\theta_1$ with SGD and drawing a random permutation $\sigma$, the barrier $B(\theta_0,\theta_1)$ between independent solutions and the barrier $B(\theta_0,\sigma\cdot\theta_0)$ within a single orbit follow similar distributions, as if independent runs differed only by a permutation.
On the theoretical side, wide-enough shallow ReLU networks satisfy the conjecture under standard initialization conditions.
\citet{ainsworth2023} make the permutation search more practical, aligning the two networks by coordinate descent over layerwise assignment problems.
They exhibit a zero-barrier linear interpolation between independently trained ResNets, and applications to merging models trained on disjoint data splits in the spirit of federated learning.
Consistently with the invariance requirement of this section, \citet{pittorino2022jsm} account for two symmetries and evaluate barriers only after rescaling-representative extraction and permutation alignment, using geodesic paths in place of linear ones in their product-of-spheres parameter space.
However, the conjecture is not universally true: \citet{zhao2025a} use the rescaling symmetry to produce a counterexample.
For homogeneous activations, comparing the segment joining $(W_1,\dots,\alpha W_{L-1},\alpha^{-1}W_L)$ to $(W_1,\dots,\alpha^{-1}W_{L-1},\alpha W_L)$ as $\alpha\to0$ yields two functionally equivalent minima whose linear barrier is arbitrarily large, and (under assumptions) no permutation fixes it.
Finally, \citet{tran2025anips} cover Mixture-of-Experts, whose symmetry group is $G(n)=\mathbb{R}^{d+1}\times S_n$, the continuous factor coming from the softmax over gating scores (\Cref{sec:descr-param-space-symmetries}).
They show that, perhaps surprisingly, $B$ is invariant under this $\mathbb{R}^{d+1}$ factor: the mechanism of \Cref{eq:path-height} survives the restriction to linear paths, so that permutations alone govern LMC for MoEs.
Since the experts carry symmetries of their own, they propose a hierarchical alignment: experts are first matched (as blocks of parameters), e.g.\ by their gating scores, then the neurons inside each matched pair of experts are matched in turn.
The procedure reduces the linear barrier for both vision transformers and language models.

\begin{Remark}
    Alignment is not the only difficulty.
    \citet{jordan2023} observe a \emph{variance collapse} along the interpolating path: two aligned units are only imperfectly correlated, so their midpoint has lower variance, compounding with depth until deep activations are nearly constant.
    Their REPAIR algorithm mitigates the problem by using batch normalization to rescale the preactivations.
\end{Remark}

\paragraph{Flatness}

Intuitively, a minimum $\theta^\star$ is \emph{flat} when the loss stays close to $\mathcal{L}(\theta^\star)$ on a neighborhood of $\theta^\star$.
\citet{keskar2017} reported that large-batch SGD converges to sharper minima than small-batch SGD and that this sharpness tracks the generalization gap: empirically, flatness correlates with better test performance.
However, \citet{dinh2017} soon objected that the usual definitions of flatness are not well posed.
For a one-hidden-layer ReLU network, they lower bound the Hessian at a rescaled parameter $T_\alpha\theta:=(\alpha\theta_1,\alpha^{-1}\theta_2)$ by
\begin{equation}\label{eq:dinh-hessian-bound}
    \big\lVert(\nabla^2\mathcal{L})(T_\alpha\theta)\big\rVert_F\ \ge\ \alpha^{-2}\gamma,
\end{equation}
where $\gamma>0$ is a positive diagonal entry of $(\nabla^2\mathcal{L})(\theta)$, demonstrating that a minimum can be made arbitrarily sharp.

Subsequent measures address this gap \citep{tsuzuku2019,yi2019,petzka2021anips}.
For example, \citet{tsuzuku2019} propose the \emph{normalized sharpness}
\begin{equation*}
    \sum_{l}\sqrt{\lVert W^{(l)}\rVert_F^2\;H^{(l)}},
\end{equation*}
where $H^{(l)}$ is the sum of diagonal entries of the Hessian of the loss.
Experimentally, they find that their measure correlates better with the generalization gap than unnormalized measures like the trace of the Hessian.
Instead of redefining, \citet{pittorino2022jsm} canonicalize: flatness is evaluated after the rescaling-representative extraction of \Cref{sec:quotient-repr-align}.
\citet{stutz2021} adapt this principle for adversarial training, where the quantity of interest is the \emph{robust generalization} gap, measured on adversarial test examples.
They propose an invariant version of the worst-case increase of the robust loss in a neighborhood of the parameter.
For scale-invariant losses, \citet{lyu2023} propose the \emph{spherical sharpness} $\lambda_1\big(\nabla^2\mathcal{L}(\theta/\lVert\theta\rVert_2)\big)$, where $\lambda_1$ extracts the largest eigenvalue.
For \citet{kristiadi2023anips} (see \Cref{sec:symmetry-optimization}), flatness appears non-invariant only because the Riemannian metric is implicitly taken to be the identity, $M=I$.
They build an invariant Hessian $E:=M^{-1}H$ from the Hessian $H$ by seeing it as a linear operator on the tangent space $T_\theta\Theta$.
Because $M$ and $H$ are both bilinear forms, they transform in the same way, so that $\det E$, for example, is invariant.
\citet{petzka2021anips} look closer at the correlation reported by \citet{keskar2017}.
They study models of the form $f(W_L\,f_\theta^{(L-1)}(x))$, with $f_\theta^{(L-1)}$ the penultimate feature map and $W_L$ the weights of the final linear layer, and find that, for the flatness-generalization connection to hold, feature smoothness is important ($f_\theta^{(L-1)}$ must have locally constant labels, which must be stable to small perturbations in feature space).
They also propose a flatness measure invariant to rescalings, defined as a Hessian trace weighted by the weights of the model.

Sharpness-Aware Minimization (SAM) \citep{foret2021} treats classic flatness as a part of the objective, replacing the training loss by its worst case over a ball,
\begin{equation*}
    \mathcal{L}^{\mathrm{SAM}}(\theta):=\max_{\lVert\delta\rVert_p\le r}\mathcal{L}(\theta+\delta),
\end{equation*}
with $r$ a radius hyperparameter and $p\in[1,\infty]$.
\citet{kwon2021p3icml} correct it by proposing the \emph{adaptive SAM} objective
\begin{equation*}
    \max_{\lVert T_\theta^{-1}\delta\rVert_p\le r}\mathcal{L}(\theta+\delta)-\mathcal{L}(\theta),
\end{equation*}
where $T_\theta^{-1}$ is a normalization operator responsible for canceling the effect of rescaling on sharpness.
For ReLU networks, $T_\theta=\operatorname{diag}(|\theta|)$ (entrywise absolute values) is a valid choice.
Adaptive sharpness is shown to correlate better with generalization.

\subsubsection{Optimization}\label{sec:symmetry-optimization}
Moving from \emph{static} measurements to \emph{dynamic} training, we turn to symmetries of the loss $\mathcal L$ and of the update rule $U$ (optimization locus): the obstacles they create and their remedies.

\paragraph{Symmetry-Induced Obstacles}
A first type of obstacle is analytical.
In asymmetric matrix factorization, the $GL(r,\mathbb{R})$ symmetry makes the loss $(A,B)\mapsto\lVert AB^{\top}-M^{*}\rVert_F^2$ non-smooth: along $(\alpha A,\alpha^{-1}B)$ the loss stays constant while the Lipschitz constant of its gradient blows up.
As seen in \Cref{sec:conservation-laws}, \citet{du2018anips} contain the pathology with decreasing stepsizes under which the approximately conserved balancedness $\lVert A^{\top}A-B^{\top}B\rVert_F$ remains bounded.
Iterates are thereby confined to a region where smoothness holds, which eases the analysis.
A dynamical obstacle arises from weakly broken symmetry in the dynamic word embeddings of \citet{bamler2018p3icml}.
A coupling regularizer $\psi(Z)$ breaks the per-time-step rotations $(R_1,\dots,R_T)$ down to a single global rotation $R_1=\dots=R_T$.
However, the broken orbit directions are nearly flat (\emph{Goldstone modes} in spontaneous symmetry breaking in physics), and GD undergoes two phases: it first converges to an almost symmetric configuration optimizing the unregularized loss, then shifts very slowly along the approximate orbit to decrease $\psi$.
Their \emph{Goldstone GD} decouples the two dynamics by alternating descent on the total loss with a minimization of $\psi(Z)$ restricted to the symmetry subspace, that is, over the rotations $(R_1,\dots,R_T)$.

\paragraph{Optimizers Carry Their Own Symmetries}
Equivariance of the update map is defined as for any other map \citep{neyshabur2015anips,yi2022ptcuai,yen2025iclr,zhao2025,yang2025}: for $G$ acting on $\Theta$, the update $U:\Theta\to\Theta$ is equivariant to $G$ when
\begin{equation}\label{eq:equivariant-optimizer}
    U(g\cdot\theta)=g\cdot U(\theta)\qquad\forall g\in G,\ \theta\in\Theta.
\end{equation}
Optimizers have different invariance and equivariance properties.
GD is equivariant to orthogonal reparametrizations $\theta\mapsto Q\theta$, $Q\in O(d)$, and Newton's method to affine reparametrizations \citep{kristiadi2023anips}.
Natural gradient descent is infinitesimally invariant under any smooth invertible reparametrization \citep{amari1998nc,liang2019pticais,kristiadi2023anips,zhao2025}.
Adam \citep{kingma2015} is invariant to positive diagonal scalings of the gradients.
This last property is \emph{scale-freeness}: the iterates $\{\theta_t\}_{t\geq0}$ are unchanged when every gradient is rescaled as $\nabla\mathcal{L}(\theta_t)\gets D\nabla\mathcal{L}(\theta_t)$ for a positive diagonal matrix $D$, and it turns out to separate variants of Adam \citep{zhuang2022tmlr}.
AdamW postpones the weight-decay term $\lambda\theta_{t-1}$, adding it after the adaptive step instead of folding it into the gradient as Adam-$\ell_2$ does, and thereby does not treat the $L_2$ penalty like the task loss.
\citet{zhuang2022tmlr} show that this preserves scale-freeness while Adam-$\ell_2$ loses it, and argue that scale-freeness carries optimization benefits (improving conditioning).

\paragraph{Equivariant Optimizers}
A mismatch between the symmetries of the model and those of the optimizer means that functionally identical parameters need not train identically.
When the equivariance requirement of \Cref{eq:equivariant-optimizer} is fulfilled, however, $\theta_0\sim\tilde{\theta}_0$ implies $\theta_t\sim\tilde{\theta}_t$ for all $t$, and we may speak of \emph{the} functional trajectory unambiguously.
One strategy to avoid the mismatch is thus to employ an optimizer equivariant to the symmetry group of the model (\Cref{sec:quotient-repr-align} presents other approaches).
This symmetry-guided view would also fit in \Cref{sec:design-principle}.

Observing that SGD is not equivariant to the rescaling symmetry of ReLU networks, \citet{neyshabur2015anips} repair it through paths.
The \emph{path vector} $\pi(w)$ collects the products of weights along all input--output paths: $\pi_c(w)=\prod_{k}w_{e_k}$ for a path $c=(e_1,\dots,e_L)$ of edges.
Since every path crosses each neuron through one incoming and one outgoing edge, $\pi$ is rescaling-invariant, and so are the proposed $\ell_p$-path regularizer $\lVert\pi(w)\rVert_p$ and \emph{Path-SGD}, a variant of SGD with the coordinate-wise update
\begin{equation*}
    w_e^{(t+1)}\gets w_e^{(t)}-\frac{\eta}{\gamma_p(w^{(t)},e)}\,\frac{\partial\mathcal{L}}{\partial w_e}(w^{(t)}),
\end{equation*}
where $\gamma_p(w,e)$ is the squared $\ell_p$-norm of the path vector restricted to paths containing edge $e$, with $w_e$ itself removed from each product.
\citet{kristiadi2023anips} distinguish \emph{weight-space symmetries}, which move $\theta$ while the parameter space itself does not transform, from \emph{reparametrizations} $\phi:\Theta\to\Psi$ between distinct parameter spaces regarded as Riemannian manifolds with explicit metrics.
Under a reparametrization, transformation rules apply not only to the parameter but to every object attached to it, tangent vectors and gradients in particular.
The apparent non-invariance of GD then stems from the metric $M$ being left implicitly Euclidean.
Making it explicit in the update, $\dot{\theta}=-M(\theta)^{-1}\nabla\mathcal{L}(\theta)$, makes GD invariant.
\citet{yen2025iclr} make a similar contribution for LoRA, whose update $\Delta W=AB^{\top}$ has $GL(r,\mathbb{R})$ symmetry with $r$ the low-rank dimension, and show that most practical optimizers (Adam, RMSProp, Shampoo, \dots) are not equivariant to it.
They restore equivariance with a preconditioning matrix $P$, which they show cannot be restricted to diagonal matrices.
For efficiency, it is applied as a right preconditioner in the low $r$-dimensional row space (i.e.\ $\nabla_A\mathcal{L}\,P$ instead of $P\,\nabla_A\mathcal{L}$).

\subsubsection{Exploit the Symmetry Feature}\label{sec:exploit-symmetry-feature}
Some works treat symmetries as a resource rather than a nuisance, finding symmetry levers, for instance to accelerate training or perform pruning.

\paragraph{Overparametrization}
In linear networks $x\mapsto W_L\cdots W_1x$, \citet{arora2018p3icml} study the dynamics of the end-to-end matrix $W_e:=W_L\cdots W_1$ under gradient descent and find that overparametrization by depth acts as a preconditioner combining adaptive learning rate and momentum-like effects.
It can improve training speed, and they prove that the mechanism cannot be replicated by any regularization term.
\citet{jacot2022} exploit overparametrization along a different axis, width, as a formalizing tool.
The \emph{saddle-to-saddle conjecture} states that, under small initialization, gradient flow in a deep linear network visits saddles of increasing rank before reaching a low-rank global minimum.
Conceptually, this can be formalized using an overparametrization symmetry: the full-width network dynamics is approximated as a sequence of sub-architectural dynamics starting with a width $1$ network and progressively adding neurons.

\paragraph{Teleportation}
A line of works seeks to interleave gradient steps and movements on the orbit.
\citet{armenta2021b} instantiate this within a quiver representation formalism \citep{armenta2021a}, primarily for the rescale symmetry.
Although the function is preserved, feature maps, flatness, and gradients are not: the gradient transforms as
\begin{equation*}
    \frac{\partial \mathcal{L}}{\partial v^{(l)}_{ij}}=\frac{\tau^{(l-1)}_i}{\tau^{(l)}_j}\,\frac{\partial \mathcal{L}}{\partial w^{(l)}_{ij}},
\end{equation*}
where $v^{(l)}_{ij}=\tau^{(l)}_jw^{(l)}_{ij}/\tau^{(l-1)}_i$ are the rescaled weights and $\tau^{(l)}_i$ is the scalar value of the rescaling of neuron $i$ in layer $l$, an exact backward-pass counterpart of the forward rescaling.
Empirically, teleporting at initialization or mid-training can speed up SGD.
\citet{zhao2023} extend the idea: for $G$ acting on $w$, teleportation picks $g\leftarrow\arg\max_{g\in G}\lVert\nabla\mathcal{L}(g\cdot w_t)\rVert^2$ and updates $w_t\leftarrow g\cdot w_t$ before the usual gradient step.
For continuous groups, the optimization problem can be solved by gradient ascent on $g$.
\citet{zhao2023} give theoretical motivation.
Writing $J=\partial(g\cdot w)/\partial w$ for the Jacobian of the transformation and $\eta$ for a (possibly matrix-valued) learning rate playing the role of a metric, teleportation accelerates the loss decay if
\begin{equation*}
    \big\lVert[J^{-1}]^{\top}\nabla\mathcal{L}(w)\big\rVert_{\eta}^2>\big\lVert\nabla\mathcal{L}(w)\big\rVert_{\eta}^2,
\end{equation*}
meaning that orthogonal transformations are useless.
Furthermore, the optimal target relates to Newton's method: at the point of a level set with locally maximal gradient norm, the gradient direction coincides with the Newton direction.
In practice, however, the maximizing $g$ may lie at infinity or be costly to approximate.
\citet{zhao2024} replace the inner loop by a meta-learned optimizer for $g$, avoiding the per-step search cost.

\paragraph{Pruning}
Conservation laws (\Cref{sec:conservation-laws}) have proven useful to design pruning strategies.
\citet{tanaka2020anips} design a pruning-at-initialization scheme that avoids layer collapse.
Their algorithm computes per-edge scores $\frac{\partial \mathcal{L}}{\partial \theta}\odot \theta$ which, because of the conservation law, balance between layers and thus assign higher scores to sparser layers, thereby preventing total pruning of a layer.
The invariant set $\Theta_\Gamma(c)=\{\theta:\tilde{B}\theta^2=c\}$ of \citet{nurisso2025}, where $\tilde B$ encodes the conservation constraint, contains singularities that correspond exactly to configurations where zero-weighted edges disconnect a subset of neurons from both input and output, i.e.\ to prunable subnetworks.
Regularizing the rank of the Jacobian $J_\Gamma(\theta)=2\tilde{B}\operatorname{diag}(\theta)$ yields a differentiable pruning approach that empirically produces whole-neuron sparsification rather than the edge-level sparsity of an $\ell_1$ penalty.

\subsection{Adjacent fields}\label{sec:adjacent-fields}
We close this section by mentioning adjacent fields that directly engage with parameter-space symmetries.

\paragraph{Weight-space learning and Neural Functional Networks}
\emph{Neural Functional Networks} \citep{zhou2023} (also called \emph{metanetworks}) process the parameter vector $\theta$ (or the gradient) of another network or an implicit neural representation encoding a signal \citep[INR,][]{sitzmann2020anips}.
Since its input lives in $\Theta$ and two parameters related by a symmetry carry the same information, the parameter-space symmetry group is exactly the group a metanetwork should be invariant or equivariant to.
Weight-space learning is in this sense geometric deep learning applied one level up.

For an $L$-layer MLP, the weight space splits into
\begin{equation*}
    \Theta=\mathcal{W}\oplus\mathcal{B}=\Big(\bigoplus_{l=1}^{L}\mathcal{W}_l\Big)\oplus\Big(\bigoplus_{l=1}^{L}\mathcal{B}_l\Big)=\Big(\bigoplus_{l=1}^{L}\mathbb{R}^{d_l\times d_{l-1}}\Big)\oplus\Big(\bigoplus_{l=1}^{L}\mathbb{R}^{d_l}\Big).
\end{equation*}
\citet{navon2023p4icml} analyze the constraints of affine $G$-equivariant maps on this decomposition and propose the \emph{DWS layer} $\Lambda:\Theta\to\Theta$ structured as
\begin{equation*}
    \Lambda=\begin{pmatrix}\Lambda_{\mathcal{W}\mathcal{W}} & \Lambda_{\mathcal{B}\mathcal{W}}\\[2pt] \Lambda_{\mathcal{W}\mathcal{B}} & \Lambda_{\mathcal{B}\mathcal{B}}\end{pmatrix}:\mathcal{W}\oplus\mathcal{B}\longrightarrow\mathcal{W}\oplus\mathcal{B},
\end{equation*}
where each block ($\Lambda_{\mathcal{W}\mathcal{W}}$, etc.) splits into maps $\mathcal{W}_i\to\mathcal{W}_j$, and every such elementary block follows the set equivariance rule \citep{zaheer2018}.
Stacking DWS layers with pointwise nonlinearities yields a network expressive enough to approximate the forward pass of the original model and (under regularity conditions) any functional $\mathcal{F}\to \mathbb R$ of the implemented function $f_\theta$.
Concurrently, \citet{zhou2023} propose a similar construction.
Applications include predicting the class of an INR and domain adaptation.
Graph metanetworks \citep{lim2024iclr,kofinas2024iclr} treat the problem as a GNN task and thus benefit from permutation equivariance by design. 
\citet{kalogeropoulos2024} extend them to the rescaling and sign-flip symmetries.
The technical difficulty is that the classical GNN message function $\operatorname{MSG}$ must now receive arguments that scale differently.
Writing $h_V,h_E$ for the vertex and edge features of one message-passing round and $q_i:=q_l(i)$, $q_j:=q_{l-1}(j)$ for the rescalings of a neuron $i$ in layer $l$ and of a neighbor $j$ in layer $l-1$, the constraint reads:
\begin{equation*}
    q_i\operatorname{MSG}\big(h_V(i),h_V(j),h_E(i,j)\big)=\operatorname{MSG}\big(q_ih_V(i),\;q_jh_V(j),\;q_iq_j^{-1}h_E(i,j)\big).
\end{equation*}
For transformers, \citet{tran2025iclr} identify the symmetry group of multi-head attention with $h$ heads as
\begin{equation*}
    G:=\big(GL(d_k,\mathbb{R})^{h}\times GL(d_v,\mathbb{R})^{h}\big)\rtimes S_h,
\end{equation*}
with $S_h$ permuting the heads together with their per-head factors, under the assumption $d_k,d_v\le d$ that keys and values project to a lower dimension.
They then propose Transformer-NFN, equivariant under $G\times S_d\times S_{d_A}$, which includes permutation symmetries of the transformer outside attention ($d_A$ is the width of the downstream shallow MLP) but leaves aside some symmetries.

Note that the converse move is also possible: by removing symmetries in the model, the method of \citet{lim2024} finds improved weight-space learnability.
Finally, \emph{hypernetworks} \citep{ha2017} are a conceptual opposite, mapping a code $z$ to the parameters $\theta=\Psi(z)$ of a target network.
Parameter-space symmetries then appear in the codomain of the hypernetwork forward map $\Psi:\mathcal Z\to \Theta$, and $\Psi$ selects a point arbitrarily inside the fiber $\varphi^{-1}(f_\theta)$.

\paragraph{Algebro-Geometric Approaches}
\emph{Singular learning theory} \citep[SLT,][]{watanabe2009} is a Bayesian framework leveraging algebraic geometry to extend statistical learning theory to \emph{singular} (i.e.\ non-identifiable) models, as opposed to the \emph{regular} ones of classical statistics.
Tools such as the \emph{Real Log Canonical Threshold} (RLCT) or learning coefficient $\lambda$ are used to quantify the effective dimension of a model.
Insights on generalization are a main goal of SLT.
\citet{lau2025p2icais} build on the RLCT to propose the \emph{local learning coefficient} (LLC), a complexity measure for trained networks.
For $\theta^\star$ a minimum of $\mathcal{L}$ and $B(\theta^\star)$ a neighborhood small enough that $\mathcal{L}(\theta^\star)\le\mathcal{L}(\theta)$ for all $\theta\in B(\theta^\star)$,
\begin{equation*}
    V(\epsilon):=\operatorname{Vol}\big\{\theta\in B(\theta^\star)\ :\ \mathcal{L}(\theta)-\mathcal{L}(\theta^\star)<\epsilon\big\}
\end{equation*}
is the volume of the sublevel set at tolerance $\epsilon$.
It scales (in simple cases) like $V(\epsilon)\propto\,\epsilon^{\lambda(\theta^\star)}$ as $\epsilon\to0$, and the rational exponent $\lambda(\theta^\star)$ is the LLC.
It measures volume growth: the more degenerate the minimum, the more low-loss volume around it and the smaller the LLC.
The LLC is invariant under any local diffeomorphism, such as common continuous parameter-space symmetries.
The authors further derive a tractable estimator.

\emph{Neuroalgebraic geometry} \citep{marchetti2025a} is an emerging field studying the function spaces, or \emph{neuromanifolds} $\mathcal{M}=\varphi(\Theta)$, of \emph{algebraic} models such as polynomial neural networks (PNNs).
The program laid out by \citet{marchetti2025a} translates deep learning questions into algebro-geometric terms: dimension for expressivity, and singularities, which coincide with subnetworks and attract empirical risk minimizers, for implicit bias.
For example, for homogeneous PNNs, the dimension of the closure of the neuromanifold is maximal (equal to the number of parameters) if and only if the model is finitely identifiable \citep{usevich2025}.
While the polynomial setting is restrictive, it still approximates a broad class of functions, and represents an idealized yet analytically tractable middle ground, similarly to infinite-width regimes.



\section{Symmetry as an Unknown: Symmetry Discovery}\label{sec:sym-disco}
The previous sections typically assumed \emph{known} symmetries.
Here the symmetry is a priori \emph{unknown}: it is to be recovered, or its validity assessed, rather than imposed (\Cref{sec:design-principle}) or understood (\Cref{sec:parameter-space-symmetries}).
Quantification (\Cref{sec:sym-quantification}) actually lies at the boundary between \emph{known} and \emph{unknown}: one starts with a candidate symmetry and measures how much it holds.\newline
Symmetry discovery seeks to recover both a group and how it acts, that is, the triplet $(G,\rho_X,\rho_Y)$ for a map $f:\mathcal{X}\to \mathcal{Y}$.
The two subsections split along the object axis (\Cref{fig:sym-loci}), with $f$ setting the locus of the symmetry.
\Cref{sec:sym-disc-from-data} is mostly at the \emph{problem} level: $f$ is the ground truth $f^*$ or the data distribution $p$, observed only through a finite dataset.
The motivation is practical: supplying the a priori knowledge that equivariant design requires and avoiding the symmetry bottleneck (\Cref{sec:edl-tradeoffs}), which leads to performance gains, or simply gaining insight into a dataset.
\Cref{sec:sym-extract-model} is at the \emph{parameter} level: $f=f_\theta$ for a trained, hence fixed, $\theta$.
Unlike in \Cref{sec:parameter-space-symmetries}, $G$ acts on the pair $(\mathcal X,\mathcal Y)$, not on $\Theta$.
\pnote{Its motivation is more theoretical: the symmetries a network implicitly learns reveal its behavior and inductive biases.}
As a whole, the section is mainly \emph{descriptive}.


\subsection{Symmetry Discovery at the problem level}\label{sec:sym-disc-from-data}
% hu2025nn seems similar to krippendorf2020
% add: ko2024anips
% non-linear SD: yang2024p4icml, shaw2024
% add ouderaa2023
% seems that some GAN objective assume distribution invariance
% maybe cohen2014
% wang2024p4icml discovery by seeing how soft equivariance method reacts to the data 
In this part, the symmetry is carried by the ground truth $f^*$ or, less frequently, the data distribution $p$ or a single data point, and must be identified from finitely many samples.
This can be contrasted on the explicitness axis (\Cref{tab:taxonomy}) with automated DA (\Cref{par:auto-da}), where an augmentation policy only produces training examples, leaving the symmetry \emph{implicit}.
Discovery instead makes the symmetry \emph{explicit}, as an interpretable, reusable object for a downstream inductive bias, augmentation or diagnostic.
Some works merge discovery with EDL into a single pipeline that discovers and enforces the symmetry, the model becoming more equivariant as evidence for it accumulates \citep{benton2020anips,otto2025,escriche2025}.
Formally, given samples $(x_i,y_i)\sim p$, one seeks a group $G$ acting by $g\cdot(x,y)=(\rho_X(g)x,\rho_Y(g)y)$.
Generative approaches test the stronger condition $p\circ g^{-1}=p$, but most works target the equivariance of $f^*$, $f^*(\rho_X(g)x)=\rho_Y(g)f^*(x)$.
Since the trivial group $\{e\}$ always qualifies, works add regularizers to favor larger groups \citep{benton2020anips,yang2023p4icml,escriche2025}.
Weaker forms of discovery include estimating the range or level of a symmetry \citep{benton2020anips,ouderaa2022}.


\subsubsection{Seminal works, data level}
The problem of detecting symmetries from finite data is not new, and early works extract the symmetries of one shape (data point locus), not of a distribution or a target function.
\citet{highnam1985} and \citet{wolter1985tvc} provided algorithms for finding all rotational and reflective symmetries of planar and three-dimensional point sets.
An early machine learning approach by \citet{rao1998anips} framed invariance discovery in visual perception (problem locus) as learning Lie group operators from infinitesimally transformed image patches via an unsupervised Bayesian algorithm.


\subsubsection{Generative approaches, data distribution symmetry}
\citet{desai2022prd} introduce a GAN-based framework, where the generator learns maps that preserve the data distribution relative to a reference, formalizing a statistical notion of symmetry.
\citet{yang2023p4icml} propose \textit{LieGAN}, where a generator learns a Lie algebra basis $\{A_i\}$ such that transformations of the form $\exp(\sum_i w_i A_i)$ preserve the data distribution, with an additional discriminator enforcing \textit{distribution matching}, while regularizers on the generators prevent both identity collapse and redundancy among the learned basis elements.


\subsubsection{Discovery from task-trained models}
Here the symmetry is recovered from a model fitted to a task, either extracted post hoc or learned jointly as a parameter.

\paragraph{Continuous symmetries}
Continuous symmetries benefit from the smooth structure of Lie groups.
\citet{dehmamy2021anips} introduce Lie Algebra Convolutional Networks (L-conv), whose layers make the $n_L$ Lie algebra generators ($n_L$ a hyperparameter) learnable matrix parameters, allowing joint discovery by training L-conv layers:
\begin{equation}
    \begin{aligned}
        f \in \mathbb R^{d\times m} &\mapsto f\,W^{0\top} + \sum_{i=1}^{n_L} \hat L_i\, f\, W^i, \qquad \hat L_i \in \mathbb R^{d\times d},\; W^i \in \mathbb R^{m\times m'}.
    \end{aligned}
\end{equation}
For large dimension $d$, the $d\times d$ matrices $\hat L_i$ can be factored.
Additionally, regularization terms can be put on Lie generators \citep{escriche2025} to prevent collapse.
\citet{krippendorf2020} bet on the fact that a classifier trained on symmetric data will collapse symmetry-related inputs to nearby positions in the embedding layer.
They exploit this by performing linear regression on pairs of nearby embedding points to extract infinitesimal generators, effectively reading the symmetry off the network's representations.
\citet{vanderwilk2018anips} incorporate invariance desiderata directly into the prior of a Gaussian Process by constructing an invariant kernel that averages over an orbit of transformations.
The parameters of these transformations are then learned end-to-end by maximizing the marginal likelihood, which naturally penalizes complexity and avoids overfitting to spurious symmetries.
\citet{wieser2020anips} employ an autoencoder framework in which the latent space is partitioned into a property-sensitive and a property-invariant subspace, with a mutual information regularizer enforcing that the invariant subspace carries no information about a target quantity.
\citet{ouderaa2024} take a physics-inspired approach via Noether's theorem: they parametrize symmetries as quadratic conserved quantities of a Hamiltonian system and learn them by maximizing a variational lower bound.

\paragraph{Parameter-Sharing schemes for discrete groups}
Discrete symmetries can be approached through search over parameter-sharing schemes \citep{zhou2021,yeh2022p2icais}.
A layer's equivariance to a discrete group is determined by which weight entries are shared, so the problem of symmetry discovery can be framed as learning the partition of weights into equivalence classes, which factorizes the weight space into a structural component (the parameter-sharing scheme) and a value component (the actual shared parameters).
\citet{zhou2021} apply the idea in a multitask scenario via a reparametrization of a linear layer $W\in \mathbb R^{mn}$ as $\mathrm{vec}(W) = Uv$, where their \emph{symmetry matrix} $U\in \mathbb R^{mn\times k}$ encoding the sharing pattern is optimized in a meta-learning outer loop while $v\in \mathbb R^k$ contains learned coefficients that adapt per task.
They prove that this factorization can represent equivariance to any finite group, and they learn translation, rotation and flip patterns in $U$.
\citet{yeh2022p2icais} express parameter sharing by writing the parameter vector $\theta\in \mathbb R^n$ in terms of $\tilde \theta\in \mathbb R^n$ via a row-stochastic \emph{assignment matrix} $A\in \{0,1\}^{n\times n}$, as $\theta=A \tilde \theta$.
Parameters are shared when $A$ is not a permutation matrix.
They jointly learn $\tilde \theta$ and meta-learn $A$ on a validation set.
When $n$ discourages combinatorial search, they relax $A\in [0,1]^{n\times n}$ and add regularizers: an entropy term $H(A) = -\sum_{i,j} A_{ij}\log A_{ij}$ to push values close to $\{0,1\}$, and a nuclear norm $\|A\|_*$ to encourage low rank, hence sharing.


\subsubsection{Discovering group factors: disentangled symmetries}
A line of work frames symmetry discovery as learning not only the group but also its decomposition into independent factors, a form of disentangled discovery \citep{higgins2018}.
\citet{cohen2014} propose a probabilistic model that learns irreducible representations for compact commutative Lie groups, in particular the toroidal subgroups of $SO(D)$.
In the context of dynamic environments where observations are paired with transformations, \citet{quessard2020anips} address disentanglement with a sparsity regularizer that encourages each symmetry generator to act on as few latent dimensions as possible, thereby identifying a block-diagonal structure in the representation.
\citet{sanborn2022} introduce Bispectral Neural Networks, which learn a Fourier basis (irreducible representations) for compact commutative groups, from which the Cayley table can be read off.
\citet{ennes2025fcm} tackle the problem from a geometric perspective: given a point cloud sampled from orbits of a compact Lie group, their \textit{LieDetect} algorithm estimates the Lie algebra via a PCA-like operator and projects onto the closest valid Lie algebra for a set of candidate groups ($SO(2)$, $T^d$, $SU(2)$, $SO(3)$), with theoretical stability guarantees given in terms of Wasserstein distances.


\subsubsection{Caveat}
Despite the appeal of symmetry discovery as a path around the symmetry bottleneck or as a diagnostic tool, \citet{rath2024air} caution that, while discovery is valuable when prior knowledge is unavailable or hard to specify, automatically learned symmetries are generally unable to outperform models whose equivariance is correctly specified by design.
Misspecification can also degrade performance, as seen in \Cref{sec:edl-tradeoffs}.



\subsection{Symmetry Extraction from Learned Models}\label{sec:sym-extract-model}
Moving from the problem or data to the \emph{parameter} locus, extraction refers to a post-hoc diagnostic on a trained network: the symmetries are those of the implemented function $f_\theta$ at a given $\theta$.
\Cref{sec:disco-capability} first asks whether a network not hard-coding $G$ can learn it at all, i.e.\ whether a symmetry of $f^*$ or $p$ transfers to $f_\theta$ (\Cref{fig:sym-loci}), which requires an expressive enough architecture and that such features are learnable.
\Cref{sec:disco-identification} is discovery proper, with the least a priori knowledge: $\theta$ is studied in isolation and the group, or its Lie-algebra generators, is read off $f_\theta$.
\Cref{sec:sym-quantification} starts with more a priori knowledge: the transformation is given, and what is measured is the tolerance $\epsilon$ up to which $f_\theta$ respects it, i.e.\ the degree of \emph{approximate} symmetry.


\subsubsection{Capability}\label{sec:disco-capability}
% lenc2014a: if the linear map exists but only one transformation at a time
Before asking \emph{which} or \emph{how much} symmetry a model has acquired, a more basic question is: can a network whose architecture does \emph{not} hard-code a group $G$ incorporate, represent or learn symmetries at all, and under what conditions?
\citet{olah2020d} give a mechanistic demonstration by inspecting \emph{circuits} (human-interpretable subgraphs) of a trained convolutional network.
They find, with a combination of activation maximization and visual interpretation, clusters of ``equivariant'' neurons that respond to transformed variants (rotation, hue, scale) of the same feature.
They then show how downstream circuits combine them to synthesize invariant units at deeper layers, i.e.\ invariance can be built compositionally from equivariant parts.
\citet{barenboim2021sp} give a decodability argument: a network trained on a symmetry-agnostic decoy task nonetheless encodes the symmetry class of the data in its penultimate representations, which a secondary classifier recovers ($O(2)$, $\mathbb{Z}_2$, $T$, $T_n$ or $\emptyset$) largely above chance.
Symmetry is thus learned in the absence of any explicit symmetry label, since a structure that can be decoded must have been represented.\newline
These existence results leave open \emph{when} learning succeeds, a question that can be approached through training regimes.
In the lazy (NTK) regime, \citet{perin2025} prove a negative result: generalization to a held-out orbit point is governed entirely by local kernel geometry through the spectral ratio $\epsilon_s = \lambda_N^{-1}/\langle\lambda^{-1}\rangle$ (class separation versus orbit density), with no mechanism to transport a symmetry.
Concretely, they illustrate that rotation poses seen only on some classes cannot be extrapolated to other classes with missing poses in the training set: if not hard-coded, the symmetry is not discovered.
This places symmetry learning in the feature-learning regime.
In the active (feature-learning) regime, \citet{jacot2025} establish a positive counterpart: depth lets a network learn compositions $f_\theta = h \circ q$, and when the target is $G$-invariant the inner map $q$ can play the role of the quotient map $\pi:\mathbb{R}^{d_{\mathrm{in}}} \to \mathbb{R}^{d_{\mathrm{in}}}/G$ (\Cref{def:quotient}).

\subsubsection{Identification}\label{sec:disco-identification}
Beyond proving that neural networks can represent symmetries, another challenge is to identify the symmetry a trained model has internalized: to recover the group acting on the data or, for continuous groups, the Lie-algebra generators.
Two types of methods read the symmetry (1) off the network's \emph{representations} or, more locally, (2) off its \emph{gradients}.\newline
The representational approach inspects what the trained model has stored.
\citet{krippendorf2020} exploit the empirical observation that, when a classifier is trained on a label invariant under a group action such as the discretized value of a rotation-invariant potential $V(\|x\|)$, points on the same orbit collapse into clusters in the embedding layer that are tighter than generic same-class points.
That is, if $(x_1, y_1)$, $(x_2, y_2)$, $(x_3, y_3)$ are three samples with the same label $y_1 = y_2 = y_3$ but only $x_1$ and $x_2$ lie on the same orbit, $x_2\in G\cdot x_1$, then the embedding distance $d(f_\theta^{(l)}(x_1), f_\theta^{(l)}(x_2))$ is expected to be smaller than $d(f_\theta^{(l)}(x_1), f_\theta^{(l)}(x_3))$.
This clustering is used as a probe: orbit point clouds detected in embedding space are pulled back to the input space, and infinitesimal generators $A_a$ are recovered by linear regression on nearest-neighbor pairs $(x,x')$ under the first-order relation $x' \approx x + \epsilon_a A_a x$, yielding the $SO(n)$ and $SU(2)$ generators even from noisy, sparse point clouds.
\citet{olah2020d} work one unit at a time: \emph{activation maximization} generates inputs that strongly excite neurons, and direct visual interpretation of this preferred stimulus names the feature it encodes.
A group of neurons whose preferred stimuli are rotated, hue-shifted or rescaled copies of one feature is then read as an ``equivariant'' family.
Tracing a neuron's activation along an orbit, a \emph{tuning curve}, further confirms the symmetry: a flat curve signals invariance while shifted curves indicate equivariance.\newline
The gradient approach achieves identification by a nullspace computation.
Restricting to matrix Lie groups $G \leq GL(n,\mathbb{R})$, \citet{moskalev2023} (LieGG) read the trained network as a discriminator $F$ and invoke the criterion that a generator $A\in\mathbb{R}^{n\times n}$ is an infinitesimal symmetry if and only if $\sum_{i,j} \partial_{x_i} F(x)\,A_{ij}\, x_j = 0$ on the data, i.e.\ the Lie derivative of \Cref{sec:lie-algebra} vanishes.
Stacking this constraint over the dataset $\mathcal{D}$ yields the \emph{network polarization matrix} $\mathcal{E} \in \mathbb{R}^{|\mathcal{D}| \times n^2}$, whose nullspace is the learned Lie algebra.
In practice, the approximate nullspace is extracted from the right-singular vectors of $\mathcal{E}$ associated with near-zero singular values (SVD), the count of which gives the algebra's dimension.
\citet{hu2025nn} (LieSD) follow the same SVD-nullspace philosophy but characterize \emph{equivariance} from inputs, outputs and gradients: for connected $G$, $f$ is $G$-equivariant if and only if $d\rho_Y(A_i)\,f(x) = \partial f(x)\,d\rho_X(A_i)\,x$ for every generator $A_i$, with invariance the special case $d\rho_Y = 0$.
Assembling these constraints into a coefficient matrix, the zero singular values again count the Lie-algebra dimension and singular vectors are the generators.\newline
Notably, both \citet{krippendorf2020} and \citet{hu2025nn} rely on the assumption that a sufficiently trained network absorbs the symmetry present in its data, so identifying the model's symmetry doubles as \emph{discovery} of the data's symmetry, closing the loop with \Cref{sec:sym-disc-from-data}.

\subsubsection{Quantification}\label{sec:sym-quantification}
% mention sensitivity analysis as the broader field that studies how the variation in the output of a model can be apportioned to different sources of variation in its inputs - see storm2024prl for an example of local sensitivity analysis applied to deep networks. Symmetry quantification could be likened to sensitivity analysis to specific directions in input space.
% petrache2023 3.2 equivariance error
% yeh2022p2icais introduce the Partition Distance, a quantitative metric for comparing their learned sharing scheme against a known ground truth (an equivariance error)
% qu2024anips: sec 5.4 cosine similarity >0.99 for atom force prediction in a scaled non-equivariant model
% wang2022p3icml: equivariance error (def 3.1)
% elhag2026: equivariance metrics macro and micro averaging over group orbit (eq 8, eq 9)
Once one knows or assumes \emph{which} symmetry a model should carry, the natural follow-up is \emph{how much} of it the model has actually learned.
The works below share a common template: a transformation (or its infinitesimal generator) is assumed to be given, and one measures a notion of sensitivity of the model under it, assessing how well invariance or equivariance is respected.
For this reason, this part could also belong to \Cref{sec:incomplete-symmetries}.\newline
Among seminal works, \citet{goodfellow2009anips} define a neuron-level \emph{invariance score} $S(i) = r_{\mathrm{loc}}(i)/r_{\mathrm{glob}}(i)$, where the \emph{local firing rate} $r_{\mathrm{loc}}(i)$ is the fraction of transformed versions of a maximally activating input $z$ that still excite neuron $i$, and the \emph{global firing rate} $r_{\mathrm{glob}}(i) = \mathbb{E}_{x\sim P}[f_i(x)]$ measures indiscriminate firing under a random sampling $P$.
A neuron scores high only if it tracks transformations of its preferred stimulus while staying silent on generic inputs, because the denominator enforces selectivity.
A network-level score averages $S(i)$ over only the top-$p$ fraction of neurons in a layer, which permits a \emph{modular} detection in which one subpopulation specializes to rotation, another to translation or illumination.
The score is parametrized by a \emph{type} $\tau$ (an action), a \emph{range} $\Gamma$ (a partial orbit, not necessarily the full one) and the sampling $P$: varying these parameters instantiates different tests under the same template, e.g.\ on video $\tau$ can be forward time-translation and $\Gamma$ the maximal number of frames skipped.
Rather than measuring invariance directly, \citet{lenc2014a} test whether a learned representation $f_\theta^{(l)}$ is compatible with equivariance by fitting a linear map $M_g$ such that $M_g\,f_\theta^{(l)}(x) \approx f_\theta^{(l)}(gx)$, but only one transformation $g$ at a time.
This map is estimated by sparse regression with structured-sparsity constraints respecting the convolutional locality of $f_\theta^{(l)}$, making the problem tractable.
Two quantifications are used: (i) the reconstruction error $\mathbb{E}_x[\|f_\theta^{(l)}(gx)-M_g f_\theta^{(l)}(x)\|]$ evaluates the goodness of fit, i.e.\ the degree of linear equivariance of a representation, and (ii) invariance is assessed by measuring how close $M_g$ is to the identity.
In practice, they look at the off-diagonal weights of $M_g$ through the ratio $\frac{\|(M_g)_{i,:}\|_2}{\|(M_g)_{i,:}-(M_g)_{i,i}e_{i}^\top\|_2}$, where $e_{i}$ is the $i$-th canonical basis vector.
Applied layer-by-layer to AlexNet, it reveals that equivariance to geometric transformations (vertical flips, rotations) degrades with depth.

Among subsequent works, translation equivariance in CNNs is studied by \citet{zhang2019p3icml}, who track an internal feature distance $d\big(g(f_\theta^{(l)}(x)),\,f_\theta^{(l)}(gx)\big)$ under the translation $g$ across layers $l$, exposing how downsampling destroys shift-equivariance.
A simple performance check is used by \citet{finzi2020p3icml} (App.~B.1): the prediction gap under a group transformation of the input, which is expected to vanish up to stochasticity and numerical approximation for a truly equivariant model.
Logit-based scores sharpen this input-output analysis: \citet{hutchinson2021p3icml} report an \emph{invariance error} $\lVert \Phi(g\cdot x) - \Phi(x)\rVert_\infty$, the worst-case logit deviation under $g$.
\citet{moskalev2023p2awtagmlt} propose other invariance errors based on the consistency of logits and saliency maps along orbits.
Crucially, they distinguish \emph{genuine} invariance from a weaker \emph{learned} invariance, in which a network (e.g.\ trained by data augmentation) reaches correct predictions on an orbit while maintaining \emph{separate} internal representations per transformation.
Under an MNIST $\to$ Fashion-MNIST interpolation, their predictive-distribution invariance $DI_f = \sum_{x\sim \mathcal{D}}\sum_{g\sim G} D_{\mathrm{KL}}\big(f(x)\,\Vert\,f(gx)\big)$ finds that learned invariance deteriorates rapidly with the domain shift, a direct empirical argument for the robustness that strictly equivariant architectures guarantee by construction.
In the same vein, an equivariance error $\frac{1}{n} \sum_{i=1}^{n} \operatorname{std}(L_{R_i} \Phi(f_i) - \Phi(L_{R_i} f_i)) / \operatorname{std}(\Phi(f_i))$ is proposed by \citet{cohen2018}, measuring an average normalized difference between pre- and post-application of the group actions $L_{R_i}$ on feature maps $f_i$.
\citet{shen2026} instead define an \emph{equivariance error} as the share of the total loss attributable to non-equivariance, $\mathcal{L}_{\mathrm{equiv}}(\theta)/\mathcal{L}(\theta)$, and observe that, for augmentation-trained non-equivariant models, it drops quickly to low values ($\approx 2\%$ of the loss).\newline
A final family quantifies sensitivity infinitesimally.
\citet{gruver2024} use the \emph{Lie derivative} $\mathcal{L}_A f$ of \Cref{sec:lie-algebra}, where $A\in\mathfrak{g}$ is an infinitesimal symmetry generator, yielding a hyperparameter-free local equivariance error that quantifies local sensitivity in the symmetry direction.
\citet{leroy2025} instead estimate, in the embedding space $\mathcal{Z}$, the unit tangents $\hat v(z)$ of discrete curves $t\mapsto f(\alpha_t x)$ traced by a one-parameter attribute action $\alpha_t$ resulting from interpretable attribute modifications (e.g.\ aging a person in an image).
They define an \emph{invariance energy} as a misalignment of neighboring tangents, i.e.\ a folding of $\mathcal{Z}$, rather than as a vanishing derivative.
The nullspace methods of \Cref{sec:disco-identification} used for identification can serve as quantification: \citet{moskalev2023} (LieGG) read off a \emph{symmetry variance} $\mathcal{V}(\mathcal{E}) = \sigma_{\min}(\mathcal{E})^2/|\mathcal{D}|$ characterizing how close the network is to admitting an invariant direction in data space, and, against a known generator, a \emph{symmetry bias} measuring how faithfully that symmetry is respected.
\citet{hu2025nn} (LieSD) push the comparison to ground truth further with a \emph{space error} and an \emph{orthogonality error}.
The proposed space error is $E_{\mathrm{space}} = \min_{X} \|VX - V^*\|_F^2 + \min_{Y} \|V^*Y - V\|_F^2$, which encodes the mutual expressibility of the spans of the true ($V^*$) and recovered ($V$) generator bases.
The orthogonality error is $E_{\mathrm{orth}} = \sum_{i=1}^{k-1} \sum_{j=i+1}^{k} |\langle v_i, v_j \rangle|$ and computes a notion of redundancy among recovered generators $v_i$.



\section{Incomplete Symmetries}\label{sec:incomplete-symmetries}
% ha2017: hypernetworks in RNN break weight-sharing through time
Quantification (\Cref{sec:sym-quantification}) measured how far a trained $f_\theta$ is from a symmetry.
This section continues in that direction and gathers symmetries that are \emph{inexact}, i.e.\ not \emph{exact} along the exactness axis of \Cref{tab:taxonomy}.
They span several loci (\Cref{fig:sym-loci}), though the \emph{model} locus ($\mathcal{H}$) is the most investigated.
Symmetry breaking (\Cref{sec:symmetry-breaking-dl}) is mostly \emph{descriptive}, cataloging failure modes: why and when an expected symmetry is lost.
The two remaining subsections are more prescriptive and explain how a symmetry weakens.
They split along the exactness axis one level down: an \emph{extent-restricted} symmetry holds exactly but only for a restricted set of transformations or points (\Cref{sec:extent-restriction}), while an \emph{approximate} one holds on all of $G$ up to a tolerance (\Cref{sec:soft-equivariance}).

The two relaxations are contrasted by the \emph{$\epsilon$-stabilizer} of \citet{petrache2023}.
Let $G$ be a candidate group of transformations and write the equivariance error of a map $f$ under $G$ as:
$$
    \operatorname{ee}\colon\mathcal{F}\times G\to[0,+\infty),\qquad \operatorname{ee}(f,g)=\big\|\rho_Y(g)^{-1}f(\rho_X(g)\,\cdot)-f\big\|_\infty.
$$
The $\epsilon$-stabilizer at tolerance $\epsilon\geq 0$ keeps the transformations that preserve $f$ up to $\epsilon$:
$$
    \operatorname{Stab}_\epsilon(f)=\{g\in G:\ \operatorname{ee}(f,g)\le\epsilon\}.
$$
For a hypothesis class, it becomes $\operatorname{Stab}_\epsilon(\mathcal{H})=\cap_{h\in \mathcal{H}}\operatorname{Stab}_\epsilon(h)$.
The value on the exactness axis can be read from this criterion: the symmetry is \emph{exact} if $\operatorname{Stab}_0(f)=G$ and \emph{approximate} if $\operatorname{Stab}_0(f)\neq G=\operatorname{Stab}_\epsilon(f)$ for a small $\epsilon>0$ (\Cref{sec:soft-equivariance}).
Since $\operatorname{ee}$ takes a supremum over $\mathcal X$, $\operatorname{Stab}_0(f)$ cannot express a symmetry holding on part of the domain only, and it is always a subgroup of $G$.
\emph{Extent-restricted} symmetries are therefore defined on pairs $(g,x)$ in \Cref{sec:extent-restriction}, of which $\operatorname{Stab}_0(f)=H\subsetneq G$ is the subgroup case, whereas $\operatorname{Stab}_\epsilon(f)\subsetneq G$ with $\epsilon>0$ combines restriction and tolerance.
\Cref{fig:exactness-rotation} illustrates the three cases for rotations of a 2D image.

\begin{figure}[!tbp]
    \centering
    \includegraphics[width=\textwidth]{images/exactness-rotation.pdf}
    \caption{The exactness axis for rotations of a 2D image.
    Top: for each image, the green arrow spans the symmetry range, and waviness indicates an approximate symmetry.
    Bottom: invariance error as a function of rotation angle $\alpha$ (blue).
    In (\textbf{a}) and (\textbf{c}), $f$ is the image itself (color coordinates).
    In (\textbf{b}), $X$ is the half annulus seen as a subset of $\mathbb R^2$ and $f(X)=\min_{x\in X}\|x-p\|_2^2$ for the fixed point $p$, so that $\{\alpha:(\alpha,X)\in\mathcal S(f)\}=S=[-\pi/2,\pi/2]$ is a non-closed subset.
    Best seen in color.}
    \label{fig:exactness-rotation}
\end{figure}

Incomplete symmetries are dual to symmetry discovery (\Cref{sec:sym-disco}).
Discovery is additive: it accumulates evidence \emph{for} a transformation being a symmetry, whereas this section examines how assumed symmetries fail, a form of \emph{negative discovery}.
The two further intertwine: in partial equivariance the subset $S\subseteq G$ on which the symmetry holds can be a discovery objective, while approximate equivariance needs to calibrate how much deviation to tolerate.


\subsection{Symmetry Breaking}\label{sec:symmetry-breaking-dl}
% bamler2018p3icml
Symmetry Breaking (SB) names the situation in which an object carries less symmetry than before, than expected or than possible: a loss of symmetry relative to a group of interest $G$.
The term is ambiguous, and \Cref{table:symmetry-breaking} disambiguates it by where symmetry is lost.
\emph{Functional} SB sits between the problem and model loci: the ground truth $f^*$ is less symmetric than the hypothesis class $\mathcal{H}$.
\emph{Distributional} SB stays at the problem locus, between $f^*$ and the data density $p$: $f^*$ may be equivariant while $p$ is non-uniform along orbits.
Both are \emph{static}, with $G$ acting on $\mathcal{X},\mathcal{Y}$.
\emph{Dynamical} SB is \emph{dynamic} as the name indicates, with $G$ acting on $\Theta$: symmetry is lost along training.

The pointwise categorization of \citet{wang2023a} contrasts the two static meanings.
They call $g$ \emph{correct} at a point $x$ when $p(gx)>0$ and $f^*(gx)=gf^*(x)$ (the digit $0$ under a $180^\circ$ rotation), \emph{incorrect} when $p(gx)>0$ but $f^*(gx)\neq gf^*(x)$ (the digit $6$ sent to $9$), and \emph{extrinsic} when $p(gx)=0$, i.e.\ $g$ carries $x$ off the data support (a rotated $4$).
Incorrect equivariance is functional SB, whereas extrinsic equivariance falls in distributional SB.
Functional SB can be induced in several ways, as illustrated by the $SE(3)$-equivariant \emph{Rapidash} model of \citet{vadgama2025}.
This model has typed inputs (scalars and vectors) that transform differently under $G$, and treating vectors as tuples of scalars deliberately breaks the symmetry of $\mathcal{H}$.
The model also supports \emph{pose-informed} breaking, by supplying an extra token $Z$ encoding a reference frame, while joint equivariance $f_\theta(gX,gZ)=g\,f_\theta(X,Z)$ is retained.

\begin{table}[ht]
    \centering
    \begin{tabularx}{\textwidth}{|l|l|>{\raggedright\arraybackslash}X|>{\raggedright\arraybackslash}X|}
        \hline
        \textbf{Type} & \textbf{Broken object} & \textbf{Condition} & \textbf{Refs} \\
        \hline
        Functional & target $f^*$ vs.\ class $\mathcal{H}$ & $\operatorname{Sym}(f^*)\subsetneq\operatorname{Sym}(h)$ for $h\in \mathcal{H}$ & \citep{wang2023a,zhang2019p3icml,smidt2021prr,lawrence2024,goel2025} \\
        \hline
        Distributional & data density $p$ & $p(x)\neq p(\rho_X(g)x)$ & \citep{lawrence2026,rygiel2026,chen2020,urbano2026} \\
        \hline
        Dynamical & loss $\mathcal{L}$, optimizer $U$, state $\theta_t$ & not unique & \citep{kunin2020,tanaka2021,ziyin2025,arjevani2023} \\
        \hline
    \end{tabularx}
    \caption{Three meanings of symmetry breaking, according to where a symmetry of the group of interest $G$ is lost: in the target, in the data density, or along training.
    For the dynamical row, the conditions can be $G_{\theta_t}$ changing or broken conservation laws.}
    \label{table:symmetry-breaking}
\end{table}

\subsubsection{Functional Symmetry Breaking}\label{sec:functional-sb}
Functional SB is a property of the pair $(f^*,\mathcal{H})$: the target $f^*$ is \emph{less} equivariant than hypotheses in the model class,
\begin{align}
\operatorname{Sym}(f^*)\subsetneq\operatorname{Sym}(h)=G\quad&\text{for all }h\in\mathcal{H},\\
&\text{equivalently}\quad \exists\,x,g:\ f^*(\rho_X(g)x)\neq\rho_Y(g)f^*(x).\notag
\end{align}
It is an overspecification of equivariance and coincides pointwise with the \emph{incorrect equivariance} of \citet{wang2023a}, whose magnitude can be measured by equivariance errors (as in \Cref{eq:equiv-error-approximate} for instance).
It appears in standard architectures: \citet{zhang2019p3icml} show that the subsampling hidden inside strided convolutions and max-pooling silently destroys shift-equivariance, since it violates the Nyquist--Shannon condition $f_s>2f_{\max}$ for lossless sampling of a band-limited signal.
\citet{zhang2019p3icml} propose a low-pass (blur) filtering before subsampling to partially restore it.
Imposing an exact $G$ on a task that lacks it carries an unavoidable cost, which \citet{wang2023a} make precise through error lower bounds: for invariant classification the accuracy is capped by the best per-orbit choice, $\operatorname{err}(f_\theta)\ge\int_F k(Gx)\,dx$, where the \emph{total dissent} $k(Gx)$ is the probability mass of the points of the orbit $Gx$ that do not carry its majority label, and the \emph{fundamental domain} $F$ contains one representative per orbit such that $\mathcal{X}=\cup_{x\in F}Gx$.
With analogous bounds for invariant and equivariant regression, this motivates the partial and approximate relaxations of \Cref{sec:extent-restriction,sec:soft-equivariance}.
The excess symmetry can also enter through specific blocks of an architecture, for instance the readout layer of GNNs: \citet{talhi2026} show that any linear permutation-invariant graph readout (sum pooling, etc.) factors through the Reynolds operator, here acting on features as $\mathcal{R}(x)=\frac{1}{|G|}\sum_{g\in G}\rho(g)x$.
This means that, for every linear permutation-invariant $f$, there is an $f'$ with $f=f'\circ\mathcal{R}$.
This caps expressivity, and they propose readouts invariant only to $\operatorname{Aut}(\Gamma)$, the automorphism group of the input graph $\Gamma$, rather than to the full $S_n$, deliberately breaking the symmetry.
Reducing equivariance is a generic remedy, whether by input-injection (below) or by retyping inputs as in the vector-as-scalars trick of \citet{vadgama2025}.

\paragraph{Generative Ambiguity and Curie's Principle}\label{par:curie-principle}
A sub-instance of functional breaking arises when a symmetric input must produce a less symmetric output.
This is forbidden for a deterministic equivariant map by \emph{Curie's principle}: if $f$ is $G$-equivariant then $\operatorname{Stab}(x)\subseteq\operatorname{Stab}(f(x))$, since $g\in\operatorname{Stab}(x)\Rightarrow\rho_Y(g)f(x)=f(\rho_X(g)x)=f(x)$.
Simply put, no equivariant $f$ can map a square to a rectangle.
The obstruction arises when the problem requires $\operatorname{Stab}(f^*(x))\subsetneq\operatorname{Stab}(x)$.

\citet{smidt2021prr} seek to discover the causes of the obstruction and use an equivariant model as a ``symmetry compiler''.
Starting with an equivariant model, they augment the input $x\mapsto x\oplus x'$ with a learnable auxiliary $x'$ optimized following the gradient $\partial\mathcal{L}/\partial x'$. 
At convergence $x'$ recovers the ``missing data'', so-called \emph{order parameters}, whose absence prevented an equivariant $f_\theta$ from solving the task and which can be interpreted.
Symmetry is broken at the pipeline level while $f_\theta$ itself stays equivariant.
In a similar diagnostic-oriented approach, \citet{wang2024p4icml} fit a relaxed group convolution initialized at exact equivariance (constant filter weights $w_l$): training leaves the network equivariant only to $\operatorname{Stab}(x)\cap\operatorname{Stab}(y)$ (Prop.~3.1), and the non-constant $w_l$ then read off the symmetry-breaking factors at play.
Third, \citet{lawrence2024} take a generative modeling stance and move from equivariant \emph{functions} to equivariant conditional \emph{distributions}, $\mathbb{P}(Y\mid X{=}gx)=g\cdot\mathbb{P}(Y\mid X{=}x)$: Curie's principle applies at the distribution level, while individual samples may break it, restoring the ability to implement asymmetric solutions.

\paragraph{Pose-Informed Symmetry Breaking, Input-Injection}
One generic technique to overcome functional SB is to intentionally inject into the input a feature $v$ carrying a preferred frame, so that a $G$-equivariant backbone becomes equivariant to at most $\operatorname{Stab}(v)=\{g\in G:\rho(g)v=v\}$ (\Cref{def:orbit-stabilizer}).
It is the equivariant analogue of positional encoding.

\citet{vadgama2025} realize it with the reference-frame token $Z$ mentioned above.
\citet{lawrence2024} concatenate $x\mapsto x\oplus\tilde g\,v$ with $v$ learnable and $\tilde g$ a pose sampled uniformly over compatible ones.
Because the function $(x,\tilde g)\mapsto x\oplus\tilde g\,v$ is equivariant, an equivariant model in $x$ can be extended to one in $(x,\tilde g)$.
They stress the importance of this joint equivariance, proving the identity $R(f)-R(\bar f) = \|f^{\perp}\|_{\mathbb{P}(X,\tilde G)}^{2}$, where $\bar f$ and $f^\perp$ are respectively the equivariant and non-equivariant parts of $f$.
This identity is informative when the true map $f^*$ is equivariant: a lack of joint equivariance leads to excess risk.
\citet{goel2025} turn the mechanism into a subgroup selector (cf.\ \Cref{sec:subgroups}): from a backbone $h_\theta:\mathcal{X}\times V\to\mathcal{Y}$ equivariant to a large group $G$, concatenating $v$ to the input with $\operatorname{Stab}(v)=H$ yields $f_\theta(x)=h_\theta(x,v)$ that is exactly $H$-equivariant provided $h_\theta$ is injective in $v$.
Here, however, the ``pose'' is not interpretable: it is closer to the random node features of GNNs than to positional encoding.

\subsubsection{Distributional Symmetry Breaking}\label{sec:distributional-sb}
Distributional SB is a property of the pair $(f^*, p)$, which occurs when $p(\rho_X(g)x)\neq p(x)$ for some $g\in \operatorname{Sym}(f^*)$ and $x\in\mathcal{X}$.
In words, orbits do not carry constant density.
Performance can be affected even when $f^*$ is equivariant: a model with the correct equivariance is by construction blind to orbit-position cues, which a non-equivariant model can exploit as a shortcut.
Whenever these cues are consistent between train and test sets, they become \emph{spurious but useful} correlations.
Accordingly, it has been observed that on canonically presented data (a Dirac-like density, far from uniform on orbits), a plain model can outperform an equivariant one carrying the correct symmetry~\citep{lawrence2026,rygiel2026}.
A second reason (not found in the literature) could be that some functions are simpler to express on canonicalized inputs.
For example, computing the $SE(2)$-invariant distance between $x,y\in\mathbb{R}^2$ needs a single hidden unit after canonicalization by a translation moving $x$ to $0$, then a rotation until $y_2=0$ and $y_1\geq 0$, since it then equals $y_1$.

On the theoretical side, \citet{lawrence2026} prove for invariant ridge regression that the correctly-invariant model can generalize worse than an unconstrained one under distributional SB.
\citet{chen2020} quantify the effect for data augmentation by relaxing exact to approximate distributional invariance $X\approx_d gX$, with the degree of non-invariance measured by an optimal-transport distance $W_1$.
Compared to the distributionally invariant case, the resulting risk bounds carry uncertainty terms $\pm\,\mathbb{E}_g\,W_1(f(gx),f(x))$ that widen as the symmetry bias grows.
\citet{vadgama2025} study the impact on expressivity: the entropy of the pose distribution $\mathbb{P}(Z\mid X)$ interpolates between zero entropy (canonically aligned data, $\mathbb{P}(Z\mid X)=\delta(Z-Z_{\mathrm{true}})$) and the uniform distribution, which forces the model to be pose-agnostic and hence less expressive.
They prove that the first case has higher expressivity (Cor.~3.1).
Empirically, \citet{rygiel2026} find on computational fluid dynamics tasks that non-equivariant models outperform in highly aligned regimes (e.g.\ automotive aerodynamics) while equivariant models win in unaligned ones (e.g.\ hemodynamics), so the dataset's degree of alignment, not the symmetry alone, governs the comparison.
Likewise, a translation-equivariant ($\mathbb{R}^2$) CNN beats an $SE(2)$-equivariant one on CIFAR-10/100 \citep[Tab.~4]{romero2023}, which fits the picture.

\paragraph{Measuring Distributional Symmetry Breaking}
\citet{lawrence2026} propose two model-based metrics.
A \emph{symmetry-bias} metric compares $p_X$ to its symmetrization $\bar p_X(x)=\int_G p_X(gx)\,dg$ through the accuracy of a two-sample classifier trained to tell the two apart.
The metric is defined as:
$$m(p_X):=\operatorname{Acc}_{\mathrm{NN}}(p_X, \bar p_X)=\mathbb{E}_{x\sim \operatorname{Unif}(p_X, \bar p_X)}\big[\operatorname{Acc}(\mathrm{NN}(x))\big],$$
where $\mathrm{NN}$ is a neural network.
A second \emph{task-useful} metric then asks whether the bias is label-informative, as a ratio $\mathcal{L}_G/\mathcal{L}$ comparing a network that predicts the label from a random-but-equivariant canonical pose against a baseline.
Given a pose map $c:\mathcal{X}\to G$, with $c(x)^{-1}x$ the corresponding canonical form, the losses are defined as $\mathcal{L}=\mathbb{E}_x\big[\mathbb{I}(\mathrm{NN}(c(x))\neq f^*(x))\big]$ and $\mathcal{L}_G=\mathbb{E}_{x,g}\big[\mathbb{I}(\mathrm{NN}(c(gx))\neq gf^*(x))\big]$. Here $\mathrm{NN}:G\to\mathcal{Y}$ is another model, trained to predict the label from the pose.
They report that standard benchmarks (QM9, ModelNet40) are notably symmetry-biased.
\citet{urbano2026} instead estimate the orbit density directly with a class--pose IE (invariant--equivariant) autoencoder \citep{winter2022anips}.

\paragraph{Two extremes}
% kalogeropoulos2024: use canonicalization to design ScaleInv
Canonicalization is the limit in which the orbit density collapses to a Dirac mass \citep{ciuclea2025,urbano2026}.
\citet{urbano2026} take the canonical representative to be a \emph{natural pose} obtained from a (robust) Fréchet mean $\bar g=\arg\min_{h\in G}\sum_{i}d(h,g_i)^2$ of observed poses $g_i$ and left-multiply each orbit so that the identity element is that pose, obtaining an orbit distribution centered around the natural pose.
Extrinsic equivariance is the other extreme, in which $g$ sends $x$ entirely off-distribution, $p(gx)=0$.
Whether it helps is itself context-dependent: it is beneficial in \citet{wang2022b} but harmful in the examples of \citet{wang2023a}.

\paragraph{A hierarchical caveat}
Distributional SB may act only at some levels of a feature hierarchy: on MNIST, local features (edges) can be $E(2)$-equivariant while global ones (digits) are not.
A depth-limited equivariant model could exploit the low-level symmetry without paying for the high-level asymmetry \citep[Fig.~2c]{weiler2019anips}.
This perspective is supported by \citet{elsayed2020p3icml}, who report that breaking equivariance only in the later layers performs best on MNIST, CIFAR-10 and CelebA, compared to breaking it in earlier layers and to equivariant architectures.

\subsubsection{Dynamical Symmetry Breaking}\label{sec:dynamical-sb}
We now focus on parameter space $\Theta$ along training, for which the term ``spontaneous'' is sometimes borrowed from physics to describe the lack of intent behind SB.
\Cref{sec:parameter-space-symmetries} develops these symmetries, and we record here only the SB aspects.

\paragraph{Spontaneous SB}
Gradient Flow (GF) overstates how much symmetry survives training.
The modified gradient flow accounting for finite learning rate, momentum and weight decay of \citet{kunin2020} breaks the conservation laws implied by parameter-space symmetries (\Cref{sec:conservation-laws}).
\citet{tanaka2021} show, through a Lagrangian formulation with the loss as potential energy $V$ and the learning rule as kinetic energy $T$, that a symmetry of the loss need not be a symmetry of the \emph{dynamics}.
This can be seen from the Lagrangian $\mathfrak L = T-V$ where $T=\frac{1}{2} m \dot q^2$ is the kinetic energy and $V$ the potential energy.
Parameter-space symmetries leave $V$ invariant by definition, but $T$ is not necessarily invariant.
They give the example of a scaling transformation under which $\dot q^2$ transforms as $(1+s)\dot q^2$: the kinetic part is responsible for breaking the symmetry of $\mathfrak L$ and for altering the trajectory.
Surprisingly, this \emph{kinetic SB} is shown to turn normalization layers into an implicit RMSProp-like adaptive optimizer.
\citet{ziyin2025a} reach similar conclusions from an entropic perspective: with stochastic minibatches and finite steps, the descent follows $\dot\theta=-\eta(\nabla \mathcal{L}+\gamma\theta+\nabla S)$, whose \emph{entropic} term $S(\theta)=\tfrac14\mathbb{E}_B\|\sqrt{\Lambda}\,\mathbb{E}_{x\in B}\nabla \ell(x,\theta)\|^2$ can be invariant only to orthogonal symmetries.
When the learning rate $\Lambda$ is matrix-valued (as in common optimizers like Adam) with $\Lambda\neq\alpha I$, it can be invariant only to \emph{discrete} ones: stochasticity systematically breaks continuous symmetries.
The inverse of SB has also been studied, under the name \emph{symmetry restoration}~\citep{ziyin2025}.
\citet{ziyin2024} study mirror symmetries $\mathcal{L}(\theta)=\mathcal{L}(R\theta)$ with $R=I-2OO^\top$ (subsuming rescaling, rotation and permutation symmetries), whose reflection surface $\{\theta:O^\top \theta=0\}$ is both gradient-stable ($O^\top\nabla_\theta \mathcal{L}(\theta)=0$) and, under enough weight decay or a large enough learning rate, attractive, so training can move \emph{toward} more symmetric parameters.
\citet{ziyin2025} synthesize the two directions: for a finite subgroup $G$, the $G$-symmetric states ($G_\theta\supseteq G$) form the fixed set of the projector $P_G=\frac{1}{|G|}\sum_g\rho_\Theta(g)$, and a symmetry-breaking distance $\Delta^G=\|\theta-P_G\theta\|_2^2$ tracks how learning alternates between symmetry-breaking \citep{ziyin2025a} and symmetry-restoration \citep{ziyin2024} phases across layers and iterations.
For example, a symmetric state for the permutation symmetry occurs when two neurons $i$ and $j$ are isomorphic in the computational graph and have exactly the same input and output weights, which makes them virtually indistinguishable.
Evaluated for the transposition of the pair $(i,j)$, $\Delta^G$ simply yields $\frac12\| \theta_i - \theta_j \|_2^2$.
The distance $\Delta^G$, in practice used with a threshold, unveils a non-trivial symmetry dynamics both through space (layers) and time (iterations).
Finally, \citet{arjevani2023} identify a principle of least SB: for a shallow ReLU network $f_\theta(x)=\mathbf 1^\top\operatorname{ReLU}(Wx)$ with a teacher of stabilizer $S_d$, spurious minima reached by SGD have strictly smaller stabilizers $S_p\times S_q$ with $p+q=d$, yet these are maximal proper subgroups (for $p\neq q$): they break the teacher's symmetry as little as possible.

\paragraph{Designed SB}
Since symmetries can slow optimization and trap training at low-capacity symmetric saddles, \citet{ziyin2024b} remove them purposefully with \emph{syre}, a model-agnostic scheme that anchors a random static bias at initialization, $\mathcal{L}_r(\theta)=\mathcal{L}(\theta+\theta_0)+\gamma\|\theta\|^2$, with the penalty upgraded to $\gamma\,\theta^\top D\theta$, $D_{ii}\sim\operatorname{Unif}(1-\epsilon,1+\epsilon)$, to also discard continuous symmetries.


\subsection{Extent restriction: beyond groups}\label{sec:extent-restriction}
% Add deformation stability [bronstein2021] which is about diffeomorphisms
An extent-restricted symmetry holds for a set $S$ of transformations that falls short of a full group $G$.
To define it more precisely, we use the following algebraic structure, similar to a group except that inverses are not required.
\begin{Definition}
    A \emph{monoid} $(M,\cdot)$ is a set $M$ with an associative product and identity $e\in M$.
\end{Definition}
Let $f:\mathcal X\to\mathcal Y$ and $M$ be a monoid acting on $\mathcal{X}$ and $\mathcal{Y}$.
The pair $(M,\mathcal{X})$ defines the transformations that we consider.
The transformations that are symmetries are collected in
$$
    \mathcal S(f)=\{(m,x)\in M\times\mathcal X:\ f(\rho_X(m)x)=\rho_Y(m)f(x)\}.
$$
An \emph{extent-restricted symmetry} is then a subset $\mathcal K\subseteq \mathcal S(f)$ where composition holds:
$$
    (m,x), (m',mx)\in \mathcal K \implies (m'm,x)\in \mathcal K.
$$
We break down subcases depending on $\mathcal K$ in \Cref{tab:extent-restriction} and start the subsection with the algebraic conditions (rows~2--5), then the domain restriction (row~6).

\begin{Remark}\label{rem:extent-restr-category}
    In categorical terms, $(M,\mathcal X)$ defines a \emph{category}, with objects $x\in\mathcal X$ and arrows $(m,x)\colon x\to mx$, and so does $\mathcal K$ once the identities $(e,x)\in\mathcal S(f)$ are added.
\end{Remark}

\begin{table}[ht]
    \centering
    \small
    \setlength{\tabcolsep}{4pt}
    \begin{tabularx}{\textwidth}{|l|>{\RaggedRight\arraybackslash}X|>{\RaggedRight\arraybackslash}X|}
        \hline
        \textbf{Case} & \textbf{Condition on $\mathcal K$} & \textbf{Example / refs} \\
        \hline
        Exact & $\mathcal K=M\times \mathcal{X}$ and $M$ is a group & equivariant layer \\
        \hline
        Subgroup & $\mathcal K=H\times \mathcal{X}$ with $H\leq M$ groups & downgrade equivariance \citep{goel2025} (mechanism), \citep{weiler2019anips} (choose $G$) \\
        \hline
        Monoid/semigroup & $\mathcal K=N\times \mathcal{X}$ with $N$ monoid/semigroup & downsampling \citep{worrall2019anips} \\
        \hline
        Non-closed subset & $\exists x, \{ m:(m,x)\in \mathcal K \}$ not closed & partial equivariance of \citep{romero2023}, $\pm 90^\circ$ \\
        \hline
        Groupoid & $(m,x)\in \mathcal K \implies (m^{-1},mx)\in \mathcal K$ & overparametrization \citep{sareen2026} \\
        \hline
        Domain & $\mathcal{K}=\{(m,x):x\in \mathcal{X}_0, mx\in \mathcal{X}_0 \}$ for some $\mathcal{X}_0 \subset \mathcal{X}$ & 6/9 digits, data-dependent and instance-level PS of \Cref{sec:scope-param-space-symmetries} \\
        \hline
    \end{tabularx}
    \caption{Extent-restricted symmetries, except for the first (exact) row, which also implies $\mathcal S(f)=M\times \mathcal{X}$.
    As a concrete example, consider 2D rotations.
    The interval of rotations $I:=[-90^\circ,90^\circ]$ is non-closed since $90^\circ\!\circ\!90^\circ=180^\circ\notin I$, whereas $\{0^\circ,90^\circ,180^\circ,270^\circ\}=C_4$ is a discrete subgroup of all continuous rotations.}
    \label{tab:extent-restriction}
\end{table}

\subsubsection{Subgroups}\label{sec:subgroups}
% wang2024p4icml: relaxed equiv model initialized as G-equiv reduced to H-equiv while training due to symmetry breaking (H subgroup of G)
In \Cref{tab:extent-restriction}, a subgroup $H\le G$ is the only restriction that breaks no group axiom: ``incompleteness'' is here intended as a reduction in symmetry rather than a defect in the algebraic structure.
Two questions organize the literature: how to obtain an $H$-equivariant model, and what the subgroup structure of $G$ itself buys.

A direct construction proposed by \citet{goel2025} descends equivariance ($G\to H$): the input is augmented with an auxiliary feature $v$ whose stabilizer is exactly $\operatorname{Stab}(v)=H$, so that a $G$-equivariant backbone $h_\theta(\cdot,v)$ becomes $H$-equivariant and, if $h_\theta$ is injective in $v$, equivariant to no larger group.
A single backbone thus serves any subgroup by swapping $v$.
The descent can also unfold depthwise: \citet{weiler2019anips} keep early layers equivariant to a large group and deeper layers only to a subgroup, along a chain $G=G^{(0)}\ge G^{(1)}\ge\cdots\ge G^{(L)}$, matching the intuition that low-level edges occur at all orientations whereas global structure is more canonical (however the chain is not subgroup specific, making it less a dedicated subgroup mechanism than \citep{goel2025}).
Instead of re-deriving equivariant architectures for each subgroup, \citet{cesa2021} show that an $O(n)$-steerable kernel basis restricts to an $H$-steerable basis for any compact $H\le O(n)$, so equivariance goes from the general to the particular without additional cost.
When the appropriate subgroup is unknown, one may instead \emph{promote} symmetry by enlarging it: \citet{otto2025} (Sec.~8.4) nudge the model towards the largest subgroup consistent with the data.
For discrete groups, they use a relaxed penalty on the number of group elements that break symmetry, of the form $\sum_{g\in G} \sum_l \| g f^{(l)} - f^{(l)}\|^2$, where $f^{(l)}$ denotes the function implemented at layer $l$.
For continuous groups, they penalize the nuclear norm of the Lie derivative operator, a convex surrogate of its rank, with the objective of decreasing the number of independently broken generators.
A pragmatic use of subgroups is \emph{discretization}: a finite $H\le G$ approximates a continuous $G$, trading exact continuous equivariance for tractability \citep{cohen2016,worrall2018cv-e21ecmgs82ppv,ali2026}.

A second role of subgroups organizes the representation itself.
Direct products of subgroups are used by \citet{higgins2018} to define a representation $f\colon \mathcal{X}\to \mathcal{Z}$ \emph{disentangled} with respect to a decomposition $G=G_1\times\cdots\times G_m$ if $\mathcal{Z}$ splits as $\mathcal{Z}=\bigoplus_{i=1}^m \mathcal{Z}_i$ with $f$ equivariant and each factor $G_i$ acting nontrivially only on $\mathcal{Z}_i$.
Several methods have been proposed to learn disentangled representations.
\citet{quessard2020anips} start from trajectories $(o_1, g_1, o_2, g_2 \dots, )$ and map observations $o_i$ to a vector space $V$ and actions $g_i$ to $GL(V)$ while regularizing the number of dimensions acted on by the representation of the $g_i$.
\citet{pfau2020anips} (GEOMANCER) geometrically factorize the global manifold into disentangled submanifolds: $M=M_1\times M_2$, guaranteed whenever there exists a subspace of the tangent space left invariant by the holonomy group of $M$, the group of linear maps of a tangent space obtained by parallel transport around closed loops (under conditions on $M$).
\citet{shakerinava2022anips} formulate a disentangling loss, using a subgroup objective $\mathcal{L}=\sum_{i=1}^{m} \big[\mathcal{L}_{G_i}^{\mathrm{equiv}} + \sum_{j\neq i} \mathcal{L}_{G_j}^{\mathrm{inv}}\big]$.
However, it requires samples where subgroups act in isolation, i.e.\ datasets $\mathcal{D}_i$ where only $G_i$ acts as $\tau (e, \dots, e, g_i, e, \dots, e)$.

Finally, the loss landscape can be structured by stabilizers, as in the principle of least SB of \citet{arjevani2023} (\Cref{sec:dynamical-sb}).

\subsubsection{Partial equivariance: non-closed subsets}\label{sec:non-closed-subsets}
The non-closed row of \Cref{tab:extent-restriction} keeps the identity and inverses but drops \emph{closure} at some point $x$: the set $S=\{m:(m,x)\in\mathcal K\}$ of transformations admissible at $x$ is a non-closed subset of $G$.
\citet{romero2023} define \emph{partial equivariance} to $S\subseteq G$ as
$$f(\rho_X(g)\,x)=\rho_Y(g)\,f(x)\qquad \forall g\in S\subsetneq G,\ \forall x\in \mathcal{X},$$
with $S$ not required to be a subgroup.
Equivalently, in the formalism of \citet{petrache2023}, this corresponds to $S\subseteq\operatorname{Stab}_0(f)$, but since $\operatorname{Stab}_0(f)$ is a subgroup, \citet{romero2023} note that non-subgroup cases are necessarily also approximate (exactness on a non-closed $S$ would require $S$ to depend on $x$, see \Cref{sec:domain-restriction}).
The distinguishing feature of partial equivariance is non-iterability: because $S$ is not closed, composing admissible transformations may leave $S$.
This subtlety, noted by \citet{romero2023}, also means partial equivariance cannot be enforced exactly by stacking layers.
The motivation is that symmetries can be only \emph{ranged}: equivariance to rotations in $[-90^\circ,90^\circ]$ avoids merging the digits $6$ and $9$, or respects a preferred orientation such as a face aligned with the camera \citep{romero2023}.

In the generalization bound of \citet{petrache2023}, partial equivariance contributes a separate term, distinct from the approximate slack $2\epsilon$.
Architecturally, \citet{romero2023} extend group convolution to the case of a subset $S^{(2)}$.
The operation mapping a function $f:S^{(1)}\to\mathbb{R}$ to a function $h:S^{(2)}\to\mathbb{R}$ is defined for $u\in S^{(2)}$ as:$$h(u) = (\psi *_{\mathrm{partial}} f)(u) = \int_{S^{(1)}} \psi(v^{-1}u)f(v) p(u) d\mu_{G}(v),$$ where $p(u)$ is a probability distribution supported exactly on $S^{(2)}$, i.e.\ it disregards non-equivariance happening outside of $S^{(2)}$.
Their Partial G-CNN can stay fully equivariant on some layers while restricting to a subset on others.
On the $6$-versus-$9$ task it recovers the $\pm 90^\circ$ range, and on natural images it shrinks $S$ where full rotation equivariance is unwarranted.

A complementary line does not fix $S$ in advance but estimates the extent from data.
Augerino, introduced by \citet{benton2020anips}, makes the invariance range a learnable parameter, sampling augmentations from a Lie-algebra-parametrized distribution over a bounded interval $[-\alpha,\alpha]$ and learning $\alpha$ jointly with the weights.
\citet{vanderwilk2018anips} likewise learn a bounded transformation range by maximizing the marginal likelihood of a Gaussian process prior.
\citet{urbano2026} recover the extent as the support of a per-class pose distribution: a class-pose decomposition factors an input into $G$-invariant and equivariant parts, and the estimated density over the orbit is precisely the range of the partial symmetry.
Thus, a digit $7$ is found to be symmetric only over roughly $\pm 30^\circ$, while distinct natural poses keep $6$ and $9$ apart.
An early example is \citet{goodfellow2009anips}, who measured neuron invariance over a bounded transformation range $\Gamma$ rather than a full orbit (e.g.\ $\pm 5$ frames of time translation in video).

\subsubsection{Additional algebraic structures}
Monoids/semigroups and groupoids relax the remaining axioms.
Semigroups keep associativity and closure but drop inverses, which is exactly the structure of irreversible operations such as blurring, downsampling, or cropping, where information is destroyed and cannot be recovered.

\citet{worrall2019anips} build scale equivariance using this structure: a Gaussian scale-space generates coarser scales from finer ones but not conversely.
They extend group correlation by a \emph{semigroup correlation}
$$[\psi *_{S} f](s)=\sum_{u\in \Omega}\psi(u)\,L_{s}[f](u),$$
with $\psi$ the filter, $f$ the signal and $L_s$ the (non-invertible) semigroup action, recovering group and ordinary convolutions as special cases. 
They then adapt it for scale-spaces by setting the action $L_s=L_{(A,z)}$ as $[k_{A}^{\Sigma_{0}}*_{\mathbb{Z}^{d}}f_0](A^{-1}x+z)$, where $k_{A}^{\Sigma_{0}}$ is a Gaussian blurring kernel and the semigroup element $(A,z)$ controls dilation and translation respectively.
In self-supervised learning, \citet{gupta2023a} model the set of augmentations $\mathcal{T}$ (crop, blur, color jitter) as a semigroup.
Their equivariant contrastive objective (CARE) encourages a representation $\rho:\mathcal{T}\to O(d)$ under which input augmentations act as rotations of the spherical embedding.
On the theoretical side, some results obtained in \citep{chen2020} (Sec.~4, ``Beyond groups'') naturally extend to semigroups (e.g.\ variance reduction), but the lack of an \emph{orbit} notion ($\exists s \in S, y=sx \nRightarrow Sx=Sy$ unless $s$ is the identity) prevents their \textit{augmented loss}, which relies on orbit averaging, from generalizing directly.

A groupoid weakens compositionality: it is a small category in which every morphism is invertible, yet the composite $g\circ h$ is defined only when the target of $h$ matches the source of $g$.
\citet{sareen2026} show that some overparametrization symmetries of a neural network form precisely such a structure.
Duplicating (splitting) a neuron is always possible, but merging is only a \emph{partial} inverse, available for weights that were tied by a prior split, and composing two such moves requires compatible widths, so global composition fails.
The resulting \emph{overparametrization groupoid} is a partitioned parameter space $(\Theta_{n,n_f})_{1\le n_f\le n}$, where $\theta\in\Theta_{n,n_f}$ if it has $n$ neurons and implements a function realizable with only $n_f$ of them.
The arrows defining the compositional structure of $\Theta_{n, n_f}$ reflect neuron-duplication symmetries.
As a simple example, duplicating a row $W^{(l)}_{i,:}$ into a new row $W^{(l)}_{i',:}$ while splitting the outgoing column, $W'^{(l+1)}_{:,i}+W'^{(l+1)}_{:,i'}=W^{(l+1)}_{:,i}$, does not change the function.

Finally, quivers and category theory have been used to reframe and extend geometric deep learning ideas \citep{gavranovic2024}, define message-passing GNNs equivariant to local isomorphisms \citep{dehaan2020} or study isomorphisms of the computational graphs of feedforward architectures \citep{armenta2021a}.

\subsubsection{Domain restriction}\label{sec:domain-restriction}
The last row of \Cref{tab:extent-restriction} restricts points rather than transformations.
The instance-level and data-dependent PS (Parameter-space Symmetry) of \Cref{sec:scope-param-space-symmetries} are domain restrictions in parameter space.
For instance-level PS, $\mathcal X=\Theta$ is restricted to $\Theta_0$.
As an example, take $\Theta_0=\{\theta:\theta^+_v\le0\}$ \citep{grigsby2023p4icml}: beyond the first hidden layer of a ReLU network, a unit $v$ with non-positive incoming parameters $\theta^+_v$ (weights and bias) never activates.
Translations $T$ of its outgoing weights and positive diagonal scalings $\mathcal D_{>0}$ of $\theta^+_v$ preserve $\Theta_0$, so $\mathcal K\supseteq(T\times\mathcal D_{>0})\times\Theta_0$.
A data-dependent PS acts on $\theta$ but constrains $x$: take the pairs $\Theta\times\mathcal X$ as domain in \Cref{ex:data-PS-GL-activation}, define $F(\theta,x)=f_\theta(x)$ and let $GL(|A|,\mathbb R)$ act by $M\cdot(\theta,x)=(M\cdot\theta,x)$, trivially on $x$ and $\mathcal Y$.
The domain $\{(\theta,x):\ W_{1,A}x\in\mathbb R^{|A|}_{>0}\}$ gives $\mathcal K\subseteq\mathcal S(F)$, as $\sigma$ is the identity on positive entries.

Label-dependent refinements are domain restrictions of the ground truth: they require $f^*(mx)=f^*(x)$ only for $m\in\mathcal K_x$, where the non-global set of transformations $\mathcal K_x$ can depend on the class of $x$ and its pose relative to a natural pose \citep{urbano2026}, or on the instance itself \citep{wang2023a,miao2023p4icml} (\Cref{fig:sym-loci}).
Different classes or instances then carry different symmetries: changing the color of a leaf from green to yellow preserves its class, whereas the same change turns a lime into a lemon.

\subsection{Approximate symmetries: tolerance}\label{sec:soft-equivariance}
% bamler2018p3icml: almost symmetric subspace from weak symmetry breaking by regularizer \psi(Z)
% chen2024jsm: approximately Q-symmetric loss function
% contrast approximate EDL with deformation stability of GDL, see LLM_saved-outputs/responses-saved/GDL_EDL.md
Where the previous subsection restricted the \emph{extent} to a set $S$ falling short of $G$, this one keeps $G$ but loosens equality to a tolerance $\epsilon>0$, exploring the last value (\emph{approximate}) of the exactness axis of the taxonomy (\Cref{tab:taxonomy}).
A symmetry at the problem locus may hold only to first order, be blurred by noise, or drift over a simulation.
Models should then match this inexactness, and relaxation can also ease optimization (\Cref{sec:approx-why-tolerate}).
\emph{Prescriptive} methods split on the object axis.
Relaxed architectures (\Cref{sec:approx-relaxed}) act at the model locus, loosening constraints so that $\mathcal{H}\not\subseteq C_G(\mathcal{X},\mathcal{Y})$.
Soft penalties (\Cref{sec:approx-soft}) act at the parameter locus, selecting a $\theta$ with approximately equivariant $f_\theta$, similarly to data augmentation (\Cref{sec:data-augmentation}, arrow $f^*\to f_\theta$ in \Cref{fig:sym-loci}).
The group $G$ is usually \emph{known}, but relaxed group convolutions have been applied to discover symmetries \citep{wang2024p4icml}.

\subsubsection{Definitions}\label{sec:approx-definitions}
Following \citep{wang2022p3icml,sipiran202523scgpis,ali2026}, the main definition of a map $f$ being \emph{$\epsilon$-approximately equivariant} to $G$ is
\begin{equation}\label{eq:equiv-error-approximate}
    \sup_{g\in G,\,x\in \mathcal{X}} d_Y\big(f(\rho_X(g)\,x),\ \rho_Y(g)\,f(x)\big)\le \epsilon,
\end{equation}
with $d_Y$ a distance on $\mathcal{Y}$ and strict equivariance recovered at $\epsilon=0$.
The supremum on the left is an equivariance error $\|f\|_{\mathrm{EE}}$ measuring the functional symmetry breaking of $f$.
It equals $\sup_{g}\operatorname{ee}(f,g)$ for the per-element error of \citet{petrache2023} introduced at the start of \Cref{sec:incomplete-symmetries}, when $d_Y$ is a norm distance and $\rho_Y$ acts by isometries.
A probabilistic variant, the \emph{insensitivity} of \citet{vanderwilk2018anips}, bounds instead the probability of a large deviation, $\mathbb{P}_{g\sim p(g)}\big([f(x)-f(gx)]^2>\delta\big)<\epsilon$.
A discrete definition is found in \citet{huang2023} for a fixed graph $\Gamma$: since most real graphs have a trivial automorphism group $\operatorname{Aut}(\Gamma)$, they build an approximate symmetry group from a coarsening $\Gamma'$ chosen \textit{close} to $\Gamma$ (in graphon cut distance to be precise).
The approximate symmetry group is then $\mathcal{G}_{\Gamma\to\Gamma'}=(S_1\times\cdots\times S_M)\rtimes\operatorname{Aut}(\Gamma')$, where each $S_i$ is a within-cluster permutation group and $\operatorname{Aut}(\Gamma')$ captures the coarser global symmetry.
Note that these symmetries are to be distinguished from the (passive) permutation symmetry common in the GNN literature (relabeling the nodes, as opposed to automorphisms of a fixed graph).

\begin{Remark}
    The tolerance $\epsilon$ is usually independent of the magnitude of $g$, although a factor $\|g\|_{\mathrm{op}}$ is discussed by \citet{wang2022p3icml} and is implicit in the infinitesimal Lie-derivative formulation of \citet{pertigkiozoglou2024}.
\end{Remark}

On the implementation side, two enforcement mechanisms recur and possibly co-occur.
A \emph{relaxed architecture} loosens the weight-sharing constraint so the layer is no longer forced to $\epsilon=0$, while a \emph{soft penalty} instead leaves the model class unconstrained but adds a penalty (or a prior) that discourages large deviations.
A relaxed architecture is rarely sufficient on its own, since ordinary training favors the less constrained model and must be nudged back towards symmetry \citep{ouderaa2022,wang2022p3icml}.
In the Bayesian case, this can be done directly through the prior \citep{vanderwilk2018anips,ouderaa2023}.

\subsubsection{Why tolerate}\label{sec:approx-why-tolerate}
The governing principle is that the model's symmetry should match the data's, so that hard-coding an exact symmetry is detrimental when the data does not reflect it exactly \citep{zhou2021,finzi2021anips,elsayed2020p3icml}.

\citet{wang2022p3icml} make the cost precise: a strictly equivariant $f_\theta$ approximating an approximately equivariant target $f^*$ must commit an error $\exists x_0\in \mathcal{X},\ \|f^*(x_0)-f_\theta(x_0)\|\ge \tfrac12\|f^*\|_{\mathrm{EE}}$, lower-bounded by the target's own functional symmetry breaking measured in terms of the equivariance error (\Cref{eq:equiv-error-approximate}).
The constraint also propagates to the latent space: \citet{ali2026} prove that a strictly equivariant real-analytic encoder $E$ cannot collapse orbits (under gradient descent), i.e.\ the embedded orbit $E_\theta(G\cdot x)$ generically spans a subspace of dimension equal to that of $G x$, and that exact equivariance forces a copy of the regular representation into the latent space: there is a subspace $V\subseteq\mathcal{Z}$ on which $\rho_Z$ restricts to a representation isomorphic to $\rho_{\mathrm{reg}}$.
Only geometric compression is possible, not true quotienting.
Consequently, the optimal level of equivariance has been framed as a trade-off, neither full nor absent.
This is established for example as a bias--variance sweet spot over a lattice of groups \citep{huang2023} and as a unique optimum in the approximation--generalization analysis of \citet{petrache2023} (\Cref{sec:incomplete-symmetries}), with empirical support in \citep{wang2022p3icml,wang2024p4icml} and through methods that learn the level automatically \citep{ouderaa2022}.

\begin{Remark}
    The converse caveat is worth stating: even when the symmetry is exact, equivariant designs are not guaranteed to be the best-performing models \citep{lawrence2026,rygiel2026}.
\end{Remark}

\paragraph{Training and inference}
Tolerance during training does not necessarily mean it should be retained at test time: some works treat the two regimes differently.
\citet{escriche2025} push towards equivariance at test time by \emph{feature averaging} over the orbit (cf.\ \Cref{sec:da-fa}), and \citet{pertigkiozoglou2024} project the trained model onto the exactly equivariant subspace.
The latter also reframes the whole process as an optimization question, arguing that the equivariant hypothesis space, while not necessarily wrong, is harder to optimize.
Relaxation is meant to smooth training, while exactness can be restored at test time.
\citet{elhag2026} make the compatible empirical observation that equivariant models have a rougher loss landscape.

\subsubsection{Relaxed architectures}\label{sec:approx-relaxed}
The most direct construction generalizes the group convolution by making its filter position-dependent, $k_\theta(v^{-1}u,v)$, so that the kernel depends on the absolute element $v$ and not only on the relative one $v^{-1}u$.
\citet{ouderaa2022} implement this with a bivariate kernel $k_\theta:G\times G\to\mathbb{R}$ fed by random Fourier features of the Lie-algebra coordinates, modulated by frequency parameters $\omega,\omega'$ that control the equivariance strength ($\omega'=0$ is strict equivariance) and are themselves learnable.
\citet{wang2022p3icml} relax both group and steerable convolution by replacing the single kernel with a weighted family, $\Psi_{\mathrm{relaxed}}(g,h)=\sum_{l=1}^{L}w_l(h)\,\Psi_l(g^{-1}h)$ for the group case and $\Phi_{\mathrm{relaxed}}(y)=\sum_{l=1}^{L}w_l(y)\odot\Phi_l(y)$ for the steerable basis.
The layer is exactly equivariant when the $w_l$ are constant over $G$, which a soft penalty $\alpha\sum_{l}\sum_{g,h}\|w_l(h)-w_l(g)\|$ encourages.
An optimal strength of this nudging is observed empirically.

A second family uses an \emph{additive} relaxation, splitting a layer into an equivariant part and a free correction, $f(x)=f_e(x)+\gamma Wx$ with $f_e$ equivariant and $W$ unconstrained.
This is the residual-pathway prior of \citet{finzi2021anips}, whose linear layers are $A+B$ with $A=Q\beta$ in the equivariant subspace ($Q$ a basis of solutions of \Cref{eq:linear-equivariance}, computed as in \citet{finzi2021p3icml}) and $B$ a free residual.
The trade-off is controlled through asymmetric regularization strengths $\lambda_A\|A\|$ and $\lambda_B\|B\|$.
\citet{ouderaa2023} extend this to select the per-layer equivariance level automatically by Bayesian model selection and improve scalability through a rank-one spatial factorization and a further sparsification using (possibly learnable) anchor points.
\citet{pertigkiozoglou2024} repurpose the same additive form as an \emph{optimization} device rather than a misspecification remedy: training a non-strict model that is projected back to an exactly equivariant one at test time ($\gamma=0$), while keeping it close to equivariant throughout training by penalizing the Lie derivative of $W$ along each generator.
The regularization is $\mathcal{L}_{A_i}(W)=-\,d\rho_{\mathrm{out}}(A_i)\,W+W\,d\rho_{\mathrm{in}}(A_i)$, where $d\rho$ is the infinitesimal representation (\Cref{def:infinitesimal-generators}), and the penalization is $\mathbb{E}_x[\sum_i\|\mathcal{L}_{A_i}(W)x\|]$.
Finally, the projection distance is also penalized, with a simple term $\|W\|$.
 
The \emph{low-rank locally connected} (LRLC) layer is an architectural interpolation proposed by \citet{elsayed2020p3icml}, which sits between a convolution (full weight sharing) and a locally connected layer (none).
The filter at position $(i,j)$ is $F^{(i,j)}=\sum_{k=1}^{K}w^{(k)}_{i,j}F^{(k)}$ over a learnable filter bank, with the combining weights themselves factorized to restrict the parameter count.
\citet{zhou2021} reuse LRLC, and \citet{ouderaa2023} adopt a similar factorization strategy alongside additional mechanisms. 

For Bayesian methods, the relaxation lives in the prior.
\citet{vanderwilk2018anips} build an approximately invariant Gaussian-process kernel from a base kernel by summing over augmentations, $k(x,x')=\sum_{x_a\in\mathcal{T}(x)}\sum_{x'_a\in\mathcal{T}(x')}k_g(x_a,x'_a)$, where $\mathcal{T}(x)$ is the set of augmentations of $x$, and select the invariance by maximizing the marginal likelihood.
\citet{finzi2021anips} place different prior strengths on the equivariant and flexible pathways, and \citet{ouderaa2023} make those strengths learnable hyperparameters selected the same way.

\subsubsection{Soft penalties}\label{sec:approx-soft}
Another architecture-agnostic route imposes symmetry through the objective alone.
A first line performs ``symmetry regularization''~\citep{shakerinava2022anips} via group invariants.
\citet{shakerinava2022anips} note that equivariance is a weak condition, as any injective $f:\mathcal{X}\to \mathcal{Z}$ is equivariant for a (generally nonlinear) induced action, $\tau_Z(g,z)=f(\tau_X(g,f^{-1}(z)))$.
They therefore constrain the form of that action by aiming to preserve a chosen invariant $\operatorname{Inv}:\mathcal{Z}^2\to\mathbb{R}$, using only pairs transformed by the same unknown element, $\sum_{x, x'} L\big(\operatorname{Inv}(f(x),f(x')),\,\operatorname{Inv}(f(gx),f(gx'))\big)$ (e.g.\ distances for $E(n)$).
We stress that the element $g$ is not known: it only matters that the same $g$ is applied to both $x$ and $x'$.
\citet{gupta2023a} specialize this to orthogonal representations with an equivariance loss $\mathcal{L}_{\mathrm{equiv}}(f;a)=\mathbb{E}_{x,x'}\big[(f(a(x))^\top f(a(x'))-f(x)^\top f(x'))^2\big]$, combined with InfoNCE terms: (1) an invariance term that keeps the induced rotations small and (2) a uniformity term preventing global collapse.
They prove that a small $\mathcal{L}_{\mathrm{equiv}}$ upper-bounds the approximate-equivariance error $\sup_x\|f(a(x))-R_a f(x)\|$, where $R_a$ is an orthogonal matrix corresponding to the group element $a$.

A second line directly couples an unconstrained network to an explicit equivariance loss $\mathcal{L}_{\mathrm{equiv}}(f_\theta)=\frac1n\sum_i\mathbb{E}_{g\sim G}\big[L(f_\theta(\rho_X(g)x_i),\rho_Y(g)y_i)\big]$, used to encode a symmetry into a free backbone \citep{escriche2025,elhag2026}.
For encoder--decoder models, \citet{ali2026} decompose this into three terms: an input-to-latent term $\operatorname{MSE}(\rho_Z(g)E(x),E(gx))$, a latent-to-output term $L_{\mathrm{task}}(D(\rho_Z(g)E(x)),gy)$ with $\rho_Z(g)$ a learnable per-generator matrix, and a presentation-based penalty encouraging $\rho_Z$ towards being a genuine representation, for example $\operatorname{MSE}(\rho_Z(r)^3,I_d)$ for a $C_3$ group with rotation generator $r$.

\section{Discussion}\label{sec:discussion}

\subsection{Loose Ends and Outlook}\label{sec:loose-ends}
We gather here some occurrences of symmetry that did not enter the main narrative.

\paragraph{Beyond disentanglement}
Disentangled representations \citep{higgins2018}, and the more specific class--pose decompositions \citep{winter2022anips,marchetti2023p2icais} in which the latent space is split into invariant and equivariant parts for a given group, have an information-theoretic extension: an embedding $Z=[Z_Y,Z_\perp]$ with mutual information $I(Z_\perp;Y)=0$, so that $Z_\perp$ parametrizes the transformations under which $Y$ is invariant.
This goes beyond disentanglement since it accounts for a label-informative distributional SB (\Cref{sec:distributional-sb}) and extends beyond groups to complex transformations \citep{wieser2020anips}.

\paragraph{Framework symmetries}
Symmetrization is mentioned in several places (\Cref{sec:edl-achieving-equiv,sec:quotient-repr-align,sec:distributional-sb}) but never unified.
It gathers post-hoc or ad-hoc methods that make a module equivariant, such as canonicalization \citep{kaba2023p4icml,urbano2026,ciuclea2025,kalogeropoulos2024} and frame averaging \citep{puny2021}, an efficient version of the feature averaging of \Cref{sec:da-fa}.
Some paradigms induce symmetries.
In federated learning, the FedAvg update $\theta=\sum_k \frac{n_k}{n}\theta_k$ \citep{mcmahan2017p2icais} is invariant to relabeling the clients $k$.
In asynchronous training, stale gradients can be applied, so the result depends on their arrival order, as in \Cref{sec:data-ordering}.
In reinforcement learning, self-play \citep{silver2017,zhang2025a} trains agents in competition with duplicates of themselves, close to the branch symmetry of joint embeddings (\Cref{sec:joint-embedding}).
Generative models and joint-embedding predictive architectures (JEPA) were not covered.
Evaluations and explanations should respect symmetries too.
In the spirit of the feature-matrix dissimilarity $\mu(X,Y)$ of \citet{godfrey2022}, invariant to the permutations and rescalings of ReLU networks (\Cref{sec:quotient-repr-align}), explanations should be equivariant to the same group as the model \citep{crabbe2023}, and self-interpretable models can produce explanations that respect equivariances \citep{wang2021anips}.
In NLP, semantic-preservation tests \citep{ribeiro2020p5amacl,zarza2026,bianchi2022pnpfwebn} ask a transformation such as a paraphrase to preserve $\phi$-equivalence (\Cref{def:equivalence}), with $\phi$ the model or a text property.

\paragraph{Other}
Still in NLP, symmetries and asymmetries of language are studied: translation-invariant co-occurrences shape word representations~\citep{karkada2026}, and learning language in reverse performs worse, although the best achievable loss is invariant \citep{papadopoulos2024}.
Complex symmetries remain largely open.
Non-linear symmetries can be discovered as linear actions in a latent space \citep{yang2024p4icml}.
A 3D rotation of an object seen as a 2D image has no well-defined action \citep{park2022p3icml,wang2022b}, and context raises similar issues: a bounding box should rotate with its object while the global scene stays unrotated \citep{yu2022}.

\subsection{Underexplored areas and possible gaps}
We collect possible future investigations.

\Cref{tab:gaps} (\Cref{app:gaps}) lists the low-lift pairs (a measure of under-representation) of the coded corpus, among which the following could be plausible candidates.

\paragraph{(Framework, Explicit)}
The design of learning systems could be formalized through symmetry principles, as in our reading of joint embeddings in \Cref{sec:joint-embedding}.
More broadly, framework and data are the values of the object axis with the fewest works in the corpus.

\paragraph{(Prescriptive, Optimization)}
Symmetries in parameter space form an active field, which is still mostly descriptive.
The use of conserved quantities as interpretable training descriptors could be explored, since they influence training~\citep{tarmoun2021p3icml}.
The descriptive side is, however, unfinished: conserved quantities of transformers are only partially derived for multi-head attention and remain open for normalization layers and max-pooling \citep{marcotte2025p4icml}.
The interaction with data augmentation beyond MSE and tensorizable architectures \citep{galley2026}, whether Noether-like frameworks capture all conservation laws, and how these relate to implicit bias \citep{zhao2025} are not well understood.

\paragraph{(Model, Dynamic)}
The low lift may reflect the relatively few works on training equivariant models (though recent ones exist~\citep{nordenfors2024,pertigkiozoglou2024,elhag2026}), limited software and hardware efficiency, or missing end-of-training metrics such as scaling laws.
On when to build equivariance, arguments do not disagree but are not fully aligned either \citep{lippmann2025iclr,brehmer2025,shen2026,vadgama2025,rygiel2026,vandegeijn2025}.
It is unclear whether FLOPs are a good criterion, as equivariant models win per nominal FLOP but reach a lower hardware throughput \citep{brehmer2025}, and to what degree expressive equivariant networks are efficiently trainable \citep{pacini2026,elbaz2025}.
\citet{lawrence2026} hypothesize that locally equivariant features may help even on globally aligned datasets.
Since non-equivariant models can learn symmetries, discover-then-enforce strategies \citep{escriche2025,otto2025} could be cost-effective: an unconstrained prototype is trained, its equivariance error quantified (\Cref{sec:sym-quantification}), and hard-to-learn symmetries then hard-coded.

\paragraph{(Model, Unknown)}
Whether fibers equal orbits, i.e.\ the reverse inclusion $\varphi^{-1}(\{\varphi(\theta)\})\subseteq G\cdot\theta$ (\Cref{sec:parameter-identifiability}), matters for the zero-loss manifold and weight-space learning.
Results exist in special cases \citep{phuong2019,bona-pellissier2023,marchetti2025a,gegenfurtner2026,usevich2025}, so the current goal is to pin down the exact conditions of identifiability.
\citet{shahverdi2026} prove that the generic fiber is finite for generic polynomial activations of large degree, but whether it contains only neuron permutations is open, and \citet{usevich2025} conjecture that finite identifiability implies global identifiability for homogeneous polynomial networks.
Two other directions are (1) complete architectures, e.g.\ a full transformer with normalization, softmax and MLPs \citep{tran2025iclr,zhang2025p4icml}, where \citet{henry2025} conjecture that normalization makes deep self-attention generically identifiable, and (2) going beyond generic $\theta$ (model locus) to the instance level (parameter locus) \citep{grigsby2023p4icml,grillo2026}, for instance the volume of parameters identifiable up to permutation and rescaling \citep{grillo2026}.

\vspace{5mm}

The cell \emph{(Prescriptive, Implicit)} contains $44$ works that could provide empirical recommendations to be understood, explained or formalized.
Historically, batch normalization \citep{ioffe2015p3icml} introduced a scale invariance whose effect is still not fully understood, and SAM \citep{foret2021} did not respect parameter-space symmetries until corrected by \citet{kwon2021p3icml}.

\paragraph{Towards more general groups}
Symmetry discovery could move beyond connected or compact Lie groups and linear actions \citep{ennes2025fcm,escriche2025}, universality results beyond point-wise activations and permutation representations \citep{pacini2026}, and new structures such as groupoids or category-theoretic frameworks could be explored \citep{sareen2026,basheer2025ijdsa}.
Not all groups have been conquered by EDL, where most results concern finite groups, compact groups and semidirect products $\mathbb{R}^n\rtimes K$ with $K$ compact.
Restricting the extent of a symmetry (\Cref{sec:extent-restriction}) is also an active area \citep{goel2025,urbano2026}.
Whether and how networks abstract dataset symmetries outside the lazy regime remains open \citep{perin2025,jacot2025}.
The discretization gap \citep{ennes2025fcm,ali2026} should also be answered: when is a finite subgroup $H\le G$ a good surrogate for a continuous group $G$, i.e.\ how large can the equivariance error $\lVert f\rVert_{\mathrm{EE}}$ (\Cref{eq:equiv-error-approximate}) be for $f$ only $H$-equivariant?

\paragraph{In parameter space}
Equivariance seems less explored than invariance, e.g.\ $f_{\alpha\theta}=\alpha^L f_\theta$ for $\alpha>0$ in a bias-free ReLU network with $L$ layers.
Input symmetries realized in parameter space have started to be studied \citep{galley2026}, whereas approximate parameter-space symmetries, e.g.\ rescaling with GeLU or Swish in place of ReLU, remain open \citep{ziyin2025a}.
The effect of extent restriction on training dynamics has been investigated for specific parameter configurations \citep{ziyin2025}, less for data-dependent symmetries.
For minibatches, writing $S_B(\theta)=\{g: f_{g\cdot\theta}|_B=f_\theta|_B\}$ for a batch $B$, we have $B\subseteq B'\Rightarrow S_{B'}(\theta)\subseteq S_B(\theta)$, so smaller batches carry more symmetries, in line with the batch-dependent functional dimension of \citet{elisendagrigsby2025aim}.

\paragraph{From parameters to functions}
A further goal is to infer properties of $\mathcal{F}$ from $\Theta$ (\Cref{sec:measurements-param-space}).
For example, mode connectivity carries over to function space (\Cref{sec:mode-connectivity}), but what does a low linear barrier after alignment $\min_{g\in G}B(\theta_0,g\cdot\theta_1)\approx 0$ (\Cref{eq:linear-barrier}) mean for $f_{\theta_0}$ and $f_{\theta_1}$?

\subsection{Conclusions}

Symmetry appears at every stage of the deep learning pipeline.
Though not the only key to understanding deep learning, it is a central one.
We organized the literature with a faceted taxonomy that makes the links between results across subfields explicit.
Read as questions (which object carries the symmetry, is it known, is it exact?), the axes also guide the choice of a symmetry-aware strategy.

%%%%%%%%%%%%%%%%%%%%%%%%%%%%%%%%%%%%%%%%%%
\authorcontributions{}

\funding{}

\acknowledgments{}

\conflictsofinterest{The authors declare no conflicts of interest.} 

%%%%%%%%%%%%%%%%%%%%%%%%%%%%%%%%%%%%%%%%%%
%% Optional

\clearpage

\abbreviations{Abbreviations and notation}{%
The following abbreviations are used in this manuscript:\\
\noindent
\begin{tabular}{@{}ll}
CL & Conservation Law \\
CQ & Conserved Quantity \\
DA & Data Augmentation \\
GDL & Geometric Deep Learning \\
LLM & Large Language Model \\
SB & Symmetry Breaking \\
GF & Gradient Flow \\
GD & Gradient Descent \\
SGD & Stochastic Gradient Descent \\
EDL & Equivariant Deep Learning \\
FA & Feature Averaging \\
MC & Mode Connectivity \\
LMC & Linear Mode Connectivity \\
SS & Symmetric States \\
PS & Parameter-Space Symmetries \\
LS & Loss-Level Symmetries \\
NTK & Neural Tangent Kernel \\
GNN & Graph Neural Network \\
CNN & Convolutional Neural Network
\end{tabular}\\[8pt]
The following notation is used in this manuscript:\\
% Three blocks (groups / spaces and maps / losses) so the table can break across pages.
% mdpi.cls makes X columns vertically centered; restore top alignment locally.
\renewcommand\tabularxcolumn[1]{p{#1}}%
\newcolumntype{N}{>{\raggedright\arraybackslash}p{3.8cm}>{\raggedright\arraybackslash}p{3.5cm}>{\RaggedRight\arraybackslash}X}%
\noindent
\begin{tabularx}{\linewidth}{@{}N@{}}
\toprule
\textbf{Object} & \textbf{Notation} & \textbf{Note} \\
\midrule
Group & $G$ & Acts by ${g\cdot x}$, $gx$, ${\rho_X(g)\,x}$ and ${\rho_Y(g)\,y}$ \\
Generic subgroup & $H\le G$ & Also the structure group with gauges \\
Orbit of $x$ & $G\cdot x$ &  \\
Stabilizer of $x$ & $G_x$, $\operatorname{Stab}(x)$ & \\
Symmetry group of $f$ & $\operatorname{Sym}(f)$ & \\
Bijections of a set & $\operatorname{Sym}(\mathcal{X})$ & \\
Permutation, permutation matrix & $\sigma\in S_n$, $P_\sigma$ & ${P_\sigma e_i=e_{\sigma(i)}}$ \\
Quotient map & ${\pi:\mathcal{X}\to\mathcal{X}/G}$ & Sends $x$ to its orbit \\
Section & ${s:\mathcal{X}/G\to\mathcal{X}}$ & A representative per orbit, ${\pi\circ s=\operatorname{id}}$ \\
Canonicalization & ${c=s\circ\pi:\mathcal{X}\to\mathcal{X}}$ & \\
Lie algebra of $G$ & $\mathfrak{g}$ & Elements $A$ \\
Augmentation family & $\mathcal{T}$ & Broader than a group (crop, blur...) \\
\midrule
\end{tabularx}\\
\noindent
\begin{tabularx}{\linewidth}{@{}N@{}}
\midrule
Input space & $\mathcal{X}$ & \\
Output / target space & $\mathcal{Y}$ & \\
Latent / embedding space & $\mathcal{Z}$ & \\
Parameter space & $\Theta$ & \\
Function space & $\mathcal{F}$ & \\
Hypothesis class & $\mathcal{H}$ & Functions realizable by an architecture \\
Dataset & $\mathcal{D}$ & \\
Generic map & $f$ & \\
Ground truth & $f^*$ & Target map of the problem \\
Implementation/realization map & ${\varphi:\Theta\to\mathcal{F}}$ & ${\varphi(\Theta)=\mathcal{H}}$ \\
Implemented function & ${f_\theta=\varphi(\theta)}$ & \\
Feature map at layer $l$ & ${f_\theta^{(l)}:\mathcal{X}\to\mathcal{Z}^{(l)}}$ & ${f_\theta^{(L)}=f_\theta}$ for $L$ layers \\
\midrule
\end{tabularx}\\
\noindent
\begin{tabularx}{\linewidth}{@{}N@{}}
Per-example loss & $\ell(x,\theta)$ & \\
Loss criterion & $L(\hat y,y)$ & Compares a prediction $\hat y$ to a target $y$ \\
Single-argument loss & $\mathcal{L}$ & ${\Theta\to\mathbb{R}}$ or ${\mathcal{F}\to\mathbb{R}}$\\
Regularizer & $\mathcal{R}$ & Penalty term added to $\mathcal{L}$ \\
\bottomrule
\end{tabularx}}

%%%%%%%%%%%%%%%%%%%%%%%%%%%%%%%%%%%%%%%%%%

\appendixtitles{yes}
\appendixstart
\appendix

\pagebreak

\section{Prompts}\label{app:prompts}
The knowledge base of \Cref{fig:selection-flow} is built by the two prompts below, run once per work and in this order.
In both, \texttt{\{\}} stands for the citation key of the work, $\to$ for a tab character, and the trailing options set the model and its reasoning effort.

\subsection{Summary prompt}
\begin{lstlisting}[style=prompt]
    "Open and read the folder raw/{}, particularly the main paper raw/{}/{}.md and summarize it to a file summaries/{}.md following the structure: \
    SUMMARY, KEY DEFINITIONS AND BACKGROUND NOTIONS, RESULTS (theoretical if any, empirical if any), METHODOLOGY, LIMITATIONS, THOUGHTS, ONE-LINE TL;DR" --model claude-sonnet-5 --effort high
\end{lstlisting}

\subsection{Indexing prompt}
\begin{lstlisting}[style=prompt]
    "Read summaries/{} and add a line to index.tsv according to the column headers: citation_key(*\,$\to$\,*)Title(*\,$\to$\,*)Links(*\,$\to$\,*)Keywords(*\,$\to$\,*)Short comment" --model claude-sonnet-5 --effort medium
\end{lstlisting}

\pagebreak

\section{Automated taxonomy}\label{app:automated-taxonomy}

\subsection{Validation of the coding}\label{app:coding-validation}
Each work $w$ of the corpus is coded from its summary by the typed-output model \texttt{jev-1.13.0} (Jev), with one call per work.
For each axis $c$ of \Cref{tab:taxonomy}, with value set $V_c$, we add a default value ``not applicable'' (NA), so that a work is not forced into a value of an axis it does not concern.
Jev returns a probability distribution $q_{w,c}$ on $V_c\cup\{\mathrm{NA}\}$.
Its choice is then $\hat v_{w,c}=\operatorname{arg\,max}_{v}q_{w,c}(v)$.
NA stays rare: over the $408$ coded works, it accounts for $165$ of the $2448$ choices ($6.7\%$), with fluctuations depending on the axis (\Cref{fig:taxonomy-choices}).
Eleven works are NA on every axis.
On inspection, they consist of deep learning works with no symmetry content or methodological references, with the exception of one work \citep{xu2024vj} introducing a data shuffling scheme, which is about data ordering (framework symmetry) but very implicit.

To verify the coding, we annotated by hand $12$ randomly selected works following our intent for the taxonomy, which gives a value $v^*_{w,c}$ for each work and axis.
We keep the pairs for which neither our annotation nor Jev's choice is NA, for a total of $62$ annotations, and the agreement rate between Jev and our annotations is $93.5\%$.
This verification is not exact, and neither is the classification itself: a work can belong to several values of an axis.
This is one of the reasons why we also record the probabilities $q_{w,c}$, which keep track of such split cases.
We illustrate with one of the disagreements, a close call between two sensible values.
\citet{brea2019} study paths whose endpoints are obtained by the permutation symmetry of the architecture (\emph{Model}) through a particular symmetric state (\emph{Parameter}): we chose \emph{Parameter}, Jev chose \emph{Model}.

\subsection{Validation of the taxonomy axes}\label{app:taxonomy-validation}
A taxonomy is useful if each axis separates the corpus and if distinct axes carry distinct information.
We check both on the coded corpus.

\begin{figure}[ht]
    \centering
    \includegraphics[width=\textwidth]{images/taxonomy-choices.pdf}
    \caption{Number of works coded with each value, for each axis of \Cref{tab:taxonomy} (\Cref{app:coding-validation}), with NA in gray.
    Colors only distinguish the values within one axis.}
    \label{fig:taxonomy-choices}
\end{figure}

\paragraph{Each axis splits the corpus}
No axis assigns the same value to all works, and every value is taken by at least $16$ works (\Cref{fig:taxonomy-choices}).
The balance varies among axes, and the least split axis is \emph{A priori knowledge}, with $303$ of $372$ works ($81.5\%$) coded \emph{Known}.

\paragraph{Axes are only weakly dependent}
At the level of axes, each axis is a categorical variable over the corpus.
For two axes $A\neq B$, with $r$ and $k$ values, we keep the $N_{AB}$ works that are not NA on either axis and count the works coded $a$ ($n_a$), $b$ ($n_b$) and both ($n_{ab}$).
With $E_{ab}=n_a n_b/N_{AB}$ the expected count under independence, we report Cramér's V for all pairs in \Cref{fig:taxonomy-cramersV}: all $15$ pairs have $V_{AB}<0.44$, with median $0.19$.

\begin{figure}[ht]
    \centering
    \includegraphics[width=0.8\textwidth]{images/taxonomy-cramersV.pdf}
    \caption{Cramér's V between the axes of \Cref{tab:taxonomy}, on the works that are not NA on either axis.}
    \label{fig:taxonomy-cramersV}
\end{figure}

We also report dependence at the level of values, where we use the probabilities $q_{w,c}(v)$ of the $17$ values (NA excluded) and compute the Spearman correlation over the $408$ works (\Cref{fig:taxonomy-Pcorrelations}).

\begin{figure}[ht]
\begin{adjustwidth}{-\extralength}{0cm}
    \raggedleft
    \includegraphics[width=0.8\linewidth]{images/taxonomy-Pcorrelations.pdf}
\end{adjustwidth}
    \caption{Spearman correlation, over the $408$ works, between the probabilities $q_{w,c}(v)$ of values of distinct axes of \Cref{tab:taxonomy}.
    Blocks within one axis are left blank.}
    \label{fig:taxonomy-Pcorrelations}
\end{figure}

\subsection{Example of analysis}\label{app:taxonomy-analysis}
We report two ways we used the coded corpus to support the framing of the survey: the complementary \emph{segment} and \emph{choice} views.
The segment view selects a section, subsection or subsubsection $s$, gathers its $n_s$ citations, of works $w_1,\dots,w_{n_s}$ (a work cited twice counts twice), and plots, for each value $v\in V_c$ of each axis $c$, the mean
\begin{equation}
    \bar q_{s,c}(v)=\frac{1}{n_s}\sum_{i=1}^{n_s}q_{w_i,c}(v)
\end{equation}
with its $95\%$ $t$-interval (\Cref{fig:taxonomy-segment-view}).
The \emph{choice view} does the converse: it selects a set of values and plots their means for every segment of a given level (\Cref{fig:taxonomy-choice-view}).
For instance, \Cref{sec:sym-disco} is the only section where \emph{Unknown} dominates, as expected for symmetry discovery.

\begin{figure}[ht]
\begin{adjustwidth}{-\extralength}{0cm}
    \raggedleft
    \includegraphics[width=0.8\linewidth]{images/taxonomy-segment-4.pdf}
\end{adjustwidth}
    \caption{Segment view of \Cref{sec:design-principle}: mean probability of each value and confidence interval.}
    \label{fig:taxonomy-segment-view}
\end{figure}

\begin{figure}[ht]
\begin{adjustwidth}{-\extralength}{0cm}
    \raggedleft
    \includegraphics[width=0.8\linewidth]{images/taxonomy-segment-Known-Unknown.pdf}
\end{adjustwidth}
    \caption{Choice view of the values \emph{Known} and \emph{Unknown} of the axis \emph{A priori knowledge} at the section level: mean probability over the citations of each section with confidence intervals.}
    \label{fig:taxonomy-choice-view}
\end{figure}

\clearpage

\section{Gaps from coded corpus}\label{app:gaps}

\begin{figure}[ht]
\begin{adjustwidth}{-\extralength}{0cm}
    \raggedleft
    \includegraphics[width=0.8\linewidth]{images/cooccurrence_lift_heatmap.pdf}
\end{adjustwidth}
    \caption{Co-occurrence lift $L_{ab}$ (\Cref{app:gaps}) between values of distinct axes of \Cref{tab:taxonomy}.
    Pairs within one axis are left blank.
    Blue cells ($L_{ab}<1$) are under-represented combinations, red cells ($L_{ab}>1$) over-represented ones.
    The color scale is centered at independence, $L_{ab}=1$.}
    \label{fig:lift-heatmap}
\end{figure}

\begin{table}[ht]
    \centering
    % Column widths: edit freely; keep their sum below about 0.9\textwidth (the rest goes to padding and rules).
    \newcommand{\colwidthA}{0.3\textwidth}
    \newcommand{\colwidthB}{0.6\textwidth}
    \begin{tabular}{|>{\raggedright\arraybackslash}p{\colwidthA}|>{\raggedright\arraybackslash}p{\colwidthB}|}
        \hline
        \textbf{Pair} & \textbf{Comment} \\
        \hline
        (Exact, Framework) & - \\
        \hline
        (Optimization, Static) & discouraged by the taxonomy\\
        \hline
        (Model, Unknown) & candidate \\
        \hline
        (Problem, Implicit) & - \\
        \hline
        (Framework, Explicit) & candidate \\
        \hline
        (Approximate, Optimization) & - \\
        \hline
        (Exact, Implicit) & - \\
        \hline
        (Prescriptive, Optimization) & candidate \\
        \hline
        (Data, Explicit) & - \\
        \hline
        (Model, Dynamic) & candidate \\
        \hline
        (Exact, Data) & possibly explained by extent restriction of $\varphi$\\
        \hline
        (Optimization, Implicit) & - \\
        \hline
        (Exact, Parameter) & possibly explained by extent restriction of $\varphi$\\
        \hline
    \end{tabular}
    \caption{Candidate gaps: pairs of values $(a,b)$ of distinct axes of \Cref{tab:taxonomy} with lift $L_{ab}<0.6$ (\Cref{app:gaps,fig:lift-heatmap}), by increasing lift. Comments indicate a possible explanation for the low lift value.}
    \label{tab:gaps}
\end{table}

The corpus is coded with one value per axis of \Cref{tab:taxonomy}, or ``not applicable'' (\Cref{fig:selection-flow}).
For two distinct axes $A\neq B$, with values $a$ of $A$ and $b$ of $B$, we keep the $N_{AB}$ works for which the value is different from ``not applicable''.
Under an independence assumption, the expected count is $E_{ab}=n_a n_b/N_{AB}$, where $n_a$, $n_b$, $N_{AB}$ are defined as in \Cref{app:taxonomy-validation}.
The \emph{lift}
\begin{equation}
    L_{ab}=\frac{n_{ab}}{E_{ab}}
\end{equation}
compares the observed count to the independent baseline (\Cref{fig:lift-heatmap}), with $L_{ab}<1$ flagging an under-represented combination.
A low lift does not necessarily imply a gap.
\Cref{tab:gaps} extracts from \Cref{fig:lift-heatmap} the pairs that are most under-represented in the corpus.

%%%%%%%%%%%%%%%%%%%%%%%%%%%%%%%%%%%%%%%%%%

% \input{unfinished-thoughts}
\newpage

\newpage

\begin{adjustwidth}{-\extralength}{0cm}

\reftitle{References}
\bibliography{bibliography}

\PublishersNote{}
\end{adjustwidth}
\end{document}
